% La gestion de l'information — Allemand Tle A — Cours signé OnBuch+
% Programme officiel MINESEC. Compiler avec : tectonic AL16-la-gestion-de-l-information.tex
\newcommand{\DOCMATIERE}{Allemand}
\newcommand{\DOCNIVEAU}{Tle A}
\newcommand{\DOCMODULE}{Chapitre 4 : Média et communication}
\newcommand{\DOCLECON}{AL16}
\newcommand{\DOCTITRE}{La gestion de l'information}
% OnBuch+ — préambule « cours signés » ENRICHI (compilé avec tectonic / XeLaTeX)
% Système visuel « Pop » d'OnBuch+ : fond crème (jamais blanc), bordures encre
% épaisses, ombres « dures » décalées, titres Archivo Black en capitales.
% Version MATHÉMATIQUES : graphes (pgfplots/tikz), boîtes pédagogiques
% (définition, propriété, méthode, attention, le savais-tu, exercice résolu,
% exercice, TP, bilan), page de garde et signature OnBuch+ ancrée en bas.
%
% Macros attendues dans main.tex AVANT \input{preamble} :
%   \DOCMATIERE  \DOCNIVEAU  \DOCMODULE  \DOCLECON (numéro)  \DOCTITRE
\documentclass[11pt]{article}

\usepackage{fontspec}
\usepackage[ngerman,french]{babel} % français = langue principale ; l'allemand via \all{..} / textebox
\usepackage[a4paper,margin=2cm,top=3.4cm,bottom=2.8cm,headheight=1.4cm,footskip=1.3cm]{geometry}
\usepackage{amsmath,amssymb,amsfonts}
\usepackage{mathtools} % \xmapsto, \coloneqq... souvent utilisés par les modèles
\usepackage{cancel} % simplifications barrées (bilans de forces)
% \abs{z} (module, valeur absolue) et \norm{u} (norme), utilisés par les modèles
\providecommand{\abs}{}\let\abs\relax\DeclarePairedDelimiter{\abs}{\lvert}{\rvert}
\providecommand{\norm}{}\let\norm\relax\DeclarePairedDelimiter{\norm}{\lVert}{\rVert}
\usepackage[table]{xcolor} % [table] : fournit aussi \rowcolors (alternance de lignes)
\usepackage{pagecolor}
\usepackage{tikz}
\usetikzlibrary{arrows.meta,positioning,calc,shapes.geometric,shapes.misc,shapes.arrows,decorations.pathmorphing,decorations.pathreplacing,decorations.markings,patterns,shadows,mindmap,trees,fit,backgrounds,matrix,intersections,angles,quotes}
\usepackage{pgfplots}
\usepackage[version=4]{mhchem} % équations nucléaires \ce{^{235}_{92}U}
\usepackage[european,siunitx]{circuitikz} % schémas électriques
\pgfplotsset{compat=1.18}
\usepgfplotslibrary{fillbetween,groupplots} % aires entre courbes (calcul intégral)
\usepackage{enumitem}
\usepackage{fancyhdr}
% [most] charge toutes les bibliothèques tcolorbox (listings, minted, raster,
% documentation...) et gonfle inutilement la mémoire TeX sur un document long
% (~300 boîtes) : on ne charge que ce qui est réellement utilisé.
\usepackage[skins,breakable]{tcolorbox}
\usepackage{graphicx}
\usepackage{colortbl}
\usepackage{booktabs}
\usepackage{tabularx}
\usepackage{diagbox} % case d'en-tête diagonale des tableaux à double entrée
% Colonnes extensibles centrée / gauche / droite (C, L, R), souvent utilisées par les modèles
\newcolumntype{C}{>{\centering\arraybackslash}X}
\newcolumntype{L}{>{\raggedright\arraybackslash}X}
\newcolumntype{R}{>{\raggedleft\arraybackslash}X}
\usepackage{array}
\usepackage{multicol}
\usepackage{multirow}
\usepackage{siunitx}
% parse-numbers=false : accepte aussi \SI{2,5 \times 10^{-3}}{...}
\sisetup{locale=FR,output-decimal-marker={,},group-minimum-digits=4,parse-numbers=false}
\DeclareSIUnit\ohm{\ensuremath{\symup{\Omega}}} % U+2126 absent de la police
\DeclareSIUnit\division{div} % réglages d'oscilloscope (ms/div, V/div)
\DeclareSIUnit\tr{tr} % tours (vitesse de rotation, ex. \SI{1500}{\tr\per\minute})
\DeclareSIUnit\molar{M} % concentration molaire (ex. \SI{3}{\molar}, \SI{1}{\micro\molar})
\DeclareSIUnit\an{an} % année (le modèle confond parfois avec \year, un compteur TeX) — ex. \SI{5}{\kilo\meter\per\an}
\DeclareSIUnit\FCFA{FCFA} % franc CFA (ex. \SI{500}{\FCFA\per\kilo\gram})
\DeclareSIUnit\degreeBrix{\textdegree Brix} % teneur en sucre d'un sirop/jus (réfractométrie)
\DeclareSIUnit\month{mois} % le modèle utilise parfois \month, un compteur TeX, comme unité de durée
\usepackage{qrcode}
% \degree : commande fréquemment utilisée par les modèles pour un angle ou une
% température en dehors de \SI{}{} (non fournie par les paquets chargés).
\providecommand{\degree}{^{\circ}}
% \qed (fin de démonstration), utilisé par les modèles sans amsthm
\providecommand{\qed}{\hfill$\square$}
% \barre : double barre des valeurs interdites dans un tableau de variations (tkz-tab)
\providecommand{\barre}{\ensuremath{\|}}
% environnement proof (amsthm non chargé) utilisé par les modèles
\ifdefined\proof\else\newenvironment{proof}[1][Démonstration]{\par\noindent\textit{#1.}\ }{\qed\par}\fi
\usepackage{lastpage}
\usepackage{hyperref}
\hypersetup{hidelinks}

% ============================================================
% Palette « Pop » OnBuch+ — mêmes hex que la classe P (plus_ui.dart)
% ============================================================
\definecolor{poporange}{HTML}{F59321}
\definecolor{popdark}{HTML}{E07A0C}
\definecolor{popcream}{HTML}{FFF6E8}
\definecolor{popink}{HTML}{1C1714}
\definecolor{poppurple}{HTML}{7A5AE0}
\definecolor{popgreen}{HTML}{2E9E6A}
\definecolor{poppink}{HTML}{E0457B}
\definecolor{popblue}{HTML}{2F7DE1}
\definecolor{popgold}{HTML}{C98A12}
\definecolor{popmuted}{HTML}{8A7F76}
\definecolor{popline}{HTML}{EADFD2}
\definecolor{popgray}{HTML}{8A7F76}
\colorlet{popgrey}{popgray}
% Alias tolérants (le modèle nomme parfois les couleurs différemment)
\colorlet{popwhite}{white}
\colorlet{popblack}{popink}
\colorlet{popred}{poppink}
\colorlet{popyellow}{popgold}
% Teintes claires (fonds de tableaux/encadrés) — le modèle nomme parfois
% les couleurs avec un suffixe L (ex. popblueL) sans les avoir définies.
\colorlet{poporangeL}{poporange!20}
\colorlet{popdarkL}{popdark!20}
\colorlet{poppurpleL}{poppurple!20}
\colorlet{popgreenL}{popgreen!20}
\colorlet{poppinkL}{poppink!20}
\colorlet{popblueL}{popblue!20}
\colorlet{popgoldL}{popgold!20}
\colorlet{popredL}{popred!20}
\colorlet{popgrayL}{popgray!20}
% Teintes claires pour fonds de tableaux / zones
\colorlet{popblueL}{popblue!10!white}
\colorlet{poporangeL}{poporange!14!white}
\colorlet{popgreenL}{popgreen!12!white}
\colorlet{poppurpleL}{poppurple!12!white}
\colorlet{poppinkL}{poppink!12!white}
\colorlet{popgoldL}{popgold!14!white}

\pagecolor{popcream}

% ============================================================
% Typographie
% ============================================================
\defaultfontfeatures{Ligatures=TeX}
\setmainfont{PlusJakartaSans-Medium.ttf}[Path=../../../fonts/,BoldFont=PlusJakartaSans-Bold.ttf]
% Maths en sans-serif (Fira Math) avec chiffres et lettres droites de Plus
% Jakarta, pour que les formules se fondent dans le texte.
\usepackage[math-style=ISO,bold-style=ISO]{unicode-math}
\setmathfont{FiraMath-Regular.otf}[Path=../../../fonts/]
\setmathfont{PlusJakartaSans-Medium.ttf}[Path=../../../fonts/,range={up/num,up/{Latin,latin},\mathup}]
\usepackage{icomma} % virgule décimale sans espace en mode math
% Caractères mathématiques Unicode absents des polices de texte (Plus Jakarta,
% Archivo Black) : rendus par leur équivalent mathématique, en texte comme en
% formule — sinon ils s'affichent en carré vide (ex. « Arithmétique dans ℤ »).
\usepackage{newunicodechar}
\newunicodechar{ℝ}{\ensuremath{\mathbb{R}}}
\newunicodechar{ℤ}{\ensuremath{\mathbb{Z}}}
\newunicodechar{ℂ}{\ensuremath{\mathbb{C}}}
\newunicodechar{ℕ}{\ensuremath{\mathbb{N}}}
\newunicodechar{ℚ}{\ensuremath{\mathbb{Q}}}
\newunicodechar{𝔻}{\ensuremath{\mathbb{D}}}
\newunicodechar{α}{\ensuremath{\alpha}}
\newunicodechar{β}{\ensuremath{\beta}}
\newunicodechar{γ}{\ensuremath{\gamma}}
\newunicodechar{δ}{\ensuremath{\delta}}
\newunicodechar{ε}{\ensuremath{\varepsilon}}
\newunicodechar{θ}{\ensuremath{\theta}}
\newunicodechar{λ}{\ensuremath{\lambda}}
\newunicodechar{μ}{\ensuremath{\mu}}
\newunicodechar{σ}{\ensuremath{\sigma}}
\newunicodechar{τ}{\ensuremath{\tau}}
\newunicodechar{φ}{\ensuremath{\varphi}}
\newunicodechar{ψ}{\ensuremath{\psi}}
\newunicodechar{ω}{\ensuremath{\omega}}
\newunicodechar{Σ}{\ensuremath{\Sigma}}
% majuscules produites par \MakeUppercase dans les titres (α-aminés -> Α)
\newunicodechar{Α}{\ensuremath{\alpha}}
\newunicodechar{Β}{\ensuremath{\beta}}
\newunicodechar{Γ}{\ensuremath{\Gamma}}
\newunicodechar{Δ}{\ensuremath{\Delta}}
\newunicodechar{Ε}{\ensuremath{\varepsilon}}
\newunicodechar{Θ}{\ensuremath{\Theta}}
\newunicodechar{Λ}{\ensuremath{\Lambda}}
\newunicodechar{Μ}{\ensuremath{\mu}}
\newunicodechar{Τ}{\ensuremath{\tau}}
\newunicodechar{Φ}{\ensuremath{\Phi}}
\newunicodechar{Ψ}{\ensuremath{\Psi}}
\newunicodechar{Ω}{\ensuremath{\Omega}}
\newunicodechar{↦}{\ensuremath{\mapsto}}
\newunicodechar{∈}{\ensuremath{\in}}
\newunicodechar{∉}{\ensuremath{\notin}}
\newunicodechar{∧}{\ensuremath{\wedge}}
\newunicodechar{⇔}{\ensuremath{\Leftrightarrow}}
\newunicodechar{⟺}{\ensuremath{\Longleftrightarrow}}
\newunicodechar{⇒}{\ensuremath{\Rightarrow}}
\newunicodechar{⟹}{\ensuremath{\Longrightarrow}}
\newunicodechar{∘}{\ensuremath{\circ}}
\newunicodechar{∪}{\ensuremath{\cup}}
\newunicodechar{∩}{\ensuremath{\cap}}
\newunicodechar{≡}{\ensuremath{\equiv}}
\newunicodechar{⊂}{\ensuremath{\subset}}
\newunicodechar{⊥}{\ensuremath{\perp}}
\newunicodechar{∃}{\ensuremath{\exists}}
\newunicodechar{∀}{\ensuremath{\forall}}
\newunicodechar{∛}{\ensuremath{\sqrt[3]{\,}}}
\newunicodechar{∜}{\ensuremath{\sqrt[4]{\,}}}
\newunicodechar{⃗}{\ensuremath{{}^{\to}}}
\newunicodechar{ⁿ}{\ifmmode^{n}\else\textsuperscript{\ensuremath{n}}\fi}
\newunicodechar{ˣ}{\ifmmode^{x}\else\textsuperscript{\ensuremath{x}}\fi}
\newunicodechar{ᵇ}{\ifmmode^{b}\else\textsuperscript{\ensuremath{b}}\fi}
\newunicodechar{ʳ}{\ifmmode^{r}\else\textsuperscript{\ensuremath{r}}\fi}
\newunicodechar{ᵘ}{\ifmmode^{u}\else\textsuperscript{\ensuremath{u}}\fi}
\newunicodechar{ʸ}{\ifmmode^{y}\else\textsuperscript{\ensuremath{y}}\fi}
\newunicodechar{ᵃ}{\ifmmode^{a}\else\textsuperscript{\ensuremath{a}}\fi}
\newunicodechar{ᵉ}{\ifmmode^{e}\else\textsuperscript{\ensuremath{e}}\fi}
\newunicodechar{ᵏ}{\ifmmode^{k}\else\textsuperscript{\ensuremath{k}}\fi}
\newunicodechar{ᵖ}{\ifmmode^{p}\else\textsuperscript{\ensuremath{p}}\fi}
\newunicodechar{ᴺ}{\ifmmode^{N}\else\textsuperscript{\ensuremath{N}}\fi}
\newunicodechar{ᶜ}{\ifmmode^{c}\else\textsuperscript{\ensuremath{c}}\fi}
\newunicodechar{ʰ}{\ifmmode^{h}\else\textsuperscript{\ensuremath{h}}\fi}
\newunicodechar{ᵠ}{\ifmmode^{q}\else\textsuperscript{\ensuremath{q}}\fi}
\newunicodechar{ₙ}{\ifmmode_{n}\else\textsubscript{\ensuremath{n}}\fi}
\newunicodechar{ₐ}{\ifmmode_{a}\else\textsubscript{\ensuremath{a}}\fi}
\newunicodechar{ᵢ}{\ifmmode_{i}\else\textsubscript{\ensuremath{i}}\fi}
\newunicodechar{ᵣ}{\ifmmode_{r}\else\textsubscript{\ensuremath{r}}\fi}
\newunicodechar{ₖ}{\ifmmode_{k}\else\textsubscript{\ensuremath{k}}\fi}
\newunicodechar{ᵦ}{\ifmmode_{\beta}\else\textsubscript{\ensuremath{\beta}}\fi}
\newunicodechar{ₓ}{\ifmmode_{x}\else\textsubscript{\ensuremath{x}}\fi}
\newfontfamily\popdisplay{ArchivoBlack-Regular.ttf}[Path=../../../fonts/]
\newfontfamily\poplogo{SpaceGrotesk-Bold.ttf}[Path=../../../fonts/]
\newfontfamily\popsemi{PlusJakartaSans-SemiBold.ttf}[Path=../../../fonts/]
\newfontfamily\popxbold{PlusJakartaSans-ExtraBold.ttf}[Path=../../../fonts/]
\newfontfamily\popxboldls{PlusJakartaSans-ExtraBold.ttf}[Path=../../../fonts/,LetterSpace=3.0]

\setlist[enumerate]{leftmargin=1.7em,itemsep=5pt,topsep=4pt}
\setlist[itemize]{leftmargin=1.7em,itemsep=4pt,topsep=4pt,label={\color{poporange}\textbullet}}
\setlength{\parindent}{0pt}
\setlength{\parskip}{5pt}
\renewcommand{\arraystretch}{1.25}

% pgfplots : style Pop par défaut
\pgfplotsset{
  every axis/.append style={axis line style={popink,line width=1pt},tick style={popink},
    grid style={popline,line width=0.5pt},label style={font=\small},tick label style={font=\footnotesize}},
  popcurve/.style={poporange,line width=1.8pt,no markers,smooth},
}
% Même style disponible dans un \draw TikZ simple (hors pgfplots)
\tikzset{popcurve/.style={poporange,line width=1.8pt}}
\tikzset{popgrid/.style={popline,very thin}}
\colorlet{popcurve}{poporange} % le nom du style est parfois utilisé comme couleur

% ============================================================
% Pill / squiggle
% ============================================================
\newcommand{\poppill}[3]{%
  \tikz[baseline=-0.55ex]\node[draw=popink,line width=1.2pt,rounded corners=2.6pt,%
    fill=#2,inner xsep=6.5pt,inner ysep=3pt]{%
    \popxboldls\fontsize{7}{7}\selectfont\color{#3}\MakeUppercase{#1}};%
}
\newcommand{\popsquiggle}[1]{%
  \tikz[baseline=-0.4ex]{
    \draw[#1,line width=2.6pt,line cap=round]
      plot [smooth] coordinates {(0,0.09) (0.34,0.01) (0.68,0.21) (1.02,0.01) (1.36,0.10)};
  }%
}
% Mot-clé surligné (marqueur orange sous le texte)
\newcommand{\cle}[1]{{\popxbold\color{popdark}#1}}

% ============================================================
% En-tête / pied
% ============================================================
\pagestyle{fancy}
\fancyhf{}
\renewcommand{\headrulewidth}{0pt}
\renewcommand{\footrulewidth}{0pt}
\lhead{\raisebox{-0.15ex}{\poplogo\fontsize{13}{13}\selectfont\color{popink}OnBuch\color{poporange}+}}
\rhead{\poppill{\DOCMATIERE\ \textbullet\ \DOCNIVEAU\ \textbullet\ Leçon \DOCLECON}{white}{popink}}
\lfoot{\popxbold\fontsize{7.6}{7.6}\selectfont\color{popmuted}\MakeUppercase{Cours signé OnBuch+}\ \textbullet\ {\popsemi\DOCTITRE}}
\rfoot{\tikz[baseline=-0.6ex]\node[rounded corners=8.5pt,fill=popink,minimum height=17pt,minimum width=17pt,inner xsep=5pt,inner sep=0]{\popxbold\fontsize{8}{8}\selectfont\color{popcream}\thepage\,/\,\pageref*{LastPage}};}

% ============================================================
% Titres
% ============================================================
\newcommand{\chaptitre}[2]{%
  \begin{tcolorbox}[enhanced,colback=poporange,colframe=popink,boxrule=2.6pt,arc=11pt,
      drop shadow=popink,left=15pt,right=15pt,top=12pt,bottom=11pt,before skip=3pt,after skip=18pt]
    \poppill{#2}{popink}{popcream}\par\vspace{9pt}
    {\raggedright\hyphenpenalty=10000\exhyphenpenalty=10000\popdisplay\fontsize{22}{24}\selectfont\color{popcream}\MakeUppercase{#1}\par}
  \end{tcolorbox}
}
% Section numérotée : pastille orange avec le numéro + titre Archivo
\newcounter{popsec}
\newcommand{\coursec}[1]{%
  \stepcounter{popsec}\setcounter{popsub}{0}%
  \par\vspace{16pt}\noindent
  \begin{minipage}{\linewidth}
  \tikz[baseline=-0.7ex]\node[draw=popink,line width=1.6pt,fill=poporange,rounded corners=4pt,minimum size=22pt,inner sep=0]{\popdisplay\fontsize{12}{12}\selectfont\color{popcream}\thepopsec};%
  \hspace{8pt}{\popdisplay\fontsize{15}{17}\selectfont\color{popink}\MakeUppercase{#1}}\par
  \vspace{4pt}\noindent{\color{poporange}\rule{36pt}{3.6pt}}
  \end{minipage}\par\nopagebreak\vspace{8pt}
}
\newcounter{popsub}
\newcommand{\courssub}[1]{%
  \stepcounter{popsub}%
  \par\vspace{9pt}\noindent{\popxbold\fontsize{11.5}{13}\selectfont\color{popink}{\color{poporange}\thepopsec.\thepopsub}\hspace{6pt}#1}\par\nopagebreak\vspace{4pt}
}

% ============================================================
% Boîtes pédagogiques — même recette : fond blanc, bordure épaisse colorée,
% ombre dure encre, badge plein. Titre optionnel [#1].
% ============================================================
\tcbset{popbase/.style={breakable,enhanced,colback=white,boxrule=2.3pt,arc=9pt,
  shadow={3pt}{-3pt}{0pt}{popink},left=14pt,right=14pt,top=11pt,bottom=11pt,before skip=11pt,after skip=15pt,
  fontupper=\popsemi\fontsize{10.6}{14.5}\selectfont\color{popink},
  fontlower=\popsemi\fontsize{10.6}{14.5}\selectfont\color{popink}}}
% Coche dessinée (le glyphe ✓ n'existe pas dans les polices du document)
\newcommand{\popcheck}[1]{\tikz[baseline=-0.5ex]\draw[#1,line width=1.6pt,line cap=round,line join=round](-0.13,0.0)--(-0.04,-0.09)--(0.13,0.1);}
\providecommand{\coche}{\popcheck{popgreen}}
\providecommand{\resultat}[1]{\par\smallskip\noindent\fcolorbox{popgreen}{popgreenL}{\parbox{\dimexpr\linewidth-2\fboxsep-2\fboxrule\relax}{\textbf{Résultat.} #1}}\par\smallskip}
\newcommand{\popboxhead}[3]{\poppill{#1}{#2}{white}\ifx\relax#3\relax\else\hspace{6pt}{\popxbold\fontsize{10.6}{12}\selectfont\color{popink}#3}\fi\par\vspace{7pt}}

\newtcolorbox{aretenir}[1][]{popbase,colframe=popgreen,before upper={\popboxhead{À retenir}{popgreen}{#1}}}
% Texte en allemand (document de lecture, dialogue, extrait) : cadre
% d'étude. Un titre optionnel (source, auteur) s'affiche dans l'en-tête.
\newtcolorbox{textebox}[1][]{popbase,colframe=popblue,before upper={\popboxhead{Text}{popblue}{#1}\begin{otherlanguage}{ngerman}},after upper={\end{otherlanguage}}}
% Marques ✓ / ✗ (forme correcte / fautive) souvent utilisées par les modèles
\providecommand{\cross}{\textcolor{poppink}{\ensuremath{\times}}}
\providecommand{\xmark}{\cross}\providecommand{\crossmark}{\cross}\providecommand{\cmark}{\ensuremath{\checkmark}}
% Ligne de réponse / justification à compléter (inventée par les modèles)
\providecommand{\justification}[1]{\par\noindent\textit{Justification~:} \dotfill\par}
% \textipa (package tipa non chargé) : la police ne couvre pas l'API -> simple passage
\providecommand{\textipa}[1]{#1}
% Soulignement (ulem non chargé) : \uline, \ul
\providecommand{\uline}[1]{\underline{#1}}\providecommand{\ul}[1]{\underline{#1}}
% \glossaire{\item ...} inventée par les modèles : liste de glossaire
\providecommand{\glossaire}[1]{\begin{itemize}#1\end{itemize}}
% \alert (beamer) inventée par les modèles : texte mis en évidence
\providecommand{\alert}[1]{\textbf{\textcolor{poppink}{#1}}}
% variantes ASCII d'ordinaux écrites en macro
\providecommand{\iere}{\textsuperscript{re}}\providecommand{\ieme}{\textsuperscript{e}}\providecommand{\ere}{\textsuperscript{re}}\providecommand{\eme}{\textsuperscript{e}}\providecommand{\er}{\textsuperscript{er}}\providecommand{\ie}{\textsuperscript{e}}
% Ordinaux français écrits en macro par les modèles : 1\ière, 3\ième
\newcommand{\ière}{\textsuperscript{re}}\newcommand{\ième}{\textsuperscript{e}}
% \exemple (inventée par les modèles) : étiquette « Exemple : »
\providecommand{\exemple}{\textbf{Exemple~:}\ }
% \square utilisable aussi hors mode math (cases à cocher)
\AtBeginDocument{\let\squareMath\square\renewcommand{\square}{\ensuremath{\squareMath}}}
% Mot ou phrase en allemand à l'intérieur d'un paragraphe français
\newcommand{\all}[1]{\ifmmode\text{\textit{#1}}\else\begingroup\selectlanguage{ngerman}\itshape #1\endgroup\fi}
\let\esp\all % alias toléré
\newtcolorbox{exemplebox}[1][]{popbase,colframe=poporange,before upper={\popboxhead{Exemple}{poporange}{#1}}}
\newtcolorbox{definition}[1][]{popbase,colframe=popblue,before upper={\popboxhead{Définition}{popblue}{#1}}}
\newtcolorbox{propriete}[1][]{popbase,colframe=poppurple,before upper={\popboxhead{Propriété}{poppurple}{#1}}}
\newtcolorbox{methode}[1][]{popbase,colframe=popgold,before upper={\popboxhead{Méthode}{popgold}{#1}}}
\newtcolorbox{attention}[1][]{popbase,colframe=poppink,before upper={\popboxhead{Attention}{poppink}{#1}}}
\newtcolorbox{experience}[1][]{popbase,colframe=popblue,colback=popblueL,before upper={\popboxhead{Expérience}{popblue}{#1}}}
% « Le savais-tu ? » : carte violette pleine (contexte, culture, Cameroun)
\newtcolorbox{savaistu}[1][]{popbase,colback=poppurple,colframe=popink,
  fontupper=\popsemi\fontsize{10.4}{14.2}\selectfont\color{white},
  before upper={\poppill{Le savais-tu\,?}{white}{poppurple}\ifx\relax#1\relax\else\hspace{6pt}{\popxbold\color{white}#1}\fi\par\vspace{7pt}}}
% Exercice résolu : énoncé en haut, \tcblower puis la solution sur fond crème
\newcounter{popexo}\newcounter{popexr}
\newtcolorbox{exoresolu}[1][]{popbase,colframe=poporange,colbacklower=popcream,
  segmentation style={popink,line width=1.2pt,dashed},
  before upper={\stepcounter{popexr}\popboxhead{Exercice résolu \thepopexr}{poporange}{#1}},
  before lower={\poppill{Solution}{popink}{popcream}\par\vspace{6pt}}}
% Exercice (non résolu en place) — la correction est dans la partie Corrigés
\newtcolorbox{exercice}[1][]{popbase,colframe=popink,
  before upper={\stepcounter{popexo}\popboxhead{Exercice \thepopexo}{popink}{#1}}}
% Corrigé
\newtcolorbox{corrige}[1][]{popbase,colframe=popgreen,colback=popgreenL,
  before upper={\popboxhead{Corrigé}{popgreen}{#1}}}
% Objectifs / prérequis / bilan
\newtcolorbox{objectifs}{popbase,colframe=popink,colback=poporangeL,
  before upper={\popboxhead{Objectifs de la leçon}{popink}{}}}
\newtcolorbox{prerequis}{popbase,colframe=popink,colback=popblueL,
  before upper={\popboxhead{Prérequis}{popblue}{}}}
\newtcolorbox{bilan}{popbase,colframe=popink,colback=white,boxrule=2.8pt,
  before upper={\popboxhead{Fiche bilan}{poporange}{L'essentiel à connaître pour l'examen}}}
% Figure encadrée avec légende
\newtcolorbox{popfigure}[1][]{enhanced,breakable=false,colback=white,colframe=popline,boxrule=1.4pt,arc=8pt,
  left=10pt,right=10pt,top=10pt,bottom=8pt,before skip=10pt,after skip=12pt,halign=center,
  lower separated=false,halign lower=center,
  fontlower=\popsemi\fontsize{9}{11.5}\selectfont\color{popmuted},
  #1}
\newcommand{\legende}[1]{\par\vspace{6pt}{\popsemi\fontsize{9}{11.5}\selectfont\color{popmuted}#1}}

% ============================================================
% Signature OnBuch+
% ============================================================
\newcommand{\poplogoinline}[1]{{\poplogo\fontsize{#1}{#1}\selectfont\color{popink}OnBuch\color{poporange}+}}
\newcommand{\signatureonbuch}{%
  \begin{tcolorbox}[enhanced,colback=popink,colframe=popink,boxrule=2.2pt,arc=10pt,
      drop shadow=poporange,left=14pt,right=14pt,top=10pt,bottom=10pt,before skip=0pt,after skip=0pt]
    \tikz[baseline=-0.7ex]\node[circle,fill=popgreen,draw=popcream,line width=1.3pt,minimum size=22pt,inner sep=0]{\popcheck{white}};%
    \hspace{9pt}\begin{minipage}[c]{0.62\linewidth}
      {\popxbold\fontsize{12}{13}\selectfont\color{popcream}Signé\ \textbullet\ {\poplogo OnBuch\color{poporange}+}}\par\vspace{2pt}
      {\raggedright\popsemi\fontsize{8}{10.5}\selectfont\color{popline}\MakeUppercase{Équipe pédagogique OnBuch+}\\\MakeUppercase{Conforme au programme officiel MINESEC}\par}
    \end{minipage}\hfill
    {\poplogo\fontsize{22}{22}\selectfont\color{popcream}OnBuch\color{poporange}+}
  \end{tcolorbox}%
}

% ============================================================
% Page de garde
% #1 titre · #2 matière · #3 niveau · #4 module · #5 n° leçon · #6 durée ·
% #7 description / objectifs (texte ou itemize)
% ============================================================
\newcommand{\pagegarde}[7]{%
  \thispagestyle{empty}
  \newgeometry{a4paper,margin=1.7cm,top=1.6cm,bottom=1.4cm}%
  \begin{tikzpicture}[remember picture,overlay]
    \fill[poporange!18] ([xshift=-3.2cm,yshift=-3.6cm]current page.north east) circle (4.6cm);
    \fill[poppurple!14] ([xshift=2.4cm,yshift=7.6cm]current page.south west) circle (3.8cm);
  \end{tikzpicture}%
  {\poplogo\fontsize{30}{30}\selectfont\color{popink}OnBuch\color{poporange}+}\hfill
  \raisebox{0.6ex}{\poppill{Cours signé \textbullet\ Premium}{popink}{popcream}}\par
  \vspace{1pt}\popsquiggle{poppurple}\par
  \vspace{8pt}
  \begin{tcolorbox}[enhanced,colback=poporange,colframe=popink,boxrule=2.8pt,arc=13pt,
      drop shadow=popink,left=18pt,right=18pt,top=12pt,bottom=13pt,after skip=10pt]
    \poppill{#2 \textbullet\ #3}{popink}{popcream}\hspace{5pt}\poppill{#4}{white}{popink}\par\vspace{8pt}
    {\popdisplay\fontsize{14}{14}\selectfont\color{popink}LEÇON #5}\par\vspace{4pt}
    {\raggedright\hyphenpenalty=10000\exhyphenpenalty=10000\popdisplay\fontsize{25}{27}\selectfont\color{popcream}\MakeUppercase{#1}\par}
    \vspace{8pt}
    {\popxbold\fontsize{9.5}{9.5}\selectfont\color{popink}\MakeUppercase{Durée conseillée : #6}}
  \end{tcolorbox}
  \begin{tcolorbox}[enhanced,colback=white,colframe=popink,boxrule=2.2pt,arc=10pt,
      drop shadow=popink,left=16pt,right=16pt,top=10pt,bottom=10pt,after skip=8pt]
    \poppill{Dans cette leçon}{poppurple}{white}\par\vspace{5pt}
    \popsemi\fontsize{9.3}{12.6}\selectfont\color{popink}#7
  \end{tcolorbox}
  \noindent\begin{minipage}[t]{0.487\linewidth}
    \begin{tcolorbox}[enhanced,colback=popgreen,colframe=popink,boxrule=2.2pt,arc=10pt,
        drop shadow=popink,left=11pt,right=11pt,top=10pt,bottom=10pt,height=3.1cm,valign=center]
      \tcbox[colback=white,colframe=popink,boxrule=1.3pt,arc=5pt,boxsep=0pt,nobeforeafter,
        left=3pt,right=3pt,top=3pt,bottom=3pt,valign=center]{\qrcode[height=1.9cm]{https://play.google.com/store/apps/details?id=com.luvvix.onbuch&hl=fr}}%
      \hspace{8pt}\begin{minipage}[c]{0.5\linewidth}
        {\popdisplay\fontsize{11}{12.5}\selectfont\color{white}\MakeUppercase{Télécharge OnBuch}}\par\vspace{4pt}
        {\raggedright\popsemi\fontsize{8.4}{11}\selectfont\color{white}Gratuit \textbullet\ Google Play\\Tuteur IA, annales, concours, résultats\par}
      \end{minipage}
    \end{tcolorbox}
  \end{minipage}\hfill
  \begin{minipage}[t]{0.487\linewidth}
    \begin{tcolorbox}[enhanced,colback=poppurple,colframe=popink,boxrule=2.2pt,arc=10pt,
        drop shadow=popink,left=11pt,right=11pt,top=10pt,bottom=10pt,height=3.1cm,valign=center]
      \tikz[baseline=-0.7ex]\node[circle,fill=white,draw=popink,line width=1.3pt,minimum size=1.6cm,inner sep=0]{\popdisplay\fontsize{24}{24}\selectfont\color{poppurple}+};%
      \hspace{8pt}\begin{minipage}[c]{0.6\linewidth}
        {\popdisplay\fontsize{11}{12.5}\selectfont\color{white}\MakeUppercase{Passe à OnBuch+}}\par\vspace{4pt}
        {\raggedright\popsemi\fontsize{8.4}{11}\selectfont\color{white}Tous les cours signés, Tuteur IA illimité, fiches PDF — onglet OnBuch+ de l'app\par}
      \end{minipage}
    \end{tcolorbox}
  \end{minipage}
  \vspace{8pt}
  \popcontenu
  \vfill
  \signatureonbuch
  \restoregeometry
  \newpage
}

% Rangée « Ce cours contient » (page de garde)
\newcommand{\poptile}[3]{% #1 couleur #2 gros texte #3 légende
  \begin{tcolorbox}[enhanced,colback=#1,colframe=popink,boxrule=2pt,arc=9pt,shadow={2.5pt}{-2.5pt}{0pt}{popink},
      left=6pt,right=6pt,top=9pt,bottom=9pt,halign=center,valign=center,height=1.75cm,width=\linewidth,nobeforeafter]
    {\popdisplay\fontsize{12.5}{13}\selectfont\color{white}\MakeUppercase{#2}}\par\vspace{3pt}
    {\centering\popsemi\fontsize{7.8}{9.5}\selectfont\color{white}#3\par}
  \end{tcolorbox}}
\newcommand{\popcontenu}{%
  \par\noindent{\popxboldls\fontsize{7.5}{7.5}\selectfont\color{popmuted}CE COURS SIGNÉ CONTIENT}\par\vspace{7pt}
  \noindent\begin{minipage}[t]{0.235\linewidth}\poptile{popblue}{Cours}{complet, illustré, pas à pas}\end{minipage}\hfill
  \begin{minipage}[t]{0.235\linewidth}\poptile{poporange}{Méthodes}{et exercices résolus}\end{minipage}\hfill
  \begin{minipage}[t]{0.235\linewidth}\poptile{poppink}{Exercices}{type examen + corrigés}\end{minipage}\hfill
  \begin{minipage}[t]{0.235\linewidth}\poptile{popgreen}{Fiche bilan}{l'essentiel à retenir}\end{minipage}\par}

% Sommaire
\newtcolorbox{sommaire}{enhanced,colback=white,colframe=popink,boxrule=2.2pt,arc=10pt,
  drop shadow=popink,left=18pt,right=18pt,top=13pt,bottom=13pt,before skip=6pt,after skip=20pt,
  before upper={\poppill{Sommaire}{poppurple}{white}\par\vspace{9pt}\popsemi\fontsize{10.6}{14.5}\selectfont\color{popink}}}

% Signature de fin de cours — poussée tout en bas de la dernière page
\newcommand{\findecours}{%
  \par\vfill\nopagebreak
  \begin{center}\popsquiggle{poporange}\end{center}\vspace{4pt}
  \signatureonbuch
}
\AtBeginDocument{\def\llbracket{\mathopen{[\mkern-3.5mu[}}\def\rrbracket{\mathclose{]\mkern-3.5mu]}}} % intervalles d'entiers (glyphe ⟦ absent de la police)
% Glyphes absents de Fira Math : redéfinis à partir de glyphes disponibles
\AtBeginDocument{%
  \def\setminus{\mathbin{\backslash}}\let\smallsetminus\setminus
  \def\vdots{\mathord{\vbox{\baselineskip3.2pt\lineskiplimit0pt\kern4pt\hbox{.}\hbox{.}\hbox{.}}}}%
  \def\ddots{\mathinner{\mkern1mu\raise6.5pt\hbox{.}\mkern2mu\raise3.6pt\hbox{.}\mkern2mu\raise0.7pt\hbox{.}\mkern1mu}}%
  \def\triangle{\mathord{\scalebox{0.9}{$\symup{\Delta}$}}}%
  \def\nabla{\mathord{\raisebox{\depth}{\scalebox{1}[-1]{$\symup{\Delta}$}}}}%
  \def\checkmark{\popcheck{popdark}}%
  \def\star{\mathbin{\ast}}\def\bigstar{\mathord{\ast}}%
  \def\diamond{\mathbin{\scalebox{0.6}{$\Diamond$}}}%
  \def\bigcup{\mathop{\mathchoice{\raisebox{-0.3em}{\scalebox{1.7}{$\cup$}}}{\scalebox{1.25}{$\cup$}}{\cup}{\cup}}\displaylimits}%
  \def\bigcap{\mathop{\mathchoice{\raisebox{-0.3em}{\scalebox{1.7}{$\cap$}}}{\scalebox{1.25}{$\cap$}}{\cap}{\cap}}\displaylimits}%
  \def\lesseqqgtr{\mathrel{\lessgtr}}\def\bigoplus{\mathop{\oplus}\displaylimits}%
}
\providecommand{\ssi}{si et seulement si} % abréviation fréquente
\AtBeginDocument{\providecommand{\wideparen}{\overparen}} % arc de cercle

%% FIGFIX-BEGIN v4 — correctifs globaux des figures (docs/onbuch-generation/tools/figfix)
\usepackage{adjustbox}\usepackage{etoolbox}
% toute figure TikZ/pgfplots plus large que la ligne est réduite à la largeur disponible
\BeforeBeginEnvironment{tikzpicture}{\begin{adjustbox}{max width=\linewidth}}
\AfterEndEnvironment{tikzpicture}{\end{adjustbox}}
% légendes pgfplots lisibles par-dessus les courbes
\pgfplotsset{every axis/.append style={legend style={fill=white,fill opacity=0.92,text opacity=1,draw=black!15,inner sep=3pt}}}
% étiquettes d'arêtes (syntaxe edge["..."]) avec fond blanc
\tikzset{every edge quotes/.append style={fill=white,inner sep=1.2pt}}
%% FIGFIX-END
\begin{document}

\pagegarde{La gestion de l'information}{Allemand}{Tle A}{Chapitre 4 : Média et communication}{AL16}{2 h}{Cette leçon propose la lecture et l'exploitation d'un texte journalistique allemand sur la gestion de l'information à l'ère d'Internet. Elle permet d'acquérir le vocabulaire des médias et de maîtriser les interrogatives indirectes et les verbes à complément prépositionnel.
\vspace{4pt}
\begin{multicols}{2}\small
\begin{itemize}
\item À la fin de cette leçon, je sais comprendre globalement puis en détail un texte allemand sur la gestion de l'information
\item À la fin de cette leçon, je sais employer correctement le vocabulaire des médias (die Information, die Nachricht, die Quelle, der Journalist, die Zeitung, das Internet) avec articles et pluriels
\item À la fin de cette leçon, je sais former et utiliser les interrogatives indirectes avec ob et les mots interrogatifs (w-Fragen)
\item À la fin de cette leçon, je sais conjuguer et employer les verbes à complément prépositionnel (warten auf, sich informieren über, berichten über...)
\item À la fin de cette leçon, je sais repérer dans un texte les verbes à préposition et leurs régimes (Akkusativ/Dativ)
\item À la fin de cette leçon, je sais répondre par écrit et oralement à des questions de compréhension
\end{itemize}\end{multicols}}

\begin{sommaire}
\begin{multicols}{2}
\begin{enumerate}
\item Texte support : Informationen im digitalen Zeitalter
\item Compréhension du texte
\item Lexique thématique : Medien und Information
\item Grammaire 1 : les interrogatives indirectes (ob, w-Fragen)
\item Grammaire 2 : les verbes à complément prépositionnel
\item Commentaire modèle et production guidée
\item Activité d'intégration
\item Méthodes et exercices résolus
\item Exercices
\item Corrigés des exercices
\item Fiche bilan
\end{enumerate}
\end{multicols}
\end{sommaire}

\begin{objectifs}
\begin{itemize}
\item À la fin de cette leçon, je sais comprendre globalement puis en détail un texte allemand sur la gestion de l'information
\item À la fin de cette leçon, je sais employer correctement le vocabulaire des médias (die Information, die Nachricht, die Quelle, der Journalist, die Zeitung, das Internet) avec articles et pluriels
\item À la fin de cette leçon, je sais former et utiliser les interrogatives indirectes avec ob et les mots interrogatifs (w-Fragen)
\item À la fin de cette leçon, je sais conjuguer et employer les verbes à complément prépositionnel (warten auf, sich informieren über, berichten über...)
\item À la fin de cette leçon, je sais repérer dans un texte les verbes à préposition et leurs régimes (Akkusativ/Dativ)
\item À la fin de cette leçon, je sais répondre par écrit et oralement à des questions de compréhension
\item À la fin de cette leçon, je sais rédiger un court résumé et un commentaire guidé d'un texte de presse
\item À la fin de cette leçon, je sais produire un court paragraphe d'opinion sur le rôle des médias au Cameroun et en Allemagne
\end{itemize}
\end{objectifs}

\begin{prerequis}
\begin{itemize}
\item Les phrases interrogatives directes (w-Fragen et questions fermées) vues en Seconde/Première
\item La déclinaison des articles et les cas (Nominativ, Akkusativ, Dativ)
\item La place du verbe en position finale dans les subordonnées (weil, dass)
\item Le vocabulaire de base de la communication (sagen, fragen, antworten, lesen)
\item Le Präsens et le Perfekt des verbes faibles et forts
\end{itemize}
\end{prerequis}

\newpage

\coursec{Texte support : Informationen im digitalen Zeitalter}

\courssub{Situation d'accroche et problématique}

Chaque matin à Yaoundé, des millions de Camerounais consultent leur téléphone avant même de prendre leur petit-déjeuner : WhatsApp, Facebook, TikTok, les sites d'information, CRTV en ligne... Une photo circule dans un groupe WhatsApp du quartier Bastos : « Le match du Cameroun est reporté ! » Vraie nouvelle ou rumeur ? Qui a écrit cette information ? Qui la contrôle ?

En Allemagne aussi, les jeunes s'informent surtout sur Internet, et les journalistes professionnels doivent \cle{vérifier} leurs \all{Quellen} (sources) avant de publier une \all{Nachricht} (une information). Le débat sur les fausses nouvelles — les \all{Falschmeldungen} ou \all{fake news} — est au cœur de la vie publique allemande.

\textbf{Problématique :} Comment distinguer une vraie information d'une fausse nouvelle ? Comment un journaliste allemand gère-t-il l'information ? C'est ce que nous allons découvrir à travers un texte de presse allemand, que nous lirons d'abord globalement, puis en détail.

\courssub{Méthode : comment lire un texte de presse allemand}

\begin{methode}[Lire un article de presse en quatre étapes]
\begin{enumerate}
\item \textbf{Lire le titre et la source.} Le titre annonce le thème ; la source (nom du journal, date, auteur) indique le type de texte et son degré de fiabilité.
\item \textbf{Lire une première fois sans dictionnaire.} Le but est de saisir le sens général, pas de comprendre chaque mot.
\item \textbf{Repérer les mots connus et les mots transparents.} Beaucoup de mots allemands ressemblent au français ou à l'anglais : \all{die Information}, \all{der Journalist}, \all{das Internet}, \all{der Artikel}, \all{kritisch}.
\item \textbf{Répondre aux questions globales puis détaillées.} D'abord \all{Wer? Was? Wo? Wann?} (qui ? quoi ? où ? quand ?), ensuite les idées principales de chaque paragraphe.
\end{enumerate}
\end{methode}

\begin{attention}[Piège du mot à mot]
Ne traduis jamais un texte mot à mot : l'allemand et le français n'ont pas le même ordre des mots. Attention aussi aux faux amis : \all{die Nachricht} signifie « la nouvelle, l'information » (à la radio, au journal) et non seulement « le message » au sens de SMS. Pour un SMS, on dit plutôt \all{die Mitteilung} ou \all{die SMS}. De même, \all{die Quelle} = « la source » (d'une information), mais aussi « la source » d'une rivière : le contexte décide.
\end{attention}

\courssub{Lecture globale : titre, source, type de texte}

Avant de lire, observons le document :
\begin{itemize}
\item \textbf{Titre :} \all{Informationen im digitalen Zeitalter} — « Les informations à l'ère numérique ». Le thème est donc la gestion de l'information à l'époque d'Internet.
\item \textbf{Source :} un article paru dans un journal allemand (\all{die Zeitung}), signé par un journaliste. C'est donc un \cle{article de presse}, texte informatif et argumentatif.
\item \textbf{Hypothèse de lecture :} le journaliste va probablement expliquer comment il travaille et mettre en garde contre les fausses informations.
\end{itemize}

\courssub{Le texte : Informationen im digitalen Zeitalter}

\begin{textebox}[Informationen im digitalen Zeitalter — Artikel aus einer deutschen Zeitung]
Heute bekommen wir jede Minute neue Nachrichten aus dem Internet. Aber sind alle Informationen richtig? Viele Journalisten fragen sich, ob ihre Quellen zuverlässig sind. Bevor sie einen Artikel schreiben, überprüfen sie die Fakten. Ein guter Journalist berichtet immer über die Wahrheit und wartet auf eine Bestätigung.

Ich arbeite seit zehn Jahren bei einer großen Zeitung in Hamburg. Jeden Tag erhalte ich Hunderte von Informationen: aus dem Internet, von Agenturen, von Lesern. Meine erste Frage ist immer: Woher kommt diese Nachricht? Dann suche ich nach einer zweiten Quelle. Erst wenn zwei unabhängige Quellen dieselbe Information bestätigen, schreibe ich meinen Artikel.

Das Internet hat unsere Arbeit verändert. Früher warteten die Leser auf die Zeitung am Morgen; heute lesen sie die Nachrichten sofort auf dem Handy. Das ist schnell, aber auch gefährlich. Leider verbreiten viele Menschen im Internet Falschmeldungen, ohne darüber nachzudenken. Manche wollen mit diesen „fake news“ Politikern schaden oder einfach Geld verdienen.

Deshalb ist es wichtig, dass Leser kritisch fragen: Wer hat diese Information geschrieben? Woher kommt sie? Gibt es andere Quellen? Gute Zeitungen und seriöse Internetseiten nennen immer ihre Quellen. Ich weiß nicht, ob das Problem der Falschmeldungen bald verschwinden wird. Aber ich glaube an die Kraft der Wahrheit — und an Leser, die nachdenken, bevor sie eine Nachricht weitergeben.
\end{textebox}

\textbf{Glossaire (mots difficiles) :}

\begin{center}
\begin{tabularx}{\linewidth}{|l|X|}
\hline
\rowcolor{popblueL} \textbf{Deutsch} & \textbf{Français} \\
\hline
die Quelle, -n & la source (d'une information) \\
\hline
überprüfen & vérifier, contrôler \\
\hline
zuverlässig & fiable, digne de confiance \\
\hline
die Bestätigung, -en & la confirmation \\
\hline
die Falschmeldung, -en & la fausse nouvelle, la fausse information \\
\hline
verbreiten & répandre, diffuser, propager \\
\hline
das Zeitalter, - & l'ère, l'époque (n) \\
\hline
die Agentur, -en & l'agence (de presse) (f) \\
\hline
schaden (+ dat.) & nuire à, faire du tort à \\
\hline
weitergeben (gibt weiter, gab weiter, hat weitergegeben) & transmettre, faire circuler \\
\hline
\end{tabularx}
\end{center}

\courssub{Première compréhension globale : Wer? Was? Wo? Wann?}

Répondons aux quatre questions fondamentales du journalisme (les fameuses « W-Fragen ») :

\begin{center}
\begin{tabularx}{\linewidth}{|l|X|X|}
\hline
\rowcolor{popblueL} \textbf{Frage} & \textbf{Antwort (Deutsch)} & \textbf{Réponse (français)} \\
\hline
Wer? & Ein deutscher Journalist, der bei einer Zeitung in Hamburg arbeitet. & Un journaliste allemand qui travaille pour un journal à Hambourg. \\
\hline
Was? & Er erklärt, wie er Informationen und Quellen überprüft. & Il explique comment il vérifie les informations et les sources. \\
\hline
Wo? & In Deutschland (Hamburg) und im Internet. & En Allemagne (Hambourg) et sur Internet. \\
\hline
Wann? & Heute, im digitalen Zeitalter. & Aujourd'hui, à l'ère numérique. \\
\hline
\end{tabularx}
\end{center}

On retient l'idée générale : \cle{un journaliste professionnel ne publie jamais une information sans l'avoir vérifiée}, et il invite les lecteurs à être critiques face aux fausses nouvelles sur Internet.

\courssub{Compréhension détaillée : les idées principales par paragraphe}

\begin{itemize}
\item \textbf{Paragraphe 1 :} Le problème posé — nous recevons des nouvelles en permanence, mais toutes ne sont pas vraies ; le journaliste vérifie les faits avant d'écrire. (\all{Bevor sie einen Artikel schreiben, überprüfen sie die Fakten.})
\item \textbf{Paragraphe 2 :} La méthode du journaliste — il se présente (dix ans d'expérience à Hambourg) et explique sa règle d'or : \cle{deux sources indépendantes} doivent confirmer l'information avant la publication.
\item \textbf{Paragraphe 3 :} Le rôle d'Internet — Internet a rendu l'information immédiate, mais aussi dangereuse : beaucoup de gens répandent des \all{Falschmeldungen} sans réfléchir, parfois pour nuire ou pour gagner de l'argent.
\item \textbf{Paragraphe 4 :} Le conseil au lecteur et la conclusion — il faut se poser des questions critiques (qui a écrit ? d'où vient l'information ?) ; l'auteur reste optimiste : il croit en la force de la vérité.
\end{itemize}

\courssub{Relevé du champ lexical des médias}

Cherchons dans le texte tous les mots qui appartiennent au champ lexical des médias et de l'information :

\begin{center}
\begin{tabularx}{\linewidth}{|X|X|X|}
\hline
\rowcolor{popblueL} \textbf{Substantive (noms)} & \textbf{Verben (verbes)} & \textbf{Adjektive (adjectifs)} \\
\hline
die Nachricht, die Information, die Quelle, der Journalist, der Artikel, die Zeitung, das Internet, das Handy, die Falschmeldung, die Wahrheit, die Agentur, der Leser & überprüfen, berichten, verbreiten, bestätigen, schreiben, verändern, weitergeben, nachdenken & richtig, zuverlässig, kritisch, gefährlich, seriös, unabhängig, schnell \\
\hline
\end{tabularx}
\end{center}

Ce relevé sera enrichi dans la section lexique (section 3). Pour l'instant, retenons la famille de mots : \all{die Nachricht} $\rightarrow$ \all{die Nachrichten} (les informations, les actualités) ; \all{berichten} $\rightarrow$ \all{der Bericht} (le reportage, le compte rendu).

\begin{popfigure}
\begin{tikzpicture}[mindmap, grow cyclic, every node/.style={concept, align=center, font=\small},
root concept/.append style={concept color=popblue, font=\bfseries},
level 1/.append style={level distance=3.2cm, sibling angle=72}]
\node[root concept, text=white]{die Information}
  child[concept color=poporange] { node[concept, align=center]{die Quellen\\ \emph{les sources}} }
  child[concept color=popgreen] { node[concept, align=center]{die Medien\\ Zeitung, Radio, TV} }
  child[concept color=poppurple] { node[concept, align=center]{der Journalist\\ \emph{il vérifie}} }
  child[concept color=popgold] { node[concept, align=center]{das Internet\\ Handy, sofort} }
  child[concept color=poppink] { node[concept, align=center]{fake news\\ Falschmeldungen} };
\end{tikzpicture}
\legende{Carte mentale : le champ lexical de l'information dans le texte.}
\end{popfigure}

\courssub{Relevé grammatical : préparation aux sections 4 et 5}

Le texte contient des structures que nous étudierons en grammaire. Repérons-les dès maintenant.

\textbf{1. Les interrogatives indirectes.} Une question directe comme \all{Sind die Quellen zuverlässig?} peut être intégrée dans une phrase : le verbe conjugué passe alors \cle{à la fin}.

\begin{itemize}
\item \all{Viele Journalisten fragen sich, \textbf{ob} ihre Quellen zuverlässig \textbf{sind}.} — Beaucoup de journalistes se demandent si leurs sources sont fiables.
\item \all{Ich weiß nicht, \textbf{ob} das Problem bald verschwinden \textbf{wird}.} — Je ne sais pas si le problème disparaîtra bientôt.
\item \all{Meine erste Frage ist immer: \textbf{Woher kommt} diese Nachricht?} (question directe) $\rightarrow$ \all{Ich frage, \textbf{woher} diese Nachricht \textbf{kommt}.} (indirecte) — Je demande d'où vient cette nouvelle.
\end{itemize}

Règle à retenir : \all{ob} introduit une question indirecte fermée (oui/non), et les mots interrogatifs en \all{w-} (\all{wer, was, woher, wann, warum}) introduisent une question indirecte ouverte ; dans les deux cas, le verbe conjugué se place en \cle{dernière position}.

\textbf{2. Les verbes à complément prépositionnel.} Certains verbes allemands se construisent avec une préposition fixe, comme en français « penser à », « croire en » :

\begin{itemize}
\item \all{Ein guter Journalist \textbf{berichtet über} die Wahrheit.} — Un bon journaliste rend compte de la vérité. (\all{berichten über} + accusatif)
\item \all{Er \textbf{wartet auf} eine Bestätigung.} — Il attend une confirmation. (\all{warten auf} + accusatif)
\item \all{Viele Menschen verbreiten Falschmeldungen, \textbf{ohne darüber nachzudenken}.} — ...sans y réfléchir. (\all{nachdenken über} + accusatif)
\item \all{Ich \textbf{glaube an} die Kraft der Wahrheit.} — Je crois en la force de la vérité. (\all{glauben an} + accusatif)
\item \all{Die Leser \textbf{warten auf} die Zeitung.} — Les lecteurs attendent le journal.
\end{itemize}

Note bien la forme \all{darüber} : quand le complément prépositionnel désigne une chose (et non une personne), la préposition se combine avec \all{da(r)-} : \all{darüber} = « à ce sujet, là-dessus ».

\begin{exoresolu}[Compréhension globale du texte]
\textbf{Énoncé.} Réponds en allemand, par une phrase complète, aux questions suivantes :
\begin{enumerate}
\item Wo arbeitet der Journalist?
\item Was macht ein guter Journalist, bevor er einen Artikel schreibt?
\item Warum ist das Internet gefährlich?
\end{enumerate}
\tcblower
\textbf{Solution détaillée.}
\begin{enumerate}
\item \all{Er arbeitet bei einer großen Zeitung in Hamburg.} — Il travaille pour un grand journal à Hambourg. (Reprise du paragraphe 2 ; attention : \all{bei} + datif pour « travailler chez/pour ».)
\item \all{Bevor er einen Artikel schreibt, überprüft er die Fakten und wartet auf eine Bestätigung.} — Avant d'écrire un article, il vérifie les faits et attend une confirmation. (Le verbe \all{überprüft} est placé juste après le sujet dans la proposition principale ; dans la subordonnée avec \all{bevor}, le verbe \all{schreibt} va à la fin.)
\item \all{Das Internet ist gefährlich, weil viele Menschen Falschmeldungen verbreiten, ohne darüber nachzudenken.} — Internet est dangereux parce que beaucoup de gens répandent de fausses nouvelles sans y réfléchir. (Subordonnée de cause avec \all{weil} : verbe à la fin.)
\end{enumerate}
\end{exoresolu}

\begin{attention}[Erreurs fréquentes des francophones]
\begin{itemize}
\item Ne dis pas \all{*Ich weiß nicht, ob das Problem wird verschwinden} : dans une interrogative indirecte, le verbe conjugué va \textbf{à la fin} : \all{..., ob das Problem bald verschwinden wird.}
\item \all{warten} se construit avec \all{auf} : on dit \all{Ich warte auf die Zeitung} (j'attends le journal), jamais \all{*Ich warte die Zeitung}.
\item N'oublie pas la majuscule à tous les noms : \all{die Zeitung, das Internet, die Wahrheit}.
\end{itemize}
\end{attention}

\begin{savaistu}[« Lügenpresse » et le débat allemand sur les fake news]
En Allemagne, le débat sur les fausses informations est très vif depuis 2016. Le mot anglais \all{fake news} est entré dans la langue courante, à côté du mot allemand \all{die Falschmeldung}. Le mot \all{Lügenpresse} (littéralement « presse mensongère », utilisé pour discréditer les médias) a même été élu \all{Unwort des Jahres} — « le pire mot de l'année » — par un jury de linguistes allemands. Chaque année, ce jury choisit un mot jugé contraire à la dignité humaine ou à la démocratie : une tradition qui montre combien les Allemands prennent la langue et l'information au sérieux. Et au Cameroun ? Les rumeurs sur WhatsApp avant un match des Lions Indomptables montrent que la vigilance face aux fausses nouvelles est un sujet universel !
\end{savaistu}

\begin{exercice}[À toi de jouer : repérage dans le texte]
\begin{enumerate}
\item Relève dans le texte trois noms du champ lexical des médias avec leur article et leur pluriel.
\item Souligne dans le texte deux interrogatives indirectes et recopie-les en les traduisant.
\item Trouve dans le texte trois verbes à complément prépositionnel et note leur préposition.
\item Réponds en allemand : \all{Was müssen Leser kritisch fragen?}
\end{enumerate}
\end{exercice}

\begin{corrige}[Exercice : repérage dans le texte]
\begin{enumerate}
\item Par exemple : \all{die Nachricht, -en} ; \all{die Zeitung, -en} ; \all{der Artikel, -} ; \all{die Quelle, -n}.
\item \all{Viele Journalisten fragen sich, ob ihre Quellen zuverlässig sind.} — Les journalistes se demandent si leurs sources sont fiables. \all{Ich weiß nicht, ob das Problem der Falschmeldungen bald verschwinden wird.} — Je ne sais pas si le problème des fausses nouvelles disparaîtra bientôt.
\item \all{berichten über} (+ acc.), \all{warten auf} (+ acc.), \all{nachdenken über} (+ acc.), \all{glauben an} (+ acc.).
\item \all{Leser müssen kritisch fragen: Wer hat diese Information geschrieben? Woher kommt sie? Gibt es andere Quellen?} — Les lecteurs doivent se demander : qui a écrit cette information ? D'où vient-elle ? Y a-t-il d'autres sources ?
\end{enumerate}
\end{corrige}

\begin{aretenir}[L'essentiel de cette section]
\begin{itemize}
\item Un article de presse se lit en quatre étapes : titre et source, lecture globale sans dictionnaire, mots transparents, questions globales puis détaillées.
\item Le texte présente le travail de vérification d'un journaliste allemand : deux sources indépendantes avant toute publication.
\item Champ lexical clé : \all{die Information, die Nachricht, die Quelle, der Journalist, die Zeitung, das Internet, die Falschmeldung}.
\item Le texte contient des \cle{interrogatives indirectes} (\all{ob} + verbe final) et des \cle{verbes à complément prépositionnel} (\all{warten auf, berichten über, nachdenken über, glauben an}), que nous étudierons en détail dans les sections 4 et 5.
\end{itemize}
\end{aretenir}

\coursec{Compréhension du texte}

Nous venons de lire et d'écouter le texte support \all{„Informationen im digitalen Zeitalter“} : nous savons désormais que les Allemands (comme les Camerounais) reçoivent leurs nouvelles surtout par Internet et les réseaux sociaux, que les journalistes vérifient leurs sources avant d'écrire un article, et que les \all{Falschmeldungen} (fausses informations) se diffusent très vite. Il s'agit maintenant de vérifier et d'approfondir cette compréhension, étape par étape : d'abord une compréhension \cle{globale} (le sens général), puis une compréhension \cle{détaillée} (les informations précises), en apprenant à \cle{justifier} ses réponses par une citation du texte et à \cle{reformuler} avec ses propres mots.

\begin{popfigure}
\begin{center}
\begin{tikzpicture}[
  node distance=0.55cm,
  etape/.style={rectangle, rounded corners=4pt, draw=popblue, fill=popblueL, thick, align=center, text width=4.6cm, minimum height=1.05cm, font=\small},
  fleche/.style={-{Stealth[length=3mm]}, thick, popdark}
]
\node[etape, align=center] (e1){\textbf{1. Lecture globale}\\ \emph{Worum geht es?}\\ De quoi parle le texte ?};
\node[etape, below=of e1, align=center] (e2){\textbf{2. Lecture détaillée}\\ \emph{Wer? Was? Wo? Warum?}\\ Qui ? Quoi ? Où ? Pourquoi ?};
\node[etape, below=of e2, align=center] (e3){\textbf{3. Justification}\\ \emph{Zitat, Zeile 3: „\ldots“}\\ Citer la ligne du texte};
\node[etape, below=of e3, align=center] (e4){\textbf{4. Reformulation}\\ \emph{mit eigenen Worten}\\ Dire autrement avec ses mots};
\draw[fleche] (e1) -- (e2);
\draw[fleche] (e2) -- (e3);
\draw[fleche] (e3) -- (e4);
\end{tikzpicture}
\end{center}
\legende{Les quatre étapes de la compréhension d'un texte : du sens général vers la reformulation.}
\end{popfigure}

\courssub{La méthode : comment répondre à une question de compréhension}

Avant de travailler sur les questions elles-mêmes, observons deux bonnes réponses données par des élèves de Terminale :

\begin{exemplebox}[Observer deux réponses modèles]
Question : \all{Woher bekommen viele Menschen heute ihre Nachrichten?} (D'où beaucoup de gens reçoivent-ils leurs nouvelles aujourd'hui ?)

Réponse d'élève : \all{Viele Menschen bekommen ihre Nachrichten heute aus dem Internet, besonders aus den sozialen Netzwerken.} « Beaucoup de gens reçoivent leurs nouvelles aujourd'hui d'Internet, surtout des réseaux sociaux. »

Question : \all{Was machen Journalisten, bevor sie einen Artikel schreiben?} (Que font les journalistes avant d'écrire un article ?)

Réponse d'élève : \all{Bevor sie einen Artikel schreiben, überprüfen die Journalisten ihre Quellen. Das steht in Zeile 8: „Der Journalist überprüft seine Quellen.“} « Avant d'écrire un article, les journalistes vérifient leurs sources. C'est écrit à la ligne 8 : „Le journaliste vérifie ses sources.“ »
\end{exemplebox}

Qu'observe-t-on ? La réponse \textbf{reprend les mots de la question} (\all{Menschen — bekommen — Nachrichten}), elle est rédigée en \textbf{phrase complète}, et elle peut s'appuyer sur une \textbf{citation} du texte avec le numéro de la ligne.

\begin{methode}[Répondre à une question de compréhension en 5 étapes]
\begin{enumerate}
\item \textbf{Lire la question} et souligner le mot interrogatif : \all{wer} (qui), \all{was} (quoi), \all{wo} (où), \all{wann} (quand), \all{warum} (pourquoi), \all{wie} (comment), \all{woher} (d'où).
\item \textbf{Repérer dans le texte} la phrase qui contient la réponse (chercher les mots-clés de la question ou leurs synonymes).
\item \textbf{Construire une phrase complète} qui reprend les mots de la question : on ne répond jamais par un seul mot.
\item \textbf{Justifier par une citation} : \all{Das steht in Zeile 5: „\ldots“} ou \all{Im Text steht: „\ldots“} (C'est écrit dans le texte : « … »).
\item \textbf{Vérifier la grammaire} : verbe conjugué en 2\up{e} position, majuscules aux noms, ponctuation.
\end{enumerate}
\end{methode}

\begin{attention}[La place du verbe dans la réponse]
En allemand, la réponse commence souvent par un complément (reprise de la question), mais le \textbf{verbe conjugué reste en 2\up{e} position} : \all{Bevor sie einen Artikel schreiben, \textbf{überprüfen} die Journalisten ihre Quellen.} Et surtout : ne jamais commencer une justification par \all{Weil\ldots} sans renvoyer le verbe conjugué \textbf{à la fin} de la subordonnée. On écrit : \all{Falschmeldungen sind gefährlich, \textbf{weil} viele Menschen sie \textbf{glauben}.} « Les fausses informations sont dangereuses parce que beaucoup de gens les croient. » (et non : *\all{weil viele Menschen glauben sie}).
\end{attention}

\courssub{Compréhension globale : choix multiple et vrai ou faux}

La compréhension globale répond à la question \all{Worum geht es im Text?} (De quoi parle le texte ?). Elle se travaille avec des questions à choix multiple (QCM) et des affirmations à valider.

\begin{exercice}[QCM — compréhension globale]
Entoure la bonne réponse (\all{Kreuze die richtige Antwort an}).

\textbf{1. Worum geht es im Text?}
\begin{itemize}
\item[a)] um den Sport in Kamerun
\item[b)] um die Information und die Medien im digitalen Zeitalter
\item[c)] um das Wetter in Deutschland
\end{itemize}

\textbf{2. Was ist das Hauptproblem des Textes?}
\begin{itemize}
\item[a)] Die Zeitungen sind zu teuer.
\item[b)] Es gibt zu viele Journalisten.
\item[c)] Falschmeldungen verbreiten sich sehr schnell im Internet.
\end{itemize}

\textbf{3. Der Text ist\ldots}
\begin{itemize}
\item[a)] ein Zeitungsartikel (un article de journal)
\item[b)] ein Gedicht (un poème)
\item[c)] ein Brief (une lettre)
\end{itemize}
\end{exercice}

\begin{corrige}[QCM — compréhension globale]
\textbf{1. b)} Le texte parle de l'information et des médias à l'ère numérique : \all{Es geht um die Information und die Medien im digitalen Zeitalter.}

\textbf{2. c)} Le problème principal est la diffusion rapide des fausses informations : \all{Falschmeldungen verbreiten sich sehr schnell im Internet.}

\textbf{3. a)} C'est un article de journal (\all{ein Zeitungsartikel}) : texte informatif, titre, paragraphes courts, ton objectif.
\end{corrige}

\begin{exercice}[Vrai ou faux justifié — \all{Richtig oder falsch?}]
Réponds par \all{richtig} (vrai) ou \all{falsch} (faux) et justifie avec une citation du texte : \all{Das steht in Zeile \ldots: „\ldots“}

\textbf{1.} \all{Alle Menschen lesen heute noch die Zeitung auf Papier.} (Tout le monde lit encore le journal papier.)

\textbf{2.} \all{Journalisten überprüfen ihre Quellen, bevor sie einen Artikel schreiben.} (Les journalistes vérifient leurs sources avant d'écrire un article.)

\textbf{3.} \all{Falschmeldungen verbreiten sich langsamer als gute Nachrichten.} (Les fausses informations se diffusent plus lentement que les bonnes nouvelles.)

\textbf{4.} \all{Man soll eine Nachricht erst lesen und prüfen, bevor man sie weiterschickt.} (Il faut lire et vérifier une nouvelle avant de la transférer.)
\end{exercice}

\begin{corrige}[Vrai ou faux justifié]
\textbf{1. Falsch.} \all{Das steht im Text: „Immer weniger Menschen lesen die Zeitung auf Papier.“} « De moins en moins de gens lisent le journal papier. »

\textbf{2. Richtig.} \all{Das steht in Zeile 8: „Der Journalist überprüft seine Quellen.“} « Le journaliste vérifie ses sources. »

\textbf{3. Falsch.} \all{Im Text steht: „Viele Menschen verbreiten Falschmeldungen, und sie verbreiten sich sehr schnell.“} « Beaucoup de gens diffusent de fausses informations, et elles se diffusent très vite. »

\textbf{4. Richtig.} \all{Der Text sagt: „Lies die Nachricht zuerst und prüfe die Quelle!“} « Lis d'abord la nouvelle et vérifie la source ! »
\end{corrige}

\courssub{Compréhension détaillée : répondre en phrases complètes}

La compréhension détaillée porte sur des informations précises : personnes, lieux, actions, raisons. Voici les questions types de cette leçon.

\begin{textebox}[Questions types sur le texte]
1) \all{Woher bekommen wir heute Nachrichten?}

2) \all{Was fragen sich viele Journalisten?}

3) \all{Was machen Journalisten, bevor sie einen Artikel schreiben?}

4) \all{Wer verbreitet Falschmeldungen?}
\end{textebox}

\begin{exoresolu}[Questions de compréhension détaillée]
Énoncé : réponds aux quatre questions types en phrases complètes, en allemand, avec une justification quand c'est possible.

\tcblower

\textbf{1.} \all{Woher bekommen wir heute Nachrichten?}

\all{Wir bekommen heute Nachrichten vor allem aus dem Internet, aus den sozialen Netzwerken, aus dem Radio und aus dem Fernsehen.} « Nous recevons aujourd'hui les nouvelles surtout d'Internet, des réseaux sociaux, de la radio et de la télévision. »

\textbf{2.} \all{Was fragen sich viele Journalisten?}

\all{Viele Journalisten fragen sich, wie sie die Leser im digitalen Zeitalter erreichen können.} « Beaucoup de journalistes se demandent comment ils peuvent atteindre les lecteurs à l'ère numérique. » (Remarque : la question \all{Was fragen sich\ldots?} appelle une réponse avec une \textbf{interrogative indirecte} : \all{wie sie\ldots erreichen können} — le verbe conjugué passe à la fin.)

\textbf{3.} \all{Was machen Journalisten, bevor sie einen Artikel schreiben?}

\all{Bevor sie einen Artikel schreiben, überprüfen die Journalisten ihre Quellen und recherchieren die Fakten. Das steht im Text: „Der Journalist überprüft seine Quellen.“} « Avant d'écrire un article, les journalistes vérifient leurs sources et recherchent les faits. »

\textbf{4.} \all{Wer verbreitet Falschmeldungen?}

\all{Viele Menschen verbreiten Falschmeldungen, oft ohne es zu wissen, weil sie die Nachrichten nicht prüfen.} « Beaucoup de gens diffusent de fausses informations, souvent sans le savoir, parce qu'ils ne vérifient pas les nouvelles. »
\end{exoresolu}

\begin{exemplebox}[Phrases clés du texte à retenir]
\all{Der Journalist überprüft seine Quellen.} « Le journaliste vérifie ses sources. »

\all{Viele Menschen verbreiten Falschmeldungen.} « Beaucoup de gens diffusent de fausses informations. »

\all{Die Nachrichten kommen heute sehr schnell aus dem Internet.} « Les nouvelles arrivent aujourd'hui très vite par Internet. »

\all{Man muss jede Information kritisch lesen.} « Il faut lire chaque information avec un œil critique. »
\end{exemplebox}

\courssub{Questions ouvertes : donner son opinion}

Les questions ouvertes (\all{offene Fragen}) n'ont pas une seule réponse : elles demandent ton avis, toujours en allemand et en phrases complètes. Deux questions sont au programme de cette leçon.

\textbf{1.} \all{Was macht ein guter Journalist?} (Que fait un bon journaliste ? / Quelles sont les qualités d'un bon journaliste ?)

Modèle de réponse : \all{Ein guter Journalist überprüft immer seine Quellen, er schreibt objektiv und er informiert die Menschen ehrlich. Er verbreitet keine Falschmeldungen.} « Un bon journaliste vérifie toujours ses sources, il écrit de manière objective et il informe les gens honnêtement. Il ne diffuse pas de fausses informations. »

\textbf{2.} \all{Warum sind Falschmeldungen gefährlich?} (Pourquoi les fausses informations sont-elles dangereuses ?)

Modèle de réponse : \all{Falschmeldungen sind gefährlich, weil viele Menschen sie glauben und weiterschicken. Sie können Angst machen und die Gesellschaft spalten.} « Les fausses informations sont dangereuses parce que beaucoup de gens les croient et les transfèrent. Elles peuvent faire peur et diviser la société. »

\begin{attention}[Exprimer une cause : attention à la place du verbe]
Pour justifier une opinion, on utilise \all{weil} (parce que) ou \all{denn} (car). Avec \all{weil}, le verbe conjugué va \textbf{à la fin} ; avec \all{denn}, il reste en 2\up{e} position :

\all{Falschmeldungen sind gefährlich, \textbf{weil} viele Menschen sie \textbf{glauben}.}

\all{Falschmeldungen sind gefährlich, \textbf{denn} viele Menschen \textbf{glauben} sie.}

Les deux phrases signifient : « Les fausses informations sont dangereuses parce que beaucoup de gens les croient. » Ne traduis jamais « parce que » mot à mot sans penser à la place du verbe.
\end{attention}

\courssub{La justification par la citation et la reformulation}

\textbf{La technique du \cle{Zitat} (citation).} Pour prouver qu'une réponse est correcte, on cite la ligne exacte du texte. Les formules à connaître :

\begin{center}
\begin{tabularx}{\linewidth}{|X|X|}
\hline
\rowcolor{popblueL} \textbf{Deutsch} & \textbf{Français} \\
\hline
\all{Das steht in Zeile 5: „\ldots“} & C'est écrit à la ligne 5 : « … » \\
\hline
\all{Im Text steht: „\ldots“} & Dans le texte, il est écrit : « … » \\
\hline
\all{Der Autor sagt/schreibt: „\ldots“} & L'auteur dit/écrit : « … » \\
\hline
\all{In Zeile 8 bis 10 lesen wir: „\ldots“} & Aux lignes 8 à 10, nous lisons : « … » \\
\hline
\end{tabularx}
\end{center}

\textbf{La reformulation} (\all{mit eigenen Worten}) consiste à dire la même chose autrement, avec des synonymes ou une autre structure. C'est une compétence essentielle pour le résumé et pour l'épreuve de commentaire au Baccalauréat.

\begin{center}
\begin{tabularx}{\linewidth}{|X|X|}
\hline
\rowcolor{popblueL} \textbf{Phrase du texte} & \textbf{Reformulation possible} \\
\hline
\all{Der Journalist überprüft seine Quellen.} & \all{Der Journalist kontrolliert, woher seine Informationen kommen.} \\
\hline
\all{Viele Menschen verbreiten Falschmeldungen.} & \all{Viele Leute schicken falsche Nachrichten weiter.} \\
\hline
\all{Immer weniger Menschen lesen die Zeitung auf Papier.} & \all{Die Papierzeitung hat immer weniger Leser.} \\
\hline
\end{tabularx}
\end{center}

\begin{exoresolu}[Reformuler une phrase du texte]
Énoncé : reformule la phrase \all{„Man muss jede Information kritisch lesen.“} avec tes propres mots, puis traduis ta reformulation en français.

\tcblower

Reformulation : \all{Jeder soll die Nachrichten prüfen und nicht alles sofort glauben.} « Chacun devrait vérifier les nouvelles et ne pas tout croire immédiatement. »

Technique utilisée : remplacement de \all{man muss} (il faut) par \all{jeder soll} (chacun devrait), et de \all{kritisch lesen} (lire avec un œil critique) par \all{prüfen und nicht alles glauben} (vérifier et ne pas tout croire). Le sens reste identique, les mots changent.
\end{exoresolu}

\courssub{Correction collective et entraînement oral en binôme}

La compréhension se travaille aussi à l'oral : savoir poser une question et y répondre spontanément est une compétence visée par le programme.

\begin{experience}[Questions-réponses en binôme — \all{Partnerarbeit}]
\begin{enumerate}
\item Formez des binômes (\all{Bildet Zweiergruppen}). L'élève A lit une des quatre questions types ; l'élève B répond en phrase complète, sans regarder le corrigé.
\item L'élève A vérifie la réponse et demande la justification : \all{Wo steht das im Text?} (Où est-ce écrit dans le texte ?) L'élève B cite la ligne : \all{Das steht in Zeile \ldots: „\ldots“}
\item Changez les rôles (\all{Wechselt die Rollen}).
\item Pour finir, chaque binôme pose une question ouverte à la classe : \all{Was macht ein guter Journalist?} ou \all{Warum sind Falschmeldungen gefährlich?} Deux ou trois élèves donnent leur avis ; le professeur corrige collectivement au tableau la place du verbe et le vocabulaire.
\end{enumerate}
\end{experience}

\begin{aretenir}[L'essentiel de la compréhension de texte]
\begin{itemize}
\item Répondre toujours en \textbf{phrase complète}, en reprenant les mots de la question.
\item Justifier par une \textbf{citation} : \all{Das steht in Zeile \ldots: „\ldots“}
\item Verbe conjugué en \textbf{2\up{e} position} dans la phrase principale ; à la \textbf{fin} après \all{weil} et dans les interrogatives indirectes.
\item Reformuler, c'est garder le sens et changer les mots (\all{überprüfen} $\rightarrow$ \all{kontrollieren} ; \all{verbreiten} $\rightarrow$ \all{weiterschicken}).
\item Pour les questions ouvertes, donner son avis avec \all{Ich denke\ldots} / \all{Meiner Meinung nach\ldots} (je pense / à mon avis) et justifier avec \all{weil} ou \all{denn}.
\end{itemize}
\end{aretenir}

Cette étape de compréhension étant acquise, nous pouvons maintenant enrichir notre lexique : la section suivante présente le vocabulaire thématique \all{Medien und Information}, avec articles, pluriels et familles de mots.

\coursec{Lexique thématique : Medien und Information}

Après avoir lu et compris le texte \all{Informationen im digitalen Zeitalter}, nous avons rencontré de nombreux mots nouveaux : \all{die Quelle}, \all{der Journalist}, \all{die Nachricht}, \all{verbreiten}... Cette section les organise, les explique et les fixe durablement. En allemand, un mot de vocabulaire ne s'apprend jamais seul : on l'apprend avec son \cle{article}, son \cle{pluriel} et sa \cle{famille de mots}. C'est la règle d'or de cette leçon.

\courssub{Observer d'abord : le vocabulaire dans son contexte}

\begin{textebox}[Mini-Texte : Ein Tag in der Redaktion]
Jeden Morgen um sieben Uhr liest der Redakteur Thomas Berger die wichtigsten Nachrichten im Internet und in der Zeitung. Eine Journalistin berichtet über eine Wahl in Kamerun, aber Thomas ist vorsichtig: Ist die Quelle zuverlässig? Er überprüft die Information und spricht mit einem Zeugen. Dann veröffentlicht die Zeitung den Artikel. Die Leser kommentieren die Meldung am Nachmittag in den sozialen Netzwerken. „Man muss jede Information prüfen, bevor man sie verbreitet“, sagt Thomas. „Ein Gerücht ist keine Tatsache!“
\end{textebox}

\textbf{Glossaire :} die Redaktion (-en) : la rédaction ; vorsichtig : prudent ; zuverlässig : fiable ; der Zeuge (-n) : le témoin ; prüfen : vérifier ; das Gerücht (-e) : la rumeur ; die Tatsache (-n) : le fait.

Ce petit texte contient presque tout le champ lexical de la leçon. Observons-le champ par champ.

\courssub{Champ 1 — L'information et ses nuances}

L'allemand distingue plusieurs mots là où le français utilise souvent le seul mot « information » ou « nouvelle ». Cette précision est une richesse : chaque mot a sa nuance.

\begin{definition}[Les mots de l'information]
\begin{itemize}
\item \all{die Information, -en} : l'information (renseignement, contenu factuel). \all{Ich brauche mehr Informationen.} « J'ai besoin de plus d'informations. »
\item \all{die Nachricht, -en} : la nouvelle, l'information (au journal, à la radio) ; au pluriel \all{die Nachrichten} = les informations, le journal télévisé. \all{Hast du die Nachrichten gesehen?} « As-tu regardé les informations ? »
\item \all{die Meldung, -en} : la dépêche, le communiqué (information brève diffusée par une agence). \all{Eine Meldung der Nachrichtenagentur.} « Une dépêche de l'agence de presse. »
\item \all{die Falschmeldung, -en} : la fausse information, la « fake news ». \all{Die Falschmeldung verbreitete sich schnell.} « La fausse information s'est répandue vite. »
\item \all{die Tatsache, -n} : le fait (vérifiable). \all{Das ist eine Tatsache.} « C'est un fait. »
\item \all{das Gerücht, -e} : la rumeur (non vérifiée). \all{Das ist nur ein Gerücht.} « Ce n'est qu'une rumeur. »
\end{itemize}
\end{definition}

\begin{attention}[Ne confondez pas !]
\begin{itemize}
\item \all{die Nachricht} est une nouvelle transmise par un média ; \all{die Information} est le renseignement lui-même. On dit \all{eine Nachricht im Radio hören} (écouter une nouvelle à la radio) mais \all{Informationen über ein Thema sammeln} (rassembler des informations sur un sujet).
\item Le français « une information » peut se traduire par \all{eine Information} ou \all{eine Nachricht} selon le contexte ; mais « les informations » (le journal de 20 h) se dit toujours \all{die Nachrichten}.
\item \all{die Meldung} est un faux ami partiel : ce n'est pas une « mélodie » mais une dépêche !
\end{itemize}
\end{attention}

\courssub{Champ 2 — Les sources et les acteurs de l'information}

\begin{definition}[die Quelle]
\all{die Quelle, -n} : la source (d'une information). \all{Der Journalist nennt seine Quelle nicht.} « Le journaliste ne révèle pas sa source. » Expression clé : \all{aus zuverlässiger Quelle} = « de source fiable ».
\end{definition}

\begin{attention}[Double sens de \all{die Quelle}]
\all{die Quelle} signifie aussi « la source » d'une rivière, d'un cours d'eau : \all{die Quelle des Flusses} « la source du fleuve ». Le contexte décide : dans un texte sur les médias, il s'agit presque toujours de la source d'une information.
\end{attention}

\begin{center}
\begin{tabularx}{\linewidth}{|X|X|X|}
\hline
\rowcolor{popblueL} Allemand (article, pluriel) & Français & Exemple \\
\hline
der Journalist, -en & le journaliste & \all{Der Journalist schreibt einen Artikel.} \\
\hline
die Journalistin, -nen & la journaliste & \all{Die Journalistin interviewt den Minister.} \\
\hline
der Redakteur, -e & le rédacteur & \all{Der Redakteur prüft den Text.} \\
\hline
der Leser, - & le lecteur & \all{Die Leser schreiben Briefe an die Zeitung.} \\
\hline
der Zeuge, -n & le témoin & \all{Der Zeuge beschreibt den Unfall.} \\
\hline
die Quelle, -n & la source & \all{Die Quelle ist zuverlässig.} \\
\hline
\end{tabularx}
\end{center}

\begin{attention}[Le nom faible \all{der Journalist}]
\all{der Journalist} appartient aux \cle{noms faibles} (masculins en \all{-ist}, \all{-ent}, \all{-e} désignant des personnes) : il prend un \textbf{-en} à tous les cas sauf au nominatif singulier.
\begin{itemize}
\item Nominatif : \all{Der Journalist fragt.} « Le journaliste pose une question. »
\item Accusatif : \all{Ich sehe den Journalisten.} « Je vois le journaliste. »
\item Datif : \all{Ich gebe dem Journalisten das Dokument.} « Je donne le document au journaliste. »
\end{itemize}
Même phénomène pour \all{der Zeuge} : \all{den Zeugen}, \all{dem Zeugen}. Le féminin, lui, est régulier : \all{die Journalistin, -nen} (suffixe \all{-in}, pluriel en \all{-nen}), comme \all{die Lehrerin}, \all{die Studentin}.
\end{attention}

\courssub{Champ 3 — Les médias}

\begin{center}
\begin{tabularx}{\linewidth}{|X|X|X|}
\hline
\rowcolor{popblueL} Allemand (article, pluriel) & Français & Exemple \\
\hline
die Zeitung, -en & le journal & \all{Die Zeitung erscheint jeden Tag.} \\
\hline
die Zeitschrift, -en & le magazine, la revue & \all{Diese Zeitschrift erscheint monatlich.} \\
\hline
das Radio, -s & la radio & \all{Ich höre morgens Radio.} \\
\hline
das Fernsehen (Sg.) & la télévision & \all{Das Fernsehen zeigt einen Film.} \\
\hline
das Internet (Sg.) & Internet & \all{Im Internet findet man alles.} \\
\hline
die Website, -s / die Webseite, -n & le site web & \all{Die Website der Schule ist neu.} \\
\hline
die sozialen Netzwerke (Pl.) & les réseaux sociaux & \all{Viele Jugendliche nutzen soziale Netzwerke.} \\
\hline
\end{tabularx}
\end{center}

\begin{propriete}[Les noms en -ung sont toujours féminins]
Tous les noms allemands se terminant en \all{-ung} sont féminins, sans exception : \all{die Zeitung}, \all{die Meldung}, \all{die Information} (comme tous les noms en \all{-tion}), \all{die Meinung}, \all{die Regierung}. Leur pluriel est en \all{-en} : \all{die Zeitungen}, \all{die Meldungen}. Astuce : si vous hésitez sur le genre d'un nom en \all{-ung}, dites \all{die} — vous aurez toujours raison.
\end{propriete}

\begin{savaistu}[D'où vient le mot \all{Zeitung} ?]
\all{Zeitung} vient de \all{die Zeit} (le temps) : le journal est littéralement ce qui rapporte « les choses du temps présent », les nouvelles de l'époque. On retrouve la même logique dans \all{die Zeitschrift} (la revue : \all{Schrift} = écrit) et dans le mot anglais \emph{news} lui-même. Au Cameroun, beaucoup d'élèves s'informent aujourd'hui moins par la \all{Zeitung} que par les \all{sozialen Netzwerke} — d'où l'importance de vérifier les \all{Quellen} !
\end{savaistu}

\courssub{Champ 4 — Les verbes de la communication}

\begin{center}
\begin{tabularx}{\linewidth}{|X|X|X|}
\hline
\rowcolor{popblueL} Verbe & Français & Construction et exemple \\
\hline
sich informieren (über + Akk.) & s'informer (sur) & \all{Ich informiere mich über die Nachrichten im Internet.} \\
\hline
berichten (über + Akk.) & rendre compte de, rapporter & \all{Die Zeitung berichtet über die Wahl.} \\
\hline
überprüfen & vérifier & \all{Man muss jede Information überprüfen.} \\
\hline
verbreiten & diffuser, répandre & \all{Er verbreitet eine Nachricht auf Facebook.} \\
\hline
veröffentlichen & publier & \all{Die Zeitung veröffentlicht den Artikel.} \\
\hline
lesen (liest, las, gelesen) & lire & \all{Sie liest die Zeitung am Morgen.} \\
\hline
fragen (nach + Dat.) & demander & \all{Der Leser fragt nach der Quelle.} \\
\hline
antworten (auf + Akk.) & répondre (à) & \all{Er antwortet auf die Frage.} \\
\hline
kommentieren & commenter & \all{Viele Leser kommentieren die Meldung.} \\
\hline
\end{tabularx}
\end{center}

\begin{exemplebox}[Exemples traduits]
\begin{itemize}
\item \all{Ich informiere mich über die Nachrichten im Internet.} « Je m'informe sur les nouvelles sur Internet. »
\item \all{Die Zeitung berichtet über die Wahl.} « Le journal rend compte de l'élection. »
\item \all{Der Journalist verbreitet eine Nachricht aus zuverlässiger Quelle.} « Le journaliste diffuse une nouvelle de source fiable. »
\item \all{Bevor du eine Meldung teilst, musst du sie überprüfen.} « Avant de partager une information, tu dois la vérifier. »
\item \all{Die Schüler kommentieren den Artikel in den sozialen Netzwerken.} « Les élèves commentent l'article sur les réseaux sociaux. »
\end{itemize}
\end{exemplebox}

\begin{attention}[Verbes à préposition fixe : ne pas traduire mot à mot]
Le français « s'informer \emph{de} quelque chose » devient \all{sich informieren \emph{über} + Akkusativ} ; « répondre \emph{à} une question » devient \all{antworten \emph{auf} + Akkusativ}. La préposition allemande ne coïncide presque jamais avec la française : apprenez chaque verbe avec sa préposition, comme un bloc. Nous reverrons ce point en grammaire (section 5).
\end{attention}

\courssub{Les familles de mots : un mot en cache trois}

En allemand, les familles de mots sont très régulières. Connaître un verbe, c'est souvent connaître aussi le nom et l'adjectif correspondants.

\begin{itemize}
\item \all{informieren} (informer) $\rightarrow$ \all{die Information} (l'information) $\rightarrow$ \all{informativ} (instructif) $\rightarrow$ \all{informiert} (informé : \all{gut informiert sein} « être bien informé »).
\item \all{berichten} (rapporter) $\rightarrow$ \all{der Bericht, -e} (le reportage, le rapport) : \all{ein Bericht über das Fußballspiel} « un reportage sur le match de football ».
\item \all{schreiben} (écrire) $\rightarrow$ \all{der Schriftsteller, -} (l'écrivain) / \all{der Artikel, -} (l'article) : \all{Der Schriftsteller schreibt einen Artikel.} « L'écrivain écrit un article. »
\end{itemize}

\begin{popfigure}
\begin{center}
\begin{tikzpicture}[
  root/.style={rectangle, rounded corners, draw=popblue, fill=popblueL, align=center, font=\bfseries, minimum width=3.2cm, minimum height=1cm},
  child/.style={rectangle, rounded corners, draw=popgreen, fill=popgreenL, align=center, minimum width=3cm, minimum height=0.9cm},
  link/.style={-{Stealth[length=3mm]}, thick, popdark}
]
\node[root, align=center] (v) at (0,0){informieren\\ (verbe : informer)};
\node[child, align=center] (n) at (-4.2,-2){die Information, -en\\ (nom : l'information)};
\node[child, align=center] (a) at (0,-2){informativ\\ (adjectif : instructif)};
\node[child, align=center] (p) at (4.2,-2){informiert\\ (participe : informé)};
\draw[link] (v) -- (n);
\draw[link] (v) -- (a);
\draw[link] (v) -- (p);
\end{tikzpicture}
\end{center}
\legende{Arbre de la famille de mots : informieren}
\end{popfigure}

\begin{popfigure}
\begin{center}
\begin{tikzpicture}[
  col/.style={rectangle, rounded corners, draw=popline, fill=popcream, align=center, minimum width=3.6cm, text width=3.4cm},
  head/.style={rectangle, rounded corners, draw=poporange, fill=poporangeL, align=center, font=\bfseries, minimum width=3.6cm, minimum height=0.8cm}
]
\node[head] (h1) at (0,0) {Information};
\node[col, below=0.15cm of h1, align=center]{die Information\\ die Nachricht\\ die Meldung\\ die Tatsache\\ das Gerücht};
\node[head] (h2) at (4.6,0) {Akteure};
\node[col, below=0.15cm of h2, align=center]{die Quelle\\ der Journalist\\ die Journalistin\\ der Redakteur\\ der Leser\\ der Zeuge};
\node[head] (h3) at (9.2,0) {Medien};
\node[col, below=0.15cm of h3, align=center]{die Zeitung\\ die Zeitschrift\\ das Radio\\ das Fernsehen\\ das Internet\\ die Website};
\node[head] (h4) at (13.8,0) {Verben};
\node[col, below=0.15cm of h4, align=center]{sich informieren\\ berichten\\ überprüfen\\ verbreiten\\ veröffentlichen\\ kommentieren};
\end{tikzpicture}
\end{center}
\legende{Carte mentale du champ lexical : Medien und Information}
\end{popfigure}

\courssub{Expressions utiles à mémoriser}

\begin{aretenir}[Expressions clés de la leçon]
\begin{itemize}
\item \all{eine Nachricht verbreiten} : diffuser une nouvelle.
\item \all{sich über etwas (Akk.) informieren} : s'informer sur quelque chose.
\item \all{aus zuverlässiger Quelle} : de source fiable.
\item \all{über etwas (Akk.) berichten} : rendre compte de quelque chose.
\item \all{eine Information überprüfen} : vérifier une information.
\item \all{einen Artikel veröffentlichen} : publier un article.
\item \all{die Nachrichten sehen / hören} : regarder / écouter les informations.
\end{itemize}
\end{aretenir}

\begin{exoresolu}[Exercice de fixation du lexique]
Énoncé — Complétez les phrases avec le mot qui convient (article inclus) : \all{Quelle}, \all{Journalistin}, \all{Nachrichten}, \all{Gerücht}, \all{überprüft}, \all{Zeitung}.
\begin{enumerate}
\item \all{Die \_\_\_ berichtet über das Fußballspiel in Douala.}
\item \all{Die \_\_\_ interviewt den Bürgermeister von Yaoundé.}
\item \all{Der Redakteur nennt seine \_\_\_ nicht.}
\item \all{Bevor er den Artikel veröffentlicht, \_\_\_ er jede Information.}
\item \all{Das ist nur ein \_\_\_, keine Tatsache!}
\item \all{Abends sehe ich die \_\_\_ im Fernsehen.}
\end{enumerate}
\tcblower
Solution détaillée :
\begin{enumerate}
\item \all{Die Zeitung berichtet über das Fußballspiel in Douala.} — C'est le journal qui « rend compte » (\all{berichten über}) d'un événement.
\item \all{Die Journalistin interviewt den Bürgermeister von Yaoundé.} — Le féminin régulier en \all{-in} ; article \all{die}.
\item \all{Der Redakteur nennt seine Quelle nicht.} — La source d'une information ; ici accusatif féminin (\all{seine Quelle}).
\item \all{Bevor er den Artikel veröffentlicht, überprüft er jede Information.} — Attention à la place du verbe : après la subordonnée en \all{bevor}, la principale commence par le verbe conjugué.
\item \all{Das ist nur ein Gerücht, keine Tatsache!} — \all{das Gerücht} est neutre : \all{ein Gerücht}.
\item \all{Abends sehe ich die Nachrichten im Fernsehen.} — « Les informations » au sens de journal télévisé = \all{die Nachrichten} (pluriel).
\end{enumerate}
\end{exoresolu}

\begin{attention}[Erreurs fréquentes des francophones]
\begin{itemize}
\item Genre : on écrit \all{die Information}, \all{die Zeitung}, \all{die Quelle} (féminins), mais \all{das Internet}, \all{das Radio}, \all{das Gerücht} (neutres). Apprenez toujours le mot avec son article.
\item Majuscule : tous ces noms prennent la majuscule, même au milieu de la phrase : \all{ich lese die Zeitung}, et non \all{die zeitung}.
\item Pluriel : \all{der Journalist} fait \all{die Journalisten} ; \all{die Journalistin} fait \all{die Journalistinnen} (double \all{-n-} !).
\item Ne dites pas \all{informieren über mich} pour « m'informer » : le verbe est réfléchi — \all{ich informiere mich über...}.
\end{itemize}
\end{attention}

Ce lexique est maintenant votre boîte à outils : il servira dans les sections suivantes pour construire des questions indirectes (\all{Ich weiß nicht, ob die Quelle zuverlässig ist...}) et pour manier les verbes à complément prépositionnel. Apprenez chaque mot avec son article, son pluriel et un exemple : c'est le prix d'un vocabulaire vraiment solide.

\coursec{Grammaire 1 : les interrogatives indirectes (ob, w-Fragen)}

Dans la section précédente, nous avons constitué notre boîte à outils lexicale autour des médias et de l'information : \all{die Information}, \all{die Nachricht}, \all{die Quelle}, \all{der Journalist}, \all{die Zeitung}, \all{das Internet}. Or, qui dit information dit questions : le journaliste vérifie ses sources, le lecteur se demande si une nouvelle est vraie, le rédacteur en chef demande d'où vient une information. En allemand, toutes ces questions rapportées obéissent à une construction très précise : \cle{l'interrogative indirecte}. C'est l'une des structures les plus fréquentes de la presse germanophone — et l'une des plus fautives chez les élèves francophones. Observons, réglons, entraînons-nous.

\courssub{Observation : repérer la structure dans le texte}

\begin{textebox}[Aus der Redaktion — Textbeobachtung]
\all{Viele Journalisten fragen sich, ob ihre Quellen zuverlässig sind. Bevor ein Artikel erscheint, will der Chefredakteur wissen, woher die Information kommt. Er fragt den Reporter, wann er das Interview geführt hat und wer die Fotos gemacht hat. Die junge Praktikantin weiß nicht, ob sie den Fehler in der Überschrift korrigieren soll. Sie überlegt, wie sie es dem Redakteur erklären kann. Die Leser fragen sich oft, warum so viele Nachrichten negativ sind, und sie wissen nicht immer, wem sie glauben können.}

\emph{Glossaire :} die Quelle, -n (la source) ; zuverlässig (fiable) ; der Chefredakteur, -e (le rédacteur en chef) ; erscheinen (paraître) ; das Interview führen (mener l'interview) ; die Praktikantin, -nen (la stagiaire) ; der Fehler, - (l'erreur, pl. die Fehler) ; die Überschrift, -en (le titre) ; korrigieren (corriger) ; überlegen (réfléchir) ; glauben + Dat. (croire).
\end{textebox}

Soulignons mentalement ce qui suit \all{fragen}, \all{wissen}, \all{überlegen} :

\begin{itemize}
\item \all{Viele Journalisten fragen sich, \textbf{ob} ihre Quellen zuverlässig \textbf{sind}.} — « Beaucoup de journalistes se demandent si leurs sources sont fiables. »
\item \all{... will der Chefredakteur wissen, \textbf{woher} die Information \textbf{kommt}.} — « ... le rédacteur en chef veut savoir d'où vient l'information. »
\item \all{Er fragt den Reporter, \textbf{wann} er das Interview \textbf{geführt hat}.} — « Il demande au reporter quand il a mené l'interview. »
\item \all{Die Leser fragen sich, \textbf{warum} so viele Nachrichten negativ \textbf{sind}.} — « Les lecteurs se demandent pourquoi tant de nouvelles sont négatives. »
\end{itemize}

\textbf{Constat :} dans chaque cas, la question rapportée est introduite par \all{ob} ou par un mot interrogatif en \all{w-} (\all{woher}, \all{wann}, \all{wer}, \all{warum}), et \textbf{le verbe conjugué est rejeté à la toute fin de la subordonnée}. C'est exactement le contraire du français, où l'ordre reste « normal » : « d'où l'information \emph{vient} ».

\begin{definition}[L'interrogative indirecte]
Une \cle{question indirecte} (ou interrogative indirecte) est une \textbf{subordonnée} qui rapporte une question sans la poser directement. Elle est introduite :
\begin{itemize}
\item par \all{ob} (« si ») quand la question directe attend une réponse \textbf{oui ou non} (question fermée) ;
\item par un \textbf{mot interrogatif} (\all{wer, was, wann, wo, woher, wohin, warum, wie, wie viele...}) quand la question directe commence par ce mot (question ouverte).
\end{itemize}
Comme toute subordonnée allemande, elle envoie le \textbf{verbe conjugué en dernière position}.
\end{definition}

\courssub{La règle de construction : le verbe à la fin}

\begin{propriete}[Position du verbe dans l'interrogative indirecte]
Dans l'interrogative indirecte, le verbe conjugué se place en \textbf{dernière position} de la subordonnée :\\
\all{Ich weiß nicht, ob er \textbf{kommt}.} = « Je ne sais pas s'il vient. »\\
\all{Sie fragt, wann die Zeitung \textbf{erscheint}.} = « Elle demande quand le journal paraît. »\\
Si le verbe est composé (Perfekt, Futur I) ou modal (+ infinitif), l'\textbf{auxiliaire conjugué} (ou le modal) va à la fin et le participe/infinitif se place juste devant lui :\\
\all{Wir wissen nicht, wer den Artikel \textbf{geschrieben hat}.} = « Nous ne savons pas qui a écrit l'article. »
\end{propriete}

\begin{popfigure}
\begin{tikzpicture}[
box/.style={draw=popblue, fill=popblueL, rounded corners=3pt, align=center, font=\small, inner sep=5pt},
box2/.style={draw=poporange, fill=poporangeL, rounded corners=3pt, align=center, font=\small, inner sep=5pt},
box3/.style={draw=popgreen, fill=popgreenL, rounded corners=3pt, align=center, font=\small, inner sep=5pt},
box4/.style={draw=poppurple, fill=poppinkL, rounded corners=3pt, align=center, font=\small, inner sep=5pt},
lbl/.style={font=\footnotesize\itshape, text=popmuted}]
\node[box, align=center] (h) at (0,0){Hauptsatz\\ \all{Er fragt (sich),}\\ \all{Sie weiß,} ...};
\node[box2, align=center] (e) at (4.6,0){Einleitung\\ \all{ob} ou mot en \all{w-}\\ \all{woher, wann, wer...}};
\node[box3, align=center] (s) at (8.6,0){Sujet +\\ compléments\\ \all{die Information}};
\node[box4, align=center] (v) at (12,0){VERBE\\ conjugué\\ \all{kommt.}};
\draw[-{Stealth}, thick, popdark] (h) -- (e);
\draw[-{Stealth}, thick, popdark] (e) -- (s);
\draw[-{Stealth}, thick, popdark] (s) -- (v);
\node[lbl] at (2.3,0.95) {proposition principale};
\node[lbl] at (10.3,0.95) {position finale obligatoire};
\draw[poppurple, thick] (4,-0.95) -- (12.6,-0.95);
\node[below=7pt, font=\footnotesize, text=poppurple] at (8.3,-0.95) {subordonnée interrogative indirecte};
\end{tikzpicture}
\legende{Structure de l'interrogative indirecte : \all{Er fragt, woher die Information kommt.}}
\end{popfigure}

\courssub{Cas 1 : la question fermée (oui/non) — ob}

Quand la question directe commence par le verbe (\all{Kommt er? Stimmt die Nachricht?}), c'est-à-dire qu'elle attend « oui » ou « non », la question indirecte est introduite par \all{ob} (= « si » interrogatif).

\begin{center}
\begin{tabularx}{\linewidth}{|X|X|}
\hline
\rowcolor{popblueL} \textbf{Question directe} & \textbf{Question indirecte} \\
\hline
\all{Stimmt die Nachricht?} & \all{Er fragt, ob die Nachricht stimmt.} \\
\hline
\all{Ist die Quelle zuverlässig?} & \all{Ich weiß nicht, ob die Quelle zuverlässig ist.} \\
\hline
\all{Hat der Journalist das Interview geführt?} & \all{Sie fragt, ob der Journalist das Interview geführt hat.} \\
\hline
\all{Wird die Zeitung morgen erscheinen?} & \all{Wir fragen uns, ob die Zeitung morgen erscheinen wird.} \\
\hline
\end{tabularx}
\end{center}

\begin{exemplebox}[Exemples traduits — ob]
\all{Er fragt, ob die Nachricht stimmt.} = « Il demande si la nouvelle est exacte. »\\
\all{Ich weiß nicht, ob er kommt.} = « Je ne sais pas s'il vient. »\\
\all{Die Praktikantin weiß nicht, ob sie den Fehler korrigieren soll.} = « La stagiaire ne sait pas si elle doit corriger l'erreur. »\\
\all{Kannst du mir sagen, ob das Internet hier funktioniert?} = « Peux-tu me dire si Internet fonctionne ici ? »
\end{exemplebox}

\courssub{Cas 2 : la question ouverte — les w-Fragen}

Quand la question directe commence par un mot interrogatif, on reprend \textbf{le même mot} dans la question indirecte, puis le sujet, puis le verbe à la fin.

\begin{center}
\begin{tabularx}{\linewidth}{|l|l|X|}
\hline
\rowcolor{popblueL} \textbf{Mot interrogatif} & \textbf{Français} & \textbf{Exemple en question indirecte} \\
\hline
\all{wer} & qui (sujet) & \all{Wir wissen nicht, wer den Artikel geschrieben hat.} \\
\hline
\all{was} & quoi, que & \all{Er fragt, was die Zeitung berichtet.} \\
\hline
\all{wann} & quand & \all{Sie fragt, wann die Zeitung erscheint.} \\
\hline
\all{wo} & où & \all{Ich weiß nicht, wo der Journalist arbeitet.} \\
\hline
\all{woher} & d'où & \all{Er fragt, woher die Information kommt.} \\
\hline
\all{wohin} & vers où & \all{Sie fragt, wohin der Reporter reist.} \\
\hline
\all{warum} & pourquoi & \all{Die Leser fragen sich, warum die Nachrichten negativ sind.} \\
\hline
\all{wie} & comment & \all{Sie überlegt, wie sie es erklären kann.} \\
\hline
\all{wie viele} & combien de & \all{Er will wissen, wie viele Leser die Zeitung hat.} \\
\hline
\all{wem} (Dat.) & à qui & \all{Sie wissen nicht, wem sie glauben können.} \\
\hline
\end{tabularx}
\end{center}

\begin{attention}[Attention aux formes déclinées de l'interrogatif]
Le mot interrogatif garde sa fonction grammaticale dans la subordonnée : \all{wer} (nominatif, « qui » sujet), \all{wen} (accusatif, « qui » complément d'objet direct), \all{wem} (datif, « à qui » complément d'objet indirect).\\
\all{Wir wissen nicht, \textbf{wen} der Journalist interviewt hat.} = « Nous ne savons pas qui le journaliste a interviewé. »\\
\all{Sie wissen nicht, \textbf{wem} sie glauben können.} = « Ils ne savent pas à qui ils peuvent croire. »
\end{attention}

\courssub{Les verbes introducteurs}

L'interrogative indirecte dépend d'un verbe de la proposition principale. Voici les plus fréquents :

\begin{itemize}
\item \all{sich fragen} — se demander : \all{Ich frage mich, ob das stimmt.} (« Je me demande si c'est vrai. »)
\item \all{fragen} — demander : \all{Er fragt den Reporter, wann er das Interview geführt hat.}
\item \all{wissen} — savoir : \all{Wir wissen nicht, wer den Artikel geschrieben hat.}
\item \all{wissen wollen} — vouloir savoir : \all{Der Chefredakteur will wissen, woher die Information kommt.}
\item \all{sagen} — dire : \all{Kannst du mir sagen, ob die Zeitung heute erscheint?}
\item \all{erklären} — expliquer : \all{Der Lehrer erklärt uns, wie man eine Quelle prüft.} (« Le professeur nous explique comment on vérifie une source. »)
\item \all{überlegen} — réfléchir, se demander : \all{Sie überlegt, wie sie es dem Redakteur erklären kann.}
\end{itemize}

\begin{aretenir}[Comparaison français / allemand]
Le français dit : « Je me demande \textbf{s'il} \textbf{vient}. » — l'ordre des mots ne change pas.\\
L'allemand dit : \all{Ich frage mich, \textbf{ob} er \textbf{kommt}.} — le verbe part à la fin.\\
C'est la même logique que pour \all{weil}, \all{dass}, \all{wenn} : toute subordonnée allemande renvoie le verbe conjugué en position finale. Le « si » interrogatif français se traduit toujours par \all{ob}, jamais par \all{wenn}.
\end{aretenir}

\courssub{L'interrogative indirecte à tous les temps}

Le verbe de la subordonnée se conjugue au temps exigé par le sens. Rappel des trois temps principaux :

\begin{center}
\begin{tabular}{|l|l|l|}
\hline
\rowcolor{popblueL} \textbf{Temps} & \textbf{Exemple} & \textbf{Traduction} \\
\hline
Präsens & \all{Er fragt, ob die Nachricht stimmt.} & Il demande si la nouvelle est vraie. \\
\hline
Perfekt & \all{Sie fragt, wann er das Interview geführt hat.} & Elle demande quand il a mené l'interview. \\
\hline
Präteritum & \all{Wir wussten nicht, wer den Artikel schrieb.} & Nous ne savions pas qui écrivait l'article. \\
\hline
Futur I & \all{Ich weiß nicht, ob die Zeitung erscheinen wird.} & Je ne sais pas si le journal paraîtra. \\
\hline
\end{tabular}
\end{center}

Notez bien : au \textbf{Perfekt} et au \textbf{Futur I}, c'est l'auxiliaire (\all{hat}, \all{wird}) qui occupe la dernière place, et le participe passé (\all{geführt}) ou l'infinitif (\all{erscheinen}) se place juste avant lui. Au Präteritum, la proposition principale est souvent au passé aussi : \all{Wir wussten nicht, ...}.

\begin{methode}[Transformer une question directe en question indirecte]
Pour rapporter une question, suivez toujours les trois étapes :
\begin{enumerate}
\item \textbf{Identifier} le mot interrogatif de la question directe ; s'il n'y en a pas (question oui/non), choisir \all{ob}.
\item \textbf{Placer le sujet} juste après ce mot introducteur.
\item \textbf{Envoyer le verbe conjugué à la fin} de la subordonnée (et l'auxiliaire tout à la fin si le verbe est composé).
\end{enumerate}
Exemple : \all{Woher kommt die Information?} $\rightarrow$ 1) mot interrogatif : \all{woher} ; 2) sujet : \all{die Information} ; 3) verbe à la fin : \all{kommt}. Résultat : \all{Er fragt, woher die Information kommt.} = « Il demande d'où vient l'information. »
\end{methode}

\begin{exoresolu}[De la question directe à la question indirecte]
\textbf{Énoncé.} Transformez les questions directes suivantes en questions indirectes, en commençant par l'expression donnée.
\begin{enumerate}
\item \all{Erscheint die Zeitung heute?} — \all{Ich weiß nicht, ...}
\item \all{Wer hat diesen Artikel geschrieben?} — \all{Wir wissen nicht, ...}
\item \all{Woher kommt diese Information?} — \all{Der Journalist fragt, ...}
\item \all{Warum sind so viele Nachrichten negativ?} — \all{Die Leser fragen sich, ...}
\end{enumerate}
\tcblower
\textbf{Solution détaillée.}
\begin{enumerate}
\item Question fermée (pas de mot interrogatif) $\rightarrow$ \all{ob} ; sujet \all{die Zeitung} ; verbe \all{erscheint} à la fin : \all{Ich weiß nicht, ob die Zeitung heute erscheint.} (« Je ne sais pas si le journal paraît aujourd'hui. »)
\item Mot interrogatif \all{wer} (sujet, nominatif) ; verbe composé au Perfekt : l'auxiliaire \all{hat} à la toute fin, le participe \all{geschrieben} juste avant : \all{Wir wissen nicht, wer diesen Artikel geschrieben hat.} (« Nous ne savons pas qui a écrit cet article. »)
\item Mot interrogatif \all{woher} ; sujet \all{diese Information} ; verbe \all{kommt} à la fin : \all{Der Journalist fragt, woher diese Information kommt.} (« Le journaliste demande d'où vient cette information. »)
\item Mot interrogatif \all{warum} ; sujet \all{so viele Nachrichten} ; verbe \all{sind} à la fin : \all{Die Leser fragen sich, warum so viele Nachrichten negativ sind.} (« Les lecteurs se demandent pourquoi tant de nouvelles sont négatives. »)
\end{enumerate}
\end{exoresolu}

\begin{attention}[Les deux pièges mortels du francophone]
\textbf{Piège n° 1 — le calque de l'ordre français.} Ne dites jamais : \emph{« Ich weiß nicht, ob kommt er »}. C'est un calque mot à mot de « je ne sais pas si vient-il / s'il vient ». En allemand, le verbe conjugué va \textbf{toujours à la fin} : \all{Ich weiß nicht, ob er kommt.}

\textbf{Piège n° 2 — confondre \all{ob} et \all{wenn}.} Le français « si » a deux valeurs :
\begin{itemize}
\item « si » \textbf{interrogatif} (on peut le remplacer par « est-ce que ») $\rightarrow$ \all{ob} : \all{Ich weiß nicht, ob er kommt.} (« Je ne sais pas s'il vient. »)
\item « si » \textbf{conditionnel} (hypothèse) $\rightarrow$ \all{wenn} : \all{Wenn er kommt, freue ich mich.} (« S'il vient, je suis content. »)
\end{itemize}
Test pratique : si vous pouvez remplacer « si » par « est-ce que », c'est \all{ob}. Sinon, c'est \all{wenn}.
\end{attention}

\begin{exercice}[À vous de jouer]
\begin{enumerate}
\item Traduisez en allemand : a) « Elle demande quand le journal paraît. » b) « Je me demande si cette source est fiable. » c) « Nous ne savons pas qui a écrit l'article. »
\item Transformez : \all{Hat der Reporter das Interview geführt?} — \all{Der Chefredakteur will wissen, ...}
\item Choisissez \all{ob} ou \all{wenn} : \all{Ich frage mich, ... das Internet in Douala schnell ist.} / \all{... die Zeitung erscheint, kaufen wir sie.}
\end{enumerate}
\end{exercice}

\begin{corrige}[Exercice « À vous de jouer »]
1. a) \all{Sie fragt, wann die Zeitung erscheint.} b) \all{Ich frage mich, ob diese Quelle zuverlässig ist.} c) \all{Wir wissen nicht, wer den Artikel geschrieben hat.}\\
2. \all{Der Chefredakteur will wissen, ob der Reporter das Interview geführt hat.} (question fermée $\rightarrow$ \all{ob} ; Perfekt : participe \all{geführt} puis auxiliaire \all{hat} à la fin.)\\
3. \all{Ich frage mich, \textbf{ob} das Internet in Douala schnell ist.} (« est-ce que » possible $\rightarrow$ \all{ob}) ; \all{\textbf{Wenn} die Zeitung erscheint, kaufen wir sie.} (condition : « quand/si le journal paraît » $\rightarrow$ \all{wenn} ; attention, la principale suit avec le verbe en deuxième position : \all{kaufen wir}.)
\end{corrige}

\begin{aretenir}[L'essentiel de la section]
\begin{itemize}
\item Question fermée (oui/non) $\rightarrow$ \all{ob} ; question ouverte $\rightarrow$ mot interrogatif en \all{w-} (\all{wer, was, wann, wo, woher, warum, wie...}).
\item Dans la subordonnée interrogative indirecte, le \textbf{verbe conjugué est en dernière position} ; auxiliaire tout à la fin si le verbe est composé.
\item Verbes introducteurs typiques : \all{sich fragen, fragen, wissen, wissen wollen, sagen, erklären, überlegen}.
\item « Si » interrogatif = \all{ob} ; « si » conditionnel = \all{wenn}. Jamais de calque de l'ordre français !
\end{itemize}
\end{aretenir}

\coursec{Grammaire 2 : les verbes à complément prépositionnel}

Dans la section précédente, nous avons appris à construire les interrogatives indirectes avec \all{ob} et les mots interrogatifs en \all{w-} : \all{Ich weiß nicht, ob die Quelle zuverlässig ist.} « Je ne sais pas si la source est fiable. » Nous avons remarqué que certaines questions commençaient par des formes étranges comme \all{worauf} ou \all{worüber} : \all{Er fragt, worauf sie wartet.} « Il demande ce qu'elle attend. » Ces formes ne sont pas un hasard : elles naissent d'une construction fondamentale de l'allemand, le \cle{verbe à complément prépositionnel} (\all{das Verb mit Präpositionalobjekt}). C'est l'objet de cette section, et c'est aussi l'une des difficultés majeures pour les francophones, car le français et l'allemand ne répartissent pas les prépositions de la même façon.

\courssub{Observation : partir du texte}

Reprenons trois phrases inspirées de notre texte support \all{Informationen im digitalen Zeitalter} :

\begin{exemplebox}[Observer avant de comprendre]
\all{Der Journalist berichtet über die Wahrheit.} « Le journaliste rend compte de la vérité. »\par
\all{Er wartet auf eine Bestätigung.} « Il attend une confirmation. »\par
\all{Die Leser informieren sich über das Internet.} « Les lecteurs s'informent sur Internet. »
\end{exemplebox}

Observons attentivement :

\begin{itemize}
\item En allemand, on ne peut pas dire \all{der Journalist berichtet die Wahrheit} ni \all{er wartet eine Bestätigung} : les verbes \all{berichten} et \all{warten} exigent une \cle{préposition} (\all{über}, \all{auf}) devant leur complément.
\item En français, c'est parfois l'inverse : « attendre » se construit \emph{sans} préposition (j'attends la nouvelle), alors que « s'informer » en exige une (s'informer \emph{sur} quelque chose).
\item La préposition allemande impose en plus un \cle{cas} : \all{auf eine Bestätigung} (accusatif), \all{über die Wahrheit} (accusatif).
\end{itemize}

Conclusion de l'observation : en allemand, beaucoup de verbes sont « mariés » à une préposition fixe, et ce couple verbe + préposition + cas doit s'apprendre par cœur, comme on apprend l'article avec le nom.

\begin{propriete}[La règle fondamentale]
Certains verbes allemands se construisent \textbf{obligatoirement} avec une \cle{préposition fixe} suivie d'un \cle{cas fixe} (accusatif ou datif). Cette construction est invariable : elle fait partie de l'identité du verbe.\par
\all{warten auf + Akkusativ} = attendre (quelque chose/quelqu'un)\par
\all{sich informieren über + Akkusativ} = s'informer sur\par
La préposition et le cas s'apprennent \textbf{en même temps que le verbe}, exactement comme le genre d'un nom : on n'apprend pas \all{warten}, on apprend \all{warten auf + Akk.}.
\end{propriete}

\courssub{La liste du programme : les verbes à connaître}

Voici les verbes à complément prépositionnel exigés en Terminale A, avec leur cas et leur traduction.

\begin{center}
\begin{tabularx}{\linewidth}{|l|l|l|X|}
\hline
\rowcolor{popblueL} \textbf{Verbe} & \textbf{Präposition} & \textbf{Kasus} & \textbf{Traduction} \\
\hline
\all{warten} & \all{auf} & Akk. & attendre \\
\hline
\all{sich informieren} & \all{über} & Akk. & s'informer sur \\
\hline
\all{berichten} & \all{über} & Akk. & rapporter, rendre compte de \\
\hline
\all{sich freuen} & \all{über} & Akk. & se réjouir de (fait présent ou passé) \\
\hline
\all{sich freuen} & \all{auf} & Akk. & se réjouir à l'idée de (fait futur) \\
\hline
\all{denken} & \all{an} & Akk. & penser à \\
\hline
\all{sprechen} & \all{über} & Akk. & parler de \\
\hline
\all{bitten} & \all{um} & Akk. & demander, prier de \\
\hline
\all{sich interessieren} & \all{für} & Akk. & s'intéresser à \\
\hline
\all{Angst haben} & \all{vor} & Dat. & avoir peur de \\
\hline
\all{helfen} & \all{bei} & Dat. & aider à (faire quelque chose) \\
\hline
\end{tabularx}
\end{center}

Remarque : la grande majorité de ces verbes exigent l'\cle{accusatif} ; seuls \all{Angst haben vor} et \all{helfen bei} exigent le \cle{datif}. Quand tu hésites sur le cas, l'accusatif est statistiquement le plus probable — mais la seule vraie sécurité est d'avoir mémorisé chaque triplet verbe + préposition + cas.

\begin{exemplebox}[Exemples traduits]
\all{Ich warte auf die Nachricht.} « J'attends la nouvelle. »\par
\all{Wir sprechen über das Internet.} « Nous parlons d'Internet. »\par
\all{Er interessiert sich für die Zeitung.} « Il s'intéresse au journal. »\par
\all{Die Schüler denken an die Prüfung.} « Les élèves pensent à l'examen. »\par
\all{Sie bittet den Journalisten um eine Antwort.} « Elle demande une réponse au journaliste. »\par
\all{Das Kind hat Angst vor dem Hund.} « L'enfant a peur du chien. »\par
\all{Kannst du mir bei der Übersetzung helfen?} « Peux-tu m'aider dans la traduction ? »
\end{exemplebox}

Observe la phrase \all{Sie bittet den Journalisten um eine Antwort} : \all{bitten} peut avoir \textbf{deux} compléments — la personne interrogée à l'accusatif (\all{den Journalisten}) et la chose demandée introduite par \all{um} (\all{eine Antwort}).

\begin{popfigure}
\begin{center}
\begin{tikzpicture}[
  node distance=0.6cm,
  verb/.style={rectangle, rounded corners, draw=popblue, fill=popblueL, align=center, font=\small, minimum height=0.8cm},
  prep/.style={rectangle, draw=poporange, fill=poporangeL, align=center, font=\small},
  kasus/.style={rectangle, draw=popgreen, fill=popgreenL, align=center, font=\small},
  trad/.style={rectangle, draw=poppurple, fill=poppinkL, align=center, font=\small}
]
\node[verb] (v) {\all{warten}};
\node[prep, right=of v] (p) {\all{auf}};
\node[kasus, right=of p] (k) {+ Akkusativ};
\node[trad, right=of k] (t) {« attendre »};
\draw[-{Stealth}, thick, popdark] (v) -- (p);
\draw[-{Stealth}, thick, popdark] (p) -- (k);
\draw[-{Stealth}, thick, popdark] (k) -- (t);
\node[below=0.9cm of p, align=center, font=\small] (q1) {\textbf{Chose :} \all{Worauf wartet er?}\\ \all{— Auf eine Bestätigung.}};
\node[below=0.9cm of k, align=center, font=\small] (q2) {\textbf{Personne :} \all{Auf wen wartet er?}\\ \all{— Auf den Journalisten.}};
\draw[-{Stealth}, thick, poppurple] (k) -- (q1);
\draw[-{Stealth}, thick, poppurple] (k) -- (q2);
\end{tikzpicture}
\end{center}
\legende{Le verbe \all{warten auf} + Akkusativ et les deux façons d'interroger son complément : \all{worauf} pour une chose, \all{auf wen} pour une personne.}
\end{popfigure}

\courssub{Conjugaison : \all{sich informieren über} et \all{warten auf}}

\textbf{Präsens de \all{sich informieren über} (verbe faible réfléchi)} — attention, le pronom réfléchi est ici à l'accusatif :

\begin{center}
\begin{tabular}{|l|l|l|}
\hline
\rowcolor{popblueL} \textbf{Personne} & \textbf{Forme} & \textbf{Traduction} \\
\hline
ich & \all{informiere mich über} & je m'informe sur \\
\hline
du & \all{informierst dich über} & tu t'informes sur \\
\hline
er/sie/es & \all{informiert sich über} & il/elle s'informe sur \\
\hline
wir & \all{informieren uns über} & nous nous informons sur \\
\hline
ihr & \all{informiert euch über} & vous vous informez sur \\
\hline
sie/Sie & \all{informieren sich über} & ils/vous s'informent sur \\
\hline
\end{tabular}
\end{center}

\textbf{Perfekt de \all{sich informieren über}} : auxiliaire \all{haben} + participe \all{informiert} (verbe en \all{-ieren} : pas de \all{ge-} au participe !).

\all{Ich habe mich über die Wahl informiert.} « Je me suis informé sur l'élection. »\par
\all{Sie hat sich über die Quelle informiert.} « Elle s'est informée sur la source. »

\textbf{Präsens et Perfekt de \all{warten auf} (verbe faible)} :

\begin{center}
\begin{tabular}{|l|l|l|}
\hline
\rowcolor{popblueL} \textbf{Personne} & \textbf{Präsens} & \textbf{Perfekt} \\
\hline
ich & \all{warte auf} & \all{habe auf ... gewartet} \\
\hline
du & \all{wartest auf} & \all{hast auf ... gewartet} \\
\hline
er/sie/es & \all{wartet auf} & \all{hat auf ... gewartet} \\
\hline
wir & \all{warten auf} & \all{haben auf ... gewartet} \\
\hline
ihr & \all{wartet auf} & \all{habt auf ... gewartet} \\
\hline
sie/Sie & \all{warten auf} & \all{haben auf ... gewartet} \\
\hline
\end{tabular}
\end{center}

\all{Der Journalist hat lange auf eine Bestätigung gewartet.} « Le journaliste a longtemps attendu une confirmation. »

\begin{attention}[Ne jamais oublier le réfléchi]
\all{sich informieren}, \all{sich freuen}, \all{sich interessieren} sont des verbes \textbf{réfléchis} : le pronom \all{mich, dich, sich, uns, euch} est obligatoire. Dire \all{ich informiere über die Wahl} est une faute : il faut \all{ich informiere \textbf{mich} über die Wahl}. Au Perfekt, le réfléchi reste : \all{ich habe \textbf{mich} informiert}.
\end{attention}

\courssub{Interroger le complément : \all{worauf}, \all{worüber} ou \all{auf wen} ?}

Comment poser une question sur le complément prépositionnel ? Tout dépend de la \cle{nature} du complément : chose ou personne.

\begin{methode}[Poser une question sur le complément prépositionnel]
\begin{enumerate}
\item Identifier le couple verbe + préposition : \all{warten auf}, \all{sich informieren über}.
\item Demander : le complément est-il une \textbf{chose} ou une \textbf{personne} ?
\item Pour une \textbf{chose} : former un mot interrogatif composé \cle{wo(r) + préposition} : \all{worauf}, \all{worüber}, \all{wofür}, \all{wovor}, \all{wobei}. On ajoute un \all{-r-} quand la préposition commence par une voyelle (\all{wo + auf = worauf}, \all{wo + über = worüber}).
\item Pour une \textbf{personne} : garder la préposition et ajouter \all{wen} (accusatif) ou \all{wem} (datif) : \all{auf wen ?}, \all{an wen ?}, \all{vor wem ?}.
\end{enumerate}
\end{methode}

\begin{exemplebox}[Questions et réponses]
\all{Worauf wartet er? — Auf eine Bestätigung.} « Qu'attend-il ? — Une confirmation. » (chose)\par
\all{Auf wen wartet er? — Auf den Journalisten.} « Qui attend-il ? — Le journaliste. » (personne)\par
\all{Worüber informiert sie sich? — Über die Wahl.} « Sur quoi s'informe-t-elle ? — Sur l'élection. »\par
\all{Wofür interessiert er sich? — Für die Zeitung.} « À quoi s'intéresse-t-il ? — Au journal. »\par
\all{Wovor hat das Kind Angst? — Vor dem Hund.} « De quoi l'enfant a-t-il peur ? — Du chien. »
\end{exemplebox}

\begin{attention}[\all{worüber}, pas \all{über was}]
Les élèves francophones traduisent souvent « de quoi » par \all{über was} ou \all{was über}. C'est incorrect à l'écrit et fautif à l'examen. Pour les choses, la forme correcte est le composé en \all{wo(r)-} : \all{Worüber sprichst du?} « De quoi parles-tu ? ». La forme \all{über was} existe dans la langue parlée familière, mais elle est à bannir dans une copie.
\end{attention}

\courssub{Le lien avec les interrogatives indirectes}

Nous bouclons maintenant avec la Grammaire 1 : quand la question sur le complément prépositionnel devient \cle{interrogative indirecte}, le mot composé \all{wo(r)-} joue le rôle de subordonnant, et le verbe conjugué passe en \cle{dernière position}.

\begin{exemplebox}[De la question directe à l'interrogative indirecte]
Question directe : \all{Worauf wartet sie?} « Qu'attend-elle ? »\par
Interrogative indirecte : \all{Er fragt, worauf sie wartet.} « Il demande ce qu'elle attend. »\par
Question directe : \all{Worüber berichtet der Journalist?} « Sur quoi le journaliste rend-il compte ? »\par
Interrogative indirecte : \all{Ich weiß nicht, worüber der Journalist berichtet.} « Je ne sais pas sur quoi le journaliste rend compte. »
\end{exemplebox}

Observe bien : dans \all{er fragt, worauf sie wartet}, la préposition est « collée » au \all{wo-} et le verbe \all{wartet} va à la fin. La préposition ne se sépare jamais du \all{wo(r)-} dans ce type de phrase.

\courssub{Comparaison français / allemand : le grand décalage}

C'est le cœur du problème pour un francophone : les deux langues ne distribuent pas les prépositions de la même manière.

\begin{center}
\begin{tabularx}{\linewidth}{|X|X|X|}
\hline
\rowcolor{popblueL} \textbf{Français} & \textbf{Deutsch} & \textbf{Remarque} \\
\hline
attendre qch (sans prép.) & \all{warten auf etwas} & préposition obligatoire en allemand \\
\hline
s'informer sur qch & \all{sich informieren über etwas} & parallèle, mais réfléchi obligatoire \\
\hline
parler de qch & \all{über etwas sprechen} & parallèle \\
\hline
penser à qch & \all{an etwas denken} & la préposition change : « à » devient \all{an} \\
\hline
s'intéresser à qch & \all{sich für etwas interessieren} & « à » devient \all{für} ! \\
\hline
avoir peur de qch & \all{Angst vor etwas haben} & « de » devient \all{vor} + Dativ \\
\hline
demander qch & \all{um etwas bitten} & préposition obligatoire en allemand \\
\hline
\end{tabularx}
\end{center}

La leçon à tirer de ce tableau : \textbf{ne jamais traduire la préposition française mot à mot}. Le français « s'intéresser \emph{à} » ne donne pas \all{sich interessieren an} mais \all{sich interessieren für}. Seule la mémorisation du triplet verbe + préposition + cas garantit l'exactitude.

\begin{attention}[L'erreur n° 1 des francophones : \all{ich warte die Nachricht}]
Le français « attendre » est transitif direct (j'attends \emph{la nouvelle}), donc l'élève francophone écrit spontanément \all{ich warte die Nachricht}. C'est \textbf{faux} : l'allemand \all{warten} est intransitif et exige \all{auf} + Akkusativ. La seule forme correcte est : \all{Ich warte \textbf{auf} die Nachricht.} Même vigilance pour \all{bitten} : on ne dit pas \all{sie bittet eine Antwort}, mais \all{sie bittet \textbf{um} eine Antwort}.
\end{attention}

\begin{attention}[\all{sich freuen über} ou \all{sich freuen auf} ? — la distinction classique du Bac]
Les deux existent, avec un sens différent :\par
\all{sich freuen über + Akk.} = se réjouir \textbf{de} quelque chose qui existe déjà (présent ou passé) : \all{Ich freue mich über die gute Nachricht.} « Je me réjouis de la bonne nouvelle » (elle est arrivée).\par
\all{sich freuen auf + Akk.} = se réjouir \textbf{à l'idée de} quelque chose de futur : \all{Ich freue mich auf die Ferien.} « Je me réjouis à l'idée des vacances » (elles ne sont pas encore là).\par
Astuce : \all{auf} regarde vers l'avant (avenir), \all{über} regarde ce qui est déjà là.
\end{attention}

\courssub{Texte d'application : les verbes prépositionnels en contexte}

\begin{textebox}[Ein Interview in Yaoundé]
Gestern hat eine Journalistin der Zeitung „Die Brücke“ einen Schüler aus Yaoundé interviewt. Zuerst hat sie auf den Direktor der Schule gewartet, denn sie wollte ihn um eine Erlaubnis bitten. Der Direktor hat lange auf eine Bestätigung aus dem Ministerium gewartet, aber am Nachmittag hat er endlich zugestimmt.\par
Im Gespräch hat sich der Schüler für die Arbeit der Journalisten interessiert. „Ich informiere mich jeden Tag über die Nachrichten im Internet“, hat er gesagt, „aber ich denke oft an die Qualität der Quellen. Viele junge Leute sprechen über das Internet, aber sie prüfen die Informationen nicht.“ Die Journalistin hat über diese Antwort berichtet und sich über die klugen Ideen des Schülers gefreut. Sie hat sich auch für das Schulprojekt interessiert und um weitere Informationen gebeten.\par
Am Ende hat der Schüler gefragt: „Wovor haben die Journalisten Angst?“ Die Journalistin hat geantwortet: „Wir haben Angst vor falschen Nachrichten. Deshalb warten wir immer auf gute Quellen und helfen unseren Lesern bei der Suche nach der Wahrheit.“ Der Schüler freut sich jetzt auf den Artikel, der nächste Woche in der Zeitung erscheint.
\end{textebox}

\textbf{Glossaire} : \all{die Erlaubnis, -se} (l'autorisation) ; \all{die Bestätigung, -en} (la confirmation) ; \all{zustimmen} (accepter, donner son accord) ; \all{die Quelle, -n} (la source) ; \all{prüfen} (vérifier) ; \all{falsch} (faux) ; \all{die Suche, -n} (la recherche) ; \all{erscheinen} (paraître) ; \all{der Leser, -} (le lecteur).

\textbf{Questions de compréhension} :
\begin{enumerate}
\item \all{Auf wen wartet die Journalistin zuerst?}
\item \all{Wofür interessiert sich der Schüler?}
\item \all{Wovor haben die Journalisten Angst?}
\item \all{Worauf freut sich der Schüler am Ende?}
\end{enumerate}

\begin{corrige}[Questions de compréhension]
1. \all{Sie wartet auf den Direktor der Schule.} (personne : \all{auf wen})\par
2. \all{Er interessiert sich für die Arbeit der Journalisten.}\par
3. \all{Sie haben Angst vor falschen Nachrichten.} (\all{vor} + Dativ : \all{falschen Nachrichten})\par
4. \all{Er freut sich auf den Artikel.} (\all{auf}, car l'article n'est pas encore paru : futur !)
\end{corrige}

\courssub{Exercice résolu et entraînement}

\begin{exoresolu}[Compléter avec la préposition et le cas corrects]
Énoncé : complète avec la préposition qui convient et mets le nom au bon cas.\par
a) \all{Der Journalist berichtet \dots\ (die Wahl).}\par
b) \all{Ich warte \dots\ (die Nachricht).}\par
c) \all{Sie interessiert sich \dots\ (das Internet).}\par
d) \all{Wir haben Angst \dots\ (die Prüfung).}\par
e) \all{Er bittet den Lehrer \dots\ (eine Erklärung).}
\tcblower
Solution détaillée :\par
a) \all{Der Journalist berichtet \textbf{über die Wahl}.} — \all{berichten über} + Akk. ; \all{die Wahl} est féminin, l'article ne change pas à l'accusatif.\par
b) \all{Ich warte \textbf{auf die Nachricht}.} — \all{warten auf} + Akk. ; jamais sans préposition !\par
c) \all{Sie interessiert sich \textbf{für das Internet}.} — \all{sich interessieren für} + Akk. ; \all{das} reste \all{das} à l'accusatif ; ne pas oublier le réfléchi \all{sich}.\par
d) \all{Wir haben Angst \textbf{vor der Prüfung}.} — \all{Angst haben vor} + \textbf{Dativ} : \all{die Prüfung} devient \all{der Prüfung}. C'est l'un des rares verbes de la liste au datif.\par
e) \all{Er bittet den Lehrer \textbf{um eine Erklärung}.} — \all{bitten um} + Akk.
\end{exoresolu}

\begin{exercice}[À toi de jouer]
1. Traduis en allemand : a) « J'attends le journal. » b) « Nous parlons de la source. » c) « Elle se réjouit à l'idée des vacances. » d) « De quoi parlez-vous ? » (chose)\par
2. Mets la question directe au discours indirect : \all{Worauf wartet der Journalist?} → \all{Ich weiß nicht, \dots}\par
3. Trouve l'erreur et corrige : \all{Ich warte den Bus.} / \all{Sie informiert über die Nachrichten.}
\end{exercice}

\begin{corrige}[Exercice « À toi de jouer »]
1. a) \all{Ich warte auf die Zeitung.} b) \all{Wir sprechen über die Quelle.} c) \all{Sie freut sich auf die Ferien.} (futur → \all{auf}) d) \all{Worüber sprechen Sie?}\par
2. \all{Ich weiß nicht, worauf der Journalist wartet.} (verbe en dernière position)\par
3. \all{Ich warte \textbf{auf} den Bus.} (préposition obligatoire) ; \all{Sie informiert \textbf{sich} über die Nachrichten.} (réfléchi obligatoire).
\end{corrige}

\begin{aretenir}[L'essentiel de la section]
\begin{itemize}
\item Un verbe à complément prépositionnel s'apprend en bloc : \textbf{verbe + préposition + cas} (\all{warten auf + Akk.}).
\item La plupart exigent l'accusatif ; \all{Angst haben vor} et \all{helfen bei} exigent le datif.
\item Pour interroger le complément : \all{wo(r)-} + préposition pour les choses (\all{worauf, worüber}) ; préposition + \all{wen/wem} pour les personnes (\all{auf wen}).
\item En interrogative indirecte, le composé \all{wo(r)-} subordonne et le verbe va à la fin : \all{Er fragt, worauf sie wartet.}
\item \all{sich freuen über} = se réjouir de (présent/passé) ; \all{sich freuen auf} = se réjouir à l'idée de (futur).
\item Jamais \all{ich warte die Nachricht} : toujours \all{ich warte \textbf{auf} die Nachricht}.
\end{itemize}
\end{aretenir}

\coursec{Commentaire modèle et production guidée}

Dans les sections précédentes, nous avons appris à construire des phrases avec des \cle{verbes à complément prépositionnel} comme \all{sich informieren über} (+ Akkusativ) et des \cle{interrogatives indirectes} avec \all{ob} ou un mot interrogatif en \all{w-}. Ces deux outils grammaticaux ne sont pas de simples exercices : ce sont les instruments qui vont vous permettre de \emph{rédiger} et de \emph{parler} sur le texte. C'est exactement ce que l'on attend de vous au baccalauréat : lire, comprendre, puis \cle{commenter} et \cle{produire}. Cette dernière section vous donne la méthode complète, un modèle rédigé, et des productions guidées écrites et orales.

\courssub{La méthode du commentaire de texte en quatre étapes}

\begin{methode}[Rédiger un commentaire de texte (Textkommentar)]
Un commentaire de texte en allemand suit toujours la même architecture en quatre étapes. Respectez cet ordre, il est évalué au Bac.
\begin{enumerate}
\item \textbf{Einleitung (introduction)} : présentez le texte en une ou deux phrases — auteur, source, thème. Formule type : \all{Der Text ... von ... handelt von ...} (« Le texte ... de ... traite de ... »). Variantes : \all{Der Artikel erscheint in der Zeitung...} (« L'article paraît dans le journal... ») / \all{Es geht um...} (« Il s'agit de... »).
\item \textbf{Zusammenfassung (résumé)} : résumez le contenu essentiel au \cle{Präsens}, avec vos propres mots, sans recopier de longues phrases du texte. Utilisez \all{zuerst} (d'abord), \all{dann} (ensuite), \all{außerdem} (de plus).
\item \textbf{Analyse} : observez la langue et l'intention — le lexique du thème (ici : les médias), les structures grammaticales remarquables (interrogatives indirectes, verbes à préposition), et le but de l'auteur (\all{Der Autor will zeigen, dass...} « L'auteur veut montrer que... »).
\item \textbf{Eigene Meinung (avis personnel)} : terminez par votre opinion argumentée avec \all{meiner Meinung nach} (à mon avis), \all{ich denke, dass...} (je pense que...), \all{ich finde es wichtig, dass...} (je trouve important que...). Connecteur de clôture : \all{schließlich} (finalement).
\end{enumerate}
\end{methode}

\begin{popfigure}
\begin{center}
\begin{tikzpicture}[
  node distance=0.55cm,
  etape/.style={rectangle, rounded corners=4pt, draw=popblue, thick, fill=popblueL, text width=4.6cm, align=center, minimum height=1.1cm, font=\small},
  fleche/.style={-{Stealth[length=3mm]}, thick, popdark}
]
\node[etape, align=center] (e1){\textbf{1. Einleitung}\\ \all{Der Text ... handelt von ...}};
\node[etape, below=of e1, fill=popgreenL, draw=popgreen, align=center] (e2){\textbf{2. Zusammenfassung}\\ Präsens, eigene Worte\\ \all{zuerst -- dann -- außerdem}};
\node[etape, below=of e2, fill=poporangeL, draw=poporange, align=center] (e3){\textbf{3. Analyse}\\ Lexik, Grammatik, Absicht\\ \all{Der Autor will zeigen, dass ...}};
\node[etape, below=of e3, fill=poppinkL, draw=poppink, align=center] (e4){\textbf{4. Eigene Meinung}\\ \all{meiner Meinung nach ...}\\ \all{schließlich ...}};
\draw[fleche] (e1) -- (e2);
\draw[fleche] (e2) -- (e3);
\draw[fleche] (e3) -- (e4);
\end{tikzpicture}
\end{center}
\legende{Organigramme du commentaire de texte : quatre étapes dans l'ordre obligatoire}
\end{popfigure}

\begin{aretenir}[Les connecteurs du commentaire]
\begin{center}
\begin{tabularx}{\linewidth}{|l|X|X|}
\hline
\rowcolor{popblueL} \textbf{Étape} & \textbf{Connecteur allemand} & \textbf{Français} \\
\hline
Introduction & \all{Der Text handelt von ...} & Le texte traite de ... \\
\hline
Résumé & \all{zuerst / dann / außerdem} & d'abord / ensuite / de plus \\
\hline
Analyse & \all{Der Autor erklärt / zeigt / fragt sich, ob ...} & L'auteur explique / montre / se demande si ... \\
\hline
Opinion & \all{meiner Meinung nach / ich denke, dass ...} & à mon avis / je pense que ... \\
\hline
Conclusion & \all{schließlich / zusammenfassend} & finalement / en résumé \\
\hline
\end{tabularx}
\end{center}
Attention à la construction : après \all{meiner Meinung nach} (qui occupe la première place), le verbe conjugué vient immédiatement en deuxième position : \all{Meiner Meinung nach \textbf{sind} Falschmeldungen gefährlich.} « À mon avis, les fausses informations sont dangereuses. »
\end{aretenir}

\courssub{Commentaire modèle rédigé}

\begin{textebox}[Commentaire modèle sur « Informationen im digitalen Zeitalter »]
\all{Der Text handelt von der Information im digitalen Zeitalter. Der Journalist erklärt, dass er seine Quellen überprüft, bevor er einen Artikel schreibt. Er fragt sich, ob die Quellen zuverlässig sind. Außerdem zeigt er, wie sich die Menschen heute im Internet informieren. Meiner Meinung nach ist das sehr wichtig, weil viele Menschen im Internet Falschmeldungen verbreiten. Ich denke, dass man immer mehrere Quellen vergleichen muss. Schließlich informiere ich mich selbst oft im Internet, aber ich achte auf seriöse Seiten.}
\end{textebox}

\textbf{Traduction} : « Le texte traite de l'information à l'ère numérique. Le journaliste explique qu'il vérifie ses sources avant d'écrire un article. Il se demande si les sources sont fiables. De plus, il montre comment les gens s'informent aujourd'hui sur Internet. À mon avis, c'est très important, parce que beaucoup de gens répandent de fausses informations sur Internet. Je pense qu'il faut toujours comparer plusieurs sources. Finalement, je m'informe moi-même souvent sur Internet, mais je fais attention aux sites sérieux. »

Observez comment le modèle applique la méthode : la première phrase est l'\emph{Einleitung} (\all{Der Text handelt von...}), les phrases suivantes forment la \emph{Zusammenfassung} au Präsens avec \all{außerdem}, la phrase \all{Er fragt sich, ob die Quellen zuverlässig sind} exploite une \cle{interrogative indirecte} (analyse grammaticale, verbe \all{sind} en finale), et les trois dernières phrases donnent l'\emph{eigene Meinung} avec \all{meiner Meinung nach}, \all{ich denke, dass} et \all{schließlich}. Notez aussi la précision d'usage : on s'informe \emph{sur un sujet} avec \all{über} + Akkusativ (\all{ich informiere mich über die Politik}), mais \emph{par un média} avec \all{im Internet}, \all{in der Zeitung}, \all{im Radio}.

\begin{exemplebox}[Phrases d'opinion à réutiliser]
\begin{itemize}
\item \all{Meiner Meinung nach sind Falschmeldungen gefährlich.} « À mon avis, les fausses informations sont dangereuses. »
\item \all{Ich informiere mich oft im Internet über die Politik.} « Je m'informe souvent sur Internet au sujet de la politique. »
\item \all{Ich denke, dass die Zeitung zuverlässiger ist als das Internet.} « Je pense que le journal est plus fiable qu'Internet. »
\item \all{Es kommt darauf an, welche Quelle man benutzt.} « Cela dépend de la source que l'on utilise. »
\item \all{Ich zweifle daran, dass alle Nachrichten im Internet wahr sind.} « Je doute que toutes les nouvelles sur Internet soient vraies. »
\end{itemize}
\end{exemplebox}

\courssub{Le résumé guidé (Zusammenfassung)}

Le résumé est un exercice à part entière. Règles strictes : \textbf{trois phrases maximum}, au \cle{Präsens}, avec \emph{vos propres mots}, sans aucune citation longue du texte.

\begin{exemplebox}[Résumé modèle en trois phrases]
\all{Der Text beschreibt die Arbeit eines Journalisten im digitalen Zeitalter. Er muss seine Quellen prüfen, bevor er schreibt. Der Text warnt auch vor Falschmeldungen im Internet.}

« Le texte décrit le travail d'un journaliste à l'ère numérique. Il doit vérifier ses sources avant d'écrire. Le texte met aussi en garde contre les fausses informations sur Internet. »
\end{exemplebox}

\begin{attention}[Les pièges du résumé pour les francophones]
\begin{itemize}
\item \textbf{Ne recopiez pas le texte.} Un résumé qui reprend des phrases entières du texte est sanctionné. Reformulez : le texte dit \all{„Bevor ich einen Artikel schreibe, überprüfe ich jede Quelle“}, vous écrivez \all{Der Journalist prüft seine Quellen vor dem Schreiben.}
\item \textbf{Restez au Präsens}, même si le texte contient du passé : le résumé d'un texte se fait conventionnellement au présent en allemand comme en français.
\item \textbf{Pas d'opinion dans le résumé} : \all{meiner Meinung nach} n'a sa place que dans la quatrième étape du commentaire, jamais dans la Zusammenfassung.
\item \textbf{Place du verbe} : dans \all{bevor er schreibt} et \all{ob die Quellen zuverlässig sind}, le verbe conjugué va à la fin — ne calquez jamais l'ordre français.
\end{itemize}
\end{attention}

\courssub{Production écrite guidée : « Wie informierst du dich? »}

\begin{methode}[Réussir la production écrite (6 à 8 lignes)]
\begin{enumerate}
\item \textbf{Lisez les contraintes} : ici, il faut au moins \textbf{deux interrogatives indirectes} (\all{ob} ou phrase en \all{w-}) et \textbf{deux verbes à complément prépositionnel}.
\item \textbf{Préparez vos outils} : choisissez vos verbes à préposition avant d'écrire — par exemple \all{sich informieren über + Akk.}, \all{sich interessieren für + Akk.}, \all{denken an + Akk.}, \all{sich verlassen auf + Akk.} (se fier à).
\item \textbf{Structurez} : une phrase d'introduction, trois ou quatre phrases de développement (où, quand, pourquoi vous vous informez), une phrase d'opinion finale.
\item \textbf{Vérifiez} : verbe en deuxième position dans les principales, verbe à la fin après \all{ob}, \all{dass}, \all{weil} ; bonne préposition et bon cas après chaque verbe.
\end{enumerate}
\end{methode}

\begin{exoresolu}[Production écrite : « Wie informierst du dich? »]
Énoncé : Rédigez un court texte de 6 à 8 lignes sur vos habitudes d'information. Employez au moins deux interrogatives indirectes et deux verbes à complément prépositionnel.
\tcblower
Solution détaillée — modèle de production :

\all{Ich informiere mich jeden Tag im Internet und höre manchmal Radio. Zuerst frage ich mich immer, ob eine Nachricht wirklich wahr ist. Dann vergleiche ich mehrere Quellen, denn ich verlasse mich nicht auf eine einzige Seite. Ich interessiere mich besonders für Sportnachrichten und für die Politik in Kamerun. Mein Vater liest lieber die Zeitung, und wir diskutieren oft darüber, welches Medium besser ist. Ich weiß nicht, warum so viele Menschen Falschmeldungen glauben. Meiner Meinung nach muss jeder lernen, wie man gute Quellen erkennt.}

« Je m'informe chaque jour sur Internet et j'écoute parfois la radio. D'abord, je me demande toujours si une nouvelle est vraiment vraie. Ensuite, je compare plusieurs sources, car je ne me fie pas à un seul site. Je m'intéresse surtout aux nouvelles sportives et à la politique au Cameroun. Mon père préfère lire le journal, et nous discutons souvent de quel média est le meilleur. Je ne sais pas pourquoi tant de gens croient les fausses informations. À mon avis, chacun doit apprendre comment on reconnaît les bonnes sources. »

Vérification des contraintes : interrogatives indirectes — \all{ob eine Nachricht wirklich wahr ist}, \all{welches Medium besser ist}, \all{warum so viele Menschen Falschmeldungen glauben}, \all{wie man gute Quellen erkennt} (quatre, donc largement deux) ; verbes à préposition — \all{sich verlassen auf + Akk.}, \all{sich interessieren für + Akk.}, \all{diskutieren über + Akk.} (trois, donc largement deux). Tous les verbes conjugués en subordonnée sont bien en finale.
\end{exoresolu}

\courssub{Production orale : mini-débat « Internet oder Zeitung — was ist besser? »}

\begin{experience}[Mini-débat en classe]
Par groupes de quatre, organisez un débat de cinq minutes : deux élèves défendent Internet, deux défendent le journal papier. Chaque élève parle au moins deux fois. Utilisez obligatoirement la grille d'expression d'opinion ci-dessous. Un rapporteur note les arguments, puis la classe vote.

\textbf{Grille d'expression d'opinion :}
\begin{itemize}
\item Donner son avis : \all{Meiner Meinung nach ...} / \all{Ich denke, dass ...} / \all{Ich finde, dass ...}
\item Être d'accord : \all{Ich stimme dir zu.} (je suis d'accord avec toi) / \all{Das sehe ich auch so.} (je le vois aussi ainsi)
\item Ne pas être d'accord : \all{Ich bin anderer Meinung.} (je suis d'un autre avis) / \all{Das stimmt nicht ganz.} (ce n'est pas tout à fait exact)
\item Argumenter : \all{weil ...} / \all{ein Vorteil ist, dass ...} (un avantage est que...) / \all{ein Nachteil ist, dass ...} (un inconvénient est que...)
\item Nuancer : \all{einerseits ... andererseits ...} (d'une part... d'autre part...)
\end{itemize}

Exemple de réplique : \all{Ich bin anderer Meinung: Meiner Meinung nach ist die Zeitung zuverlässiger, weil Journalisten ihre Quellen überprüfen.} « Je ne suis pas d'accord : à mon avis, le journal est plus fiable, parce que les journalistes vérifient leurs sources. »
\end{experience}

\begin{center}
\begin{tabularx}{\linewidth}{|X|X|}
\hline
\rowcolor{popgreenL} \textbf{Internet — Vorteile} & \textbf{Zeitung — Vorteile} \\
\hline
\all{schnell und aktuell} (rapide et actuel) & \all{geprüfte Quellen} (sources vérifiées) \\
\hline
\all{kostenlos oder billig} (gratuit ou bon marché) & \all{seriöse Journalisten} (journalistes sérieux) \\
\hline
\all{überall verfügbar, auch in Yaoundé per Handy} (disponible partout, même à Yaoundé par téléphone) & \all{keine Ablenkung durch Werbung} (pas de distraction par la publicité) \\
\hline
\all{viele Themen: Sport, Politik, Musik} (beaucoup de sujets : sport, politique, musique) & \all{gut für die Konzentration} (bon pour la concentration) \\
\hline
\end{tabularx}
\end{center}

\courssub{Grille d'auto-évaluation}

Après chaque production (écrite ou orale), évaluez-vous honnêtement avec cette grille. Notez chaque critère de 0 (absent) à 2 (maîtrisé).

\begin{center}
\begin{tabularx}{\linewidth}{|X|c|X|}
\hline
\rowcolor{popblueL} \textbf{Critère} & \textbf{Note /2} & \textbf{Question de contrôle} \\
\hline
Vocabulaire du thème (médias, information) & & Ai-je utilisé au moins quatre mots du lexique (\all{die Quelle, die Nachricht, die Falschmeldung, der Journalist}...) ? \\
\hline
Grammaire : interrogatives indirectes & & Ai-je placé le verbe à la fin après \all{ob} et les \all{w-Fragen} ? \\
\hline
Grammaire : verbes à préposition & & Ai-je choisi la bonne préposition et le bon cas (\all{über + Akk.}, \all{für + Akk.}, \all{auf + Akk.}) ? \\
\hline
Structure (Einleitung, Résumé, Analyse, Meinung) & & Mon texte suit-il les quatre étapes avec les connecteurs ? \\
\hline
Pertinence et avis personnel & & Ai-je répondu au sujet et osé donner mon opinion ? \\
\hline
\end{tabularx}
\end{center}

Objectif : 8/10 minimum. En dessous de 6/10, reprenez la méthode en quatre étapes et la liste des verbes à préposition avant de réécrire votre production.

\begin{attention}[Erreurs fréquentes dans les productions des francophones]
\begin{itemize}
\item \all{Ich informiere mich \textbf{sur} die Politik} : impossible ! La préposition française « sur » ne se traduit pas mot à mot ; le verbe allemand impose \all{über} + Akkusativ : \all{ich informiere mich \textbf{über} die Politik}. En revanche, pour le \emph{moyen} d'information, dites \all{im Internet}, \all{in der Zeitung} (datif).
\item \all{Meiner Meinung nach \textbf{ich denke}...} : double faute. Après \all{meiner Meinung nach}, le verbe conjugué suit directement : \all{Meiner Meinung nach \textbf{ist} das Internet schneller.}
\item Oublier la majuscule aux noms : \all{die zeitung}, \all{das internet} sont fautifs — écrivez \all{die Zeitung}, \all{das Internet}.
\item Recopier le texte support dans le résumé au lieu de reformuler : c'est la faute la plus sanctionnée au Bac.
\end{itemize}
\end{attention}

\begin{savaistu}[L'expression personnelle rapporte des points au Bac]
Au baccalauréat camerounais, les correcteurs valorisent explicitement l'\emph{expression personnelle} : un candidat qui ose donner son avis avec \all{ich denke, dass ...} ou \all{meiner Meinung nach ...} obtient des points supplémentaires, même si sa phrase contient une petite erreur. Un avis simple mais personnel vaut mieux qu'un résumé parfait mais sans opinion. Alors osez : \all{Meiner Meinung nach ist das Thema sehr aktuell, weil auch in Kamerun viele junge Leute sich nur noch über das Handy informieren.} « À mon avis, le sujet est très actuel, parce qu'au Cameroun aussi, beaucoup de jeunes ne s'informent plus que par téléphone. »
\end{savaistu}

\begin{exercice}[À vous de produire]
\begin{enumerate}
\item Rédigez un commentaire complet (8 à 10 lignes) du texte support « Informationen im digitalen Zeitalter » en suivant les quatre étapes et en employant au moins trois connecteurs de la liste : \all{zuerst, dann, außerdem, schließlich, meiner Meinung nach}.
\item Écrivez un résumé du même texte en exactement trois phrases, au Präsens, sans recopier aucune phrase du texte.
\item Production libre : « Wie informierst du dich? » — 6 à 8 lignes avec deux interrogatives indirectes et deux verbes à préposition, puis auto-évaluez-vous avec la grille.
\end{enumerate}
\end{exercice}

\begin{corrige}[Exercice — éléments de correction]
\begin{enumerate}
\item Le commentaire doit s'ouvrir sur \all{Der Text handelt von der Information im digitalen Zeitalter}, contenir un résumé au Präsens avec \all{zuerst / dann / außerdem}, une analyse (par exemple l'interrogative indirecte \all{Er fragt sich, ob die Quellen zuverlässig sind}), et se clore sur un avis avec \all{meiner Meinung nach} ou \all{ich denke, dass} introduit par \all{schließlich}. Vérifiez la deuxième place du verbe après \all{meiner Meinung nach}.
\item Résumé possible : \all{Der Text beschreibt, wie Journalisten heute arbeiten. Sie prüfen ihre Quellen, bevor sie schreiben. Er warnt vor Falschmeldungen im Internet.} — trois phrases, Präsens, aucune citation, verbes en finale dans les subordonnées.
\item Production : comptez vos interrogatives indirectes (verbe en finale après \all{ob} ou le mot interrogatif) et vos verbes à préposition (préposition correcte + cas correct : \all{über/für/auf + Akkusativ}). Attribuez-vous une note sur 10 avec la grille ; toute case à 0 ou 1 doit être retravaillée.
\end{enumerate}
\end{corrige}

\newpage

\coursec{Activité d'intégration}

\courssub{Objectifs de l'activité}

Cette activité d'intégration clôt la leçon sur la gestion de l'information. Elle invite les élèves à mobiliser, dans une tâche de communication authentique et ludique, l'ensemble des acquis du chapitre :

\begin{itemize}
\item le \cle{vocabulaire thématique} des médias et de l'information (\all{die Nachricht}, \all{die Quelle}, \all{der Journalist}, \all{die Zeitung}, \all{das Internet}, \all{überprüfen}, \all{veröffentlichen}) ;
\item la \cle{compréhension écrite} : lire et évaluer une information trouvée sur Internet ;
\item la \cle{grammaire 1} : formuler des interrogatives indirectes avec \all{ob} et les pronoms interrogatifs (\all{woher}, \all{wann}, \all{wer}, \all{warum}) ;
\item la \cle{grammaire 2} : employer correctement les verbes à complément prépositionnel (\all{sich informieren über + Akk.}, \all{berichten über + Akk.}, \all{fragen nach + Dat.}, \all{zweifeln an + Dat.}) ;
\item l'\cle{expression écrite et orale} : rédiger un court article de presse et le présenter devant la classe.
\end{itemize}

Au-delà de la langue, l'activité développe l'\cle{esprit critique} : savoir vérifier une information avant de la partager est une compétence de vie essentielle à l'ère des réseaux sociaux, et une compétence explicitement visée par le programme de Terminale A.

\courssub{Situation}

Vous êtes élèves en classe de Terminale A au lycée de Yaoundé. Votre club de journalisme scolaire, „Die Brücke“, prépare son prochain numéro. Ce matin, un membre du club arrive très excité : il a lu sur Internet une information spectaculaire qu'il voudrait publier immédiatement. Mais le rédacteur en chef est prudent : avant toute publication, la rédaction doit vérifier la nouvelle. Vous formez la rédaction du journal. À vous de jouer !

\begin{textebox}[La nouvelle à vérifier — message reçu sur WhatsApp]
\textbf{Sensationelle Nachricht!!!}\\[2mm]
„Große Entdeckung in Kamerun: Ein deutscher Wissenschaftler hat in der Nähe von Douala eine Pflanze gefunden, die gegen alle Krankheiten hilft! Die Regierung will das Geheimnis verstecken. Teile diese Nachricht schnell mit allen deinen Freunden, bevor sie gelöscht wird!!!“\\[2mm]
\emph{(Quelle : unbekannter Blog, kein Autor, kein Datum, kein Foto)}
\end{textebox}

\begin{definition}[Glossaire du message]
\all{die Entdeckung, -en} : la découverte ; \all{der Wissenschaftler, -} : le scientifique, le chercheur ; \all{die Nähe} (in der Nähe von + Dat.) : à proximité de ; \all{die Pflanze, -n} : la plante ; \all{die Krankheit, -en} : la maladie ; \all{das Geheimnis, -se} : le secret ; \all{verstecken} : cacher ; \all{teilen} : partager ; \all{löschen} : effacer, supprimer.
\end{definition}

\begin{attention}[Une nouvelle suspecte]
Ce message présente tous les signes d'une \cle{fausse information} (\all{die Falschmeldung, -en} ; on dit aussi couramment \all{die Fake News}, mot anglais intégré à l'usage journalistique allemand) : source inconnue, pas d'auteur, pas de date, ton sensationnel, appel urgent au partage. Votre travail de journalistes consiste à le démontrer avec des arguments et des questions précises.
\end{attention}

\courssub{Consignes}

Travail par groupes de 4 élèves. Chaque groupe choisit ses rôles : un élève apporte la nouvelle (\all{der Informant, -en}), deux élèves sont journalistes (\all{die Journalisten}) et un élève est rédacteur en chef (\all{der Chefredakteur, -e}) qui prend la décision finale.

\begin{enumerate}
\item \textbf{Lecture et première réaction (5 min).} Lisez le message en silence. Soulignez les mots du vocabulaire des médias que vous connaissez et relevez ce qui vous semble suspect. Répondez d'abord globalement : \all{Worum geht es in dieser Nachricht?} (De quoi parle cette nouvelle ?)

\item \textbf{Préparation des questions (10 min).} Les journalistes préparent par écrit au moins \textbf{cinq questions} à poser à l'informant. Deux questions directes au moins doivent être transformées en \textbf{interrogatives indirectes} introduites par \all{ob} ou par un pronom interrogatif. Modèles : \all{Wir fragen uns, ob die Quelle zuverlässig ist.} — \all{Sag uns bitte, woher die Information kommt!} — \all{Wir möchten wissen, wer den Artikel geschrieben hat.}

\item \textbf{L'interview (10 min).} Jeu de rôle : les journalistes interrogent l'informant. Celui-ci répond en improvisant (il peut défendre sa nouvelle ou avouer ses doutes). Le rédacteur en chef écoute et note les arguments. Utilisez les verbes à préposition : \all{fragen nach + Dat.}, \all{sich informieren über + Akk.}, \all{zweifeln an + Dat.}, \all{denken an + Akk.}

\item \textbf{La décision (5 min).} Le rédacteur en chef annonce la décision du groupe et la justifie en deux ou trois phrases : \all{Wir veröffentlichen die Nachricht (nicht), weil ...} Pensez à la place du verbe après \all{weil} !

\item \textbf{Rédaction de l'article (15 min).} Chaque groupe rédige un court article d'environ \textbf{6 lignes} pour „Die Brücke“. L'article doit contenir : un titre, au moins \textbf{deux interrogatives indirectes}, au moins \textbf{deux verbes à complément prépositionnel} et au moins \textbf{cinq mots du lexique des médias}. Débutez par exemple par : \all{Wir berichten über eine Nachricht aus dem Internet ...} ou \all{Wir haben uns über eine Pflanze bei Douala informiert ...}

\item \textbf{Présentation et vote (10 min).} Chaque groupe lit son article devant la classe. La classe vote ensuite pour l'article le plus fiable et le mieux rédigé (\all{der zuverlässigste Artikel}). Le professeur corrige à l'oral les erreurs de grammaire les plus fréquentes.
\end{enumerate}

\courssub{Ressources à mobiliser}

\begin{aretenir}[Boîte à outils linguistique]
\textbf{Lexique des médias :} \all{die Information, -en} ; \all{die Nachricht, -en} ; \all{die Quelle, -n} ; \all{der Journalist, -en / die Journalistin, -nen} ; \all{die Zeitung, -en} ; \all{der Artikel, -} ; \all{das Internet} ; \all{die Recherche, -n} (l'enquête journalistique) ; \all{zuverlässig} (fiable) ; \all{überprüfen} (vérifier) ; \all{veröffentlichen} (publier) ; \all{die Falschmeldung, -en}.\\[2mm]
\textbf{Interrogative indirecte :} verbe conjugué \textbf{en fin de phrase}. \all{Ich weiß nicht, ob die Nachricht stimmt.} / \all{Er fragt, woher die Information kommt.}\\[2mm]
\textbf{Verbes à préposition :} \all{sich informieren über + Akk.} ; \all{berichten über + Akk.} ; \all{fragen nach + Dat.} ; \all{zweifeln an + Dat.} ; \all{denken an + Akk.} ; \all{sich freuen über + Akk.} ; \all{warnen vor + Dat.}\\[2mm]
\textbf{Justifier une décision :} \all{weil} + verbe en fin : \all{Wir veröffentlichen sie nicht, weil die Quelle unbekannt ist.}
\end{aretenir}

\begin{methode}[Vérifier une information en cinq questions]
\begin{enumerate}
\item \textbf{Qui ?} \all{Wer hat den Artikel geschrieben?} — Un auteur identifié est un premier signe de sérieux.
\item \textbf{D'où ?} \all{Woher kommt die Information?} — La source est-elle un média connu ou un blog anonyme ?
\item \textbf{Quand ?} \all{Wann wurde die Nachricht veröffentlicht?} — Sans date, impossible de situer les faits.
\item \textbf{Pourquoi ?} \all{Warum soll ich die Nachricht teilen?} — Un appel urgent au partage est un signal d'alarme.
\item \textbf{Recouper :} \all{Berichten andere Zeitungen auch darüber?} — Une vraie nouvelle est reprise par plusieurs médias fiables.
\end{enumerate}
\end{methode}

\begin{attention}[Pièges à éviter]
\begin{itemize}
\item Ne dites pas \all{*Ich frage, ob kommt die Information aus dem Internet*} : dans l'interrogative indirecte, le verbe conjugué va à la fin : \all{..., ob die Information aus dem Internet kommt.}
\item \all{die Information} se construit avec \all{über + Akkusativ} : \all{eine Information über die Pflanze}, jamais avec \all{von} sur le modèle français « une information de ».
\item Attention au faux ami : \all{aktualisieren} signifie « mettre à jour », et non « rendre actuel » ; pour « vérifier une nouvelle », utilisez \all{eine Nachricht überprüfen}.
\item N'oubliez pas la majuscule à tous les noms : \all{die Quelle}, \all{das Internet}, \all{der Artikel}.
\item \all{fragen nach} exige le datif : \all{Wir fragen den Informanten nach der Quelle} (et non \all{*nach die Quelle*}).
\end{itemize}
\end{attention}

\courssub{Proposition de production et corrigé commenté}

\begin{exoresolu}[Modèle d'article pour „Die Brücke“]
Rédigez un article d'environ 6 lignes qui rend compte de l'enquête de votre rédaction sur la nouvelle de la « plante miracle » de Douala.
\tcblower
\textbf{Eine „Wunderpflanze“ bei Douala? Unsere Recherche}\\[2mm]
\all{Wir berichten heute über eine Nachricht aus dem Internet: Ein unbekannter Blog schreibt, dass ein deutscher Wissenschaftler bei Douala eine Pflanze gegen alle Krankheiten gefunden hat. Wir haben uns über diese Information genau informiert und den Informanten nach der Quelle gefragt. Wir wollten wissen, wer den Artikel geschrieben hat und ob die Quelle zuverlässig ist. Der Blog nennt aber keinen Autor und kein Datum. Deshalb zweifeln wir an dieser Nachricht: Wir veröffentlichen sie nicht, weil sie eine Falschmeldung sein kann. Unser Tipp an alle Leser: Überprüft jede Information, bevor ihr sie teilt!}\\[3mm]
\emph{Commentaire des choix de langue :}
\begin{itemize}
\item \textbf{Lexique des médias} bien présent : \all{die Nachricht}, \all{der Blog}, \all{die Information}, \all{die Quelle}, \all{der Autor}, \all{die Falschmeldung}, \all{veröffentlichen}, \all{überprüfen} — plus de cinq mots du champ lexical, comme demandé.
\item \textbf{Deux interrogatives indirectes} correctement construites, verbe en finale : \all{..., wer den Artikel geschrieben hat} (au Perfekt dans une subordonnée, le participe \all{geschrieben} précède l'auxiliaire \all{hat} placé en toute fin) et \all{..., ob die Quelle zuverlässig ist}.
\item \textbf{Quatre verbes à préposition} avec le bon cas : \all{berichten über + Akk.} (\all{über eine Nachricht}), \all{sich informieren über + Akk.} (\all{über diese Information}), \all{fragen nach + Dat.} (\all{nach der Quelle}), \all{zweifeln an + Dat.} (\all{an dieser Nachricht}).
\item \textbf{Subordonnée avec \all{weil}} : verbe conjugué à la fin — \all{..., weil sie eine Falschmeldung sein kann} (l'infinitif \all{sein} précède le modal \all{kann}).
\item \textbf{Style journalistique} : titre avec point d'interrogation pour susciter l'intérêt, conclusion avec un conseil au lecteur (\all{Unser Tipp}), impératif pluriel correct \all{Überprüft ... teilt!}.
\end{itemize}
\end{exoresolu}

\courssub{Bilan de l'activité}

À l'issue de cette activité, chaque élève a mis en pratique l'ensemble des compétences de la leçon dans une situation proche de la vie réelle : lire et évaluer une information, poser des questions pertinentes, argumenter une décision et rédiger un court texte journalistique. Le tableau récapitulatif ci-dessous permet à chacun de s'auto-évaluer.

\begin{center}
\begin{tabularx}{\linewidth}{|X|X|}
\hline
\rowcolor{popblueL} \textbf{Savoir-faire} & \textbf{Je l'ai réussi si...} \\
\hline
Compréhension & j'ai identifié le sujet du message et ses éléments suspects. \\
\hline
Lexique & j'ai employé au moins cinq mots du champ des médias avec le bon article. \\
\hline
Interrogatives indirectes & mes phrases avec \all{ob} et les mots en \all{w-} ont le verbe à la fin. \\
\hline
Verbes à préposition & j'ai utilisé la bonne préposition et le bon cas (\all{über + Akk.}, \all{nach + Dat.}, \all{an + Dat.}). \\
\hline
Production écrite et orale & mon article a un titre, 6 lignes, et je l'ai présenté clairement à la classe. \\
\hline
\end{tabularx}
\end{center}

\begin{savaistu}[L'éducation aux médias en Allemagne]
En Allemagne, l'éducation aux médias (\all{die Medienkompetenz}) fait partie des priorités scolaires : les élèves apprennent à reconnaître les \all{Fake News} en vérifiant la source, l'auteur et la date d'un article. De grands journaux proposent même des éditions spéciales pour les écoles, et les radios publiques offrent des dossiers pédagogiques sur la vérification des faits (\all{der Faktencheck, -s}). Au Cameroun aussi, où WhatsApp et Facebook sont des sources d'information très utilisées, savoir vérifier une nouvelle avant de la partager est devenu un réflexe citoyen indispensable. En jouant les journalistes aujourd'hui, vous avez fait exactement le même travail que les professionnels des rédactions de Berlin, Vienne ou Yaoundé !
\end{savaistu}

Pour prolonger l'activité à la maison, chaque élève peut choisir une vraie nouvelle de la semaine (radio, télévision ou Internet) et rédiger cinq lignes en allemand en commençant par \all{Ich habe mich über ... informiert}, avec une interrogative indirecte et un verbe à préposition. Ce court texte sera lu au début de la séance suivante.

\newpage

\coursec{Méthodes et exercices résolus}

\courssub{Méthode 1 : Comprendre un texte sur la gestion de l'information}

\begin{methode}[Lire et comprendre un texte de presse allemand]
\begin{enumerate}
\item \textbf{Première lecture globale.} Lis le titre, la source et le premier paragraphe : ils annoncent le thème. Pose-toi trois questions : \all{Worum geht es?} (de quoi s'agit-il ?), \all{Wer spricht?} (qui parle ?), \all{Wo und wann?} (où et quand ?).
\item \textbf{Repérage des mots-clés du champ lexical.} Souligne les mots du thème : \all{die Information}, \all{die Nachricht}, \all{die Quelle}, \all{der Journalist}, \all{die Zeitung}, \all{das Internet}. Devine le sens des mots inconnus par le contexte et les familles de mots (\all{informieren} $\rightarrow$ \all{die Information} ; \all{berichten} $\rightarrow$ \all{der Bericht}).
\item \textbf{Lecture détaillée.} Repère les connecteurs (\all{aber}, \all{denn}, \all{deshalb}, \all{trotzdem}) : ils signalent l'argumentation. Distingue les faits (chiffres, exemples) des opinions (\all{ich meine}, \all{man sollte}, \all{meiner Meinung nach}).
\item \textbf{Réponse aux questions.} Réponds toujours par une \textbf{phrase complète} en reprenant les mots de la question, jamais par un seul mot. Cite le texte si on te le demande : \all{Im Text steht: „...“}.
\item \textbf{Résumé.} Pour résumer : une phrase par idée principale, au présent, avec tes propres mots, sans exemple ni détail.
\end{enumerate}
\end{methode}

\begin{exoresolu}[Compréhension de texte — type Bac]
\textbf{Texte :} \all{Viele Jugendliche in Kamerun informieren sich heute nicht mehr über Zeitungen, sondern über das Internet und die sozialen Netzwerke. Doch nicht jede Nachricht im Netz ist wahr. Journalisten warnen deshalb vor falschen Informationen und raten, immer die Quelle zu prüfen.}\par
\textbf{Questions :} 1. \all{Wie informieren sich viele Jugendliche heute?} 2. \all{Warum warnen Journalisten?} 3. \all{Was raten die Journalisten?}
\tcblower
\textbf{Raisonnement :} la question 1 porte sur le moyen d'information (première phrase) ; la question 2 sur la cause (signalée par \all{Doch} et \all{deshalb}) ; la question 3 sur le conseil (verbe \all{raten}).\par
\textbf{Réponses rédigées :}\par
1. \all{Viele Jugendliche informieren sich heute über das Internet und die sozialen Netzwerke.} « Beaucoup de jeunes s'informent aujourd'hui par Internet et les réseaux sociaux. »\par
2. \all{Journalisten warnen, weil nicht jede Nachricht im Netz wahr ist.} « Les journalistes mettent en garde parce que toutes les nouvelles sur le Net ne sont pas vraies. » (Subordonnée avec \all{weil} : verbe conjugué à la fin.)\par
3. \all{Die Journalisten raten, immer die Quelle zu prüfen.} « Les journalistes conseillent de toujours vérifier la source. » (Tour infinitif avec \all{zu} après \all{raten}.)\par
\textbf{Conclusion :} chaque réponse reprend les termes de la question, forme une phrase complète et reste fidèle au texte.
\end{exoresolu}

\courssub{Méthode 2 : Employer le vocabulaire des médias}

\begin{methode}[Construire et utiliser le lexique thématique]
\begin{enumerate}
\item \textbf{Apprends chaque nom avec son article et son pluriel} : \all{die Information, -en} ; \all{die Nachricht, -en} ; \all{die Quelle, -n} ; \all{der Journalist, -en} (nom faible : \all{den Journalisten} à l'accusatif !) ; \all{die Zeitung, -en} ; \all{das Internet, -s}.
\item \textbf{Mémorise les familles de mots} : \all{informieren} (informer) $\rightarrow$ \all{sich informieren über + Akk.} (s'informer sur) $\rightarrow$ \all{die Information} ; \all{berichten} (rapporter) $\rightarrow$ \all{der Bericht, -e} ; \all{prüfen} (vérifier) $\rightarrow$ \all{die Prüfung, -en}.
\item \textbf{Retiens les verbes avec leur construction} : \all{eine Nachricht verbreiten} (diffuser une nouvelle), \all{eine Quelle prüfen} (vérifier une source), \all{über etwas berichten} (faire un reportage sur quelque chose).
\item \textbf{Réemploie le lexique} dans tes propres phrases : un mot n'est acquis que s'il est utilisé dans un contexte camerounais ou personnel : \all{Ich lese jeden Tag die Zeitung.} « Je lis le journal tous les jours. »
\end{enumerate}
\end{methode}

\begin{definition}[die Nachricht ou die Information ?]
\all{die Nachricht, -en} désigne une nouvelle précise, une information ponctuelle (au pluriel : « les informations » à la radio, à la télévision) : \all{die Nachrichten um 20 Uhr} « le journal de 20 heures ». \all{die Information, -en} désigne un renseignement, une donnée ; en allemand, on l'emploie très souvent au \textbf{singulier} là où le français met « des informations » : \all{Ich brauche eine Information.} « J'ai besoin d'un renseignement. »
\end{definition}

\begin{center}
\begin{tabularx}{\linewidth}{|X|X|X|}\hline
\rowcolor{popblueL} Français & Deutsch & Remarque \\ \hline
l'information, le renseignement & die Information, -en & souvent au singulier \\ \hline
la nouvelle ; les informations (radio) & die Nachricht, -en & \all{die Nachrichten} = le journal télévisé/radiophonique \\ \hline
la source & die Quelle, -n & \all{eine Quelle prüfen} \\ \hline
le journaliste / la journaliste & der Journalist, -en / die Journalistin, -nen & nom faible au masculin ; féminin en \all{-in} \\ \hline
le journal & die Zeitung, -en & \all{eine Zeitung lesen} \\ \hline
Internet & das Internet, -s & \all{im Internet} = sur Internet \\ \hline
\end{tabularx}
\end{center}

\begin{exoresolu}[Lexique — phrases à compléter]
Complète avec le mot qui convient : \all{Quelle} — \all{Nachrichten} — \all{Journalistin} — \all{informiert} — \all{Zeitung}.\par
1. \all{Mein Vater liest jeden Morgen die \ldots.} 2. \all{Im Radio hören wir um 20 Uhr die \ldots.} 3. \all{Die \ldots{} schreibt einen Artikel über den Markt in Douala.} 4. \all{Man muss immer die \ldots{} prüfen.} 5. \all{Er \ldots{} sich über das Internet.}
\tcblower
\textbf{Raisonnement :} on identifie le champ lexical et la construction grammaticale de chaque phrase.\par
1. \all{die Zeitung} : on « lit » un journal ; accusatif féminin, article \all{die} inchangé.\par
2. \all{die Nachrichten} : au pluriel, « les informations » que l'on écoute à la radio.\par
3. \all{Die Journalistin} : sujet féminin (une journaliste) ; le suffixe \all{-in} marque le féminin des métiers.\par
4. \all{die Quelle} : complément du verbe \all{prüfen} ; collocation à retenir : \all{eine Quelle prüfen}.\par
5. \all{informiert} : verbe réfléchi \all{sich informieren über + Akk.}, conjugué à la 3\up{e} personne : \all{er informiert sich}.\par
\textbf{Conclusion :} bien choisir le mot exige de connaître article, genre, pluriel et construction verbale.
\end{exoresolu}

\courssub{Méthode 3 : Construire une interrogative indirecte}

\begin{propriete}[La question indirecte : ob ou mot en w- + verbe à la fin]
La question indirecte est une \textbf{subordonnée} introduite après un verbe comme \all{wissen}, \all{fragen}, \all{sagen}, \all{verstehen}, \all{sich fragen}.\par
$\bullet$ Question directe \textbf{totale} (réponse oui/non, sans mot interrogatif) $\rightarrow$ subordonnée introduite par \all{ob} (« si »).\par
$\bullet$ Question directe \textbf{partielle} (avec mot interrogatif) $\rightarrow$ on conserve le mot en w- (\all{wer}, \all{was}, \all{wann}, \all{wo}, \all{woher}, \all{wie}, \all{warum}, \all{wieviel}...).\par
$\bullet$ Dans les deux cas : \textbf{virgule} avant la subordonnée et \textbf{verbe conjugué en dernière position}.
\end{propriete}

\begin{exemplebox}[De la question directe à la question indirecte]
\all{Ist die Quelle sicher?} $\rightarrow$ \all{Ich weiß nicht, ob die Quelle sicher ist.} « Je ne sais pas si la source est fiable. »\par
\all{Woher kommt die Information?} $\rightarrow$ \all{Ich weiß nicht, woher die Information kommt.} « Je ne sais pas d'où vient l'information. »\par
\all{Warum prüft er die Quelle?} $\rightarrow$ \all{Er sagt nicht, warum er die Quelle prüft.} « Il ne dit pas pourquoi il vérifie la source. »
\end{exemplebox}

\begin{methode}[Former une question indirecte avec ob ou un mot en w-]
\begin{enumerate}
\item \textbf{Identifie le type de question directe.} Question totale (réponse oui/non) $\rightarrow$ \all{ob}. Question partielle (avec \all{wer}, \all{was}, \all{wann}, \all{wo}, \all{wie}, \all{warum}...) $\rightarrow$ on garde le mot interrogatif.
\item \textbf{Introduis la subordonnée} après un verbe comme \all{wissen}, \all{fragen}, \all{sagen}, \all{verstehen} : \all{Ich weiß nicht, ...}
\item \textbf{Mets le verbe conjugué à la fin} de la subordonnée : c'est la règle d'or. \all{Woher kommt die Information?} $\rightarrow$ \all{Ich weiß nicht, woher die Information kommt.}
\item \textbf{Attention à l'ordre sujet-verbe} : le sujet suit immédiatement \all{ob} ou le mot en w- ; le verbe conjugué se place en dernière position.
\item \textbf{Virgule obligatoire} avant la subordonnée, comme pour toutes les subordonnées allemandes.
\end{enumerate}
\end{methode}

\begin{exoresolu}[Grammaire — transformation en interrogative indirecte]
Transforme en question indirecte après \all{Ich weiß nicht, ...}\par
1. \all{Woher bekommt der Journalist seine Informationen?} 2. \all{Ist diese Nachricht wahr?} 3. \all{Wann erscheint die Zeitung?}
\tcblower
\textbf{Raisonnement :} phrases 1 et 3 = questions partielles (on garde \all{woher}, \all{wann}) ; phrase 2 = question totale (on introduit \all{ob}). Dans tous les cas, verbe conjugué à la fin.\par
1. \all{Ich weiß nicht, woher der Journalist seine Informationen bekommt.} « Je ne sais pas d'où le journaliste tient ses informations. » (\all{bekommt} en fin de phrase.)\par
2. \all{Ich weiß nicht, ob diese Nachricht wahr ist.} « Je ne sais pas si cette nouvelle est vraie. » (\all{ob} = « si » ; verbe \all{ist} à la fin.)\par
3. \all{Ich weiß nicht, wann die Zeitung erscheint.} « Je ne sais pas quand le journal paraît. »\par
\textbf{Conclusion :} question totale $\rightarrow$ \all{ob} ; question partielle $\rightarrow$ mot en w- ; dans les deux cas, verbe conjugué en position finale.
\end{exoresolu}

\courssub{Méthode 4 : Maîtriser les verbes à complément prépositionnel}

\begin{propriete}[Verbe + préposition + cas : un bloc à mémoriser]
Certains verbes allemands se construisent obligatoirement avec une préposition qui impose son cas. On les apprend comme un tout :\par
\all{sich informieren über + Akk.} (s'informer sur) ; \all{berichten über + Akk.} (faire un reportage sur) ; \all{warten auf + Akk.} (attendre) ; \all{bitten um + Akk.} (demander) ; \all{sich freuen über + Akk.} (se réjouir de) ; \all{warnen vor + Dat.} (mettre en garde contre).\par
Pour \textbf{interroger} sur la chose : \all{worüber}, \all{worauf}, \all{wovor} (\all{wo(r)-} + préposition). Pour \textbf{reprendre} la chose : \all{darüber}, \all{darauf}, \all{davor} (\all{da(r)-} + préposition). Le \all{-r-} s'insère devant une voyelle : \all{worauf}, \all{darauf}, mais \all{wovor}, \all{davor}.
\end{propriete}

\begin{methode}[Employer un verbe avec sa préposition]
\begin{enumerate}
\item \textbf{Apprends le verbe avec sa préposition et son cas} comme un tout : \all{sich informieren über + Akk.}, \all{warten auf + Akk.}, \all{warnen vor + Dat.}, \all{berichten über + Akk.}, \all{bitten um + Akk.}, \all{sich freuen über + Akk.}
\item \textbf{Applique le cas exigé par la préposition} : \all{über} et \all{um} $\rightarrow$ accusatif ; \all{vor} (au sens figuré) $\rightarrow$ datif. \all{Der Journalist warnt vor falschen Nachrichten.} (\all{falschen} = datif pluriel.)
\item \textbf{Pour interroger sur la chose}, utilise \all{worüber}, \all{worauf}, \all{wovor} : \all{Worüber informierst du dich?} « Sur quoi t'informes-tu ? »
\item \textbf{Pour reprendre la chose}, utilise \all{darüber}, \all{darauf}, \all{davor} : \all{Er berichtet darüber.} « Il en fait un reportage. »
\item \textbf{Ne traduis jamais mot à mot} : la préposition allemande ne correspond pas toujours au français (\all{warten auf} = « attendre », sans préposition en français).
\end{enumerate}
\end{methode}

\begin{exoresolu}[Conjugaison — verbes à complément prépositionnel]
Conjugue et complète avec la bonne préposition et le bon cas.\par
1. \all{Die Jugendlichen (sich informieren) \ldots{} das Internet.} 2. \all{Der Journalist (warnen) \ldots{} den falschen Informationen.} 3. \all{Wir (warten) \ldots{} die Nachrichten.} 4. \all{Sie (berichten) \ldots{} das Fußballspiel in Yaoundé.}
\tcblower
\textbf{Raisonnement :} on conjugue le verbe au présent, puis on place la préposition apprise avec le verbe et on met le nom au cas exigé.\par
1. \all{Die Jugendlichen informieren sich über das Internet.} « Les jeunes s'informent par Internet. » (\all{sich informieren über + Akk.} ; \all{das Internet} reste \all{das} à l'accusatif.)\par
2. \all{Der Journalist warnt vor den falschen Informationen.} « Le journaliste met en garde contre les fausses informations. » (\all{warnen vor + Dat.} ; \all{den} = datif pluriel, adjectif en \all{-en}.)\par
3. \all{Wir warten auf die Nachrichten.} « Nous attendons les informations. » (\all{warten auf + Akk.})\par
4. \all{Sie berichtet über das Fußballspiel in Yaoundé.} « Elle fait un reportage sur le match de football à Yaoundé. » (\all{berichten über + Akk.})\par
\textbf{Conclusion :} verbe + préposition + cas forment un bloc inséparable à mémoriser.
\end{exoresolu}

\courssub{Méthode 5 : Produire un court texte guidé}

\begin{methode}[Rédiger un court paragraphe sur les médias]
\begin{enumerate}
\item \textbf{Planifie en trois temps} : introduction (le thème), développement (2 ou 3 arguments ou exemples), conclusion (ton avis).
\item \textbf{Utilise les connecteurs} : \all{zuerst} (d'abord), \all{außerdem} (de plus), \all{aber} (mais), \all{deshalb} (c'est pourquoi), \all{meiner Meinung nach} (à mon avis).
\item \textbf{Intègre les points de grammaire de la leçon} : au moins une interrogative indirecte (\all{Ich weiß nicht, ob...}) et un verbe à complément prépositionnel (\all{sich informieren über}).
\item \textbf{Respecte l'ordre des mots} : verbe conjugué en deuxième position dans la principale, à la fin dans la subordonnée.
\item \textbf{Relis} : majuscules aux noms, orthographe (\all{ß}, Umlaute), accords, ponctuation.
\end{enumerate}
\end{methode}

\begin{exoresolu}[Expression écrite — court paragraphe (5 à 8 phrases)]
Rédige un court texte : \all{Wie informierst du dich? Sind alle Informationen im Internet wahr?}
\tcblower
\textbf{Raisonnement :} on suit le plan en trois temps, on réemploie le lexique (\all{sich informieren}, \all{die Quelle prüfen}, \all{die Nachricht}) et les deux points de grammaire.\par
\textbf{Production modèle :}\par
\all{Heute informieren sich viele Jugendliche über das Internet und die sozialen Netzwerke. Ich lese auch manchmal die Zeitung, aber das Handy ist schneller. Man weiß oft nicht, ob eine Nachricht im Netz wahr ist. Deshalb muss man immer die Quelle prüfen. Journalisten warnen vor falschen Informationen. Meiner Meinung nach sind die Medien sehr wichtig, aber man muss kritisch bleiben.}\par
« Aujourd'hui, beaucoup de jeunes s'informent par Internet et les réseaux sociaux. Je lis aussi parfois le journal, mais le téléphone portable est plus rapide. Souvent, on ne sait pas si une nouvelle sur le Net est vraie. C'est pourquoi il faut toujours vérifier la source. Les journalistes mettent en garde contre les fausses informations. À mon avis, les médias sont très importants, mais il faut rester critique. »\par
\textbf{Points de grammaire vérifiés :} interrogative indirecte \all{ob eine Nachricht im Netz wahr ist} (verbe \all{ist} à la fin) ; \all{sich informieren über}, \all{warnen vor + Dat.} ; verbe en deuxième position après \all{Heute}, \all{Deshalb}, \all{Meiner Meinung nach}.\par
\textbf{Conclusion :} un paragraphe court mais structuré, correct et personnel vaut mieux qu'un long texte fautif.
\end{exoresolu}

\begin{attention}[Les 5 erreurs qui coûtent des points au Bac]
\begin{enumerate}
\item \textbf{Verbe mal placé dans l'interrogative indirecte.} Écris \all{Ich weiß nicht, ob die Nachricht wahr ist.} et jamais \all{...ob die Nachricht ist wahr.} : le verbe conjugué va toujours à la fin de la subordonnée.
\item \textbf{Confondre \all{ob} et \all{wenn}.} \all{ob} = « si » (interrogation, doute) ; \all{wenn} = « quand / si » (condition, temps). \all{Ich weiß nicht, ob er kommt.} $\neq$ \all{Wenn er kommt, ...}
\item \textbf{Oublier la préposition du verbe ou se tromper de cas.} \all{sich informieren über + Akk.}, \all{warnen vor + Dat.} : \all{vor falschen Nachrichten} (datif pluriel en \all{-en}), jamais \all{vor falsche Nachrichten}.
\item \textbf{Faux amis et calques du français.} \all{die Information} est le plus souvent au singulier en allemand ; \all{bekommen} = « recevoir » (et non « devenir », qui se dit \all{werden}) ; « s'informer sur » se dit \all{sich informieren über}, jamais \all{sich informieren auf}.
\item \textbf{Majuscules et orthographe.} Tout nom prend la majuscule : \all{das Internet}, \all{die Quelle}, \all{der Journalist}. Soigne les Umlaute (\all{über}, \all{prüfen}) et n'oublie pas la virgule avant toute subordonnée.
\end{enumerate}
\end{attention}

\newpage

\coursec{Exercices}

\courssub{Je vérifie mes connaissances}

\begin{exercice}[QCM : le bon choix]
Entoure la bonne réponse. Une seule réponse est correcte.
\begin{enumerate}
\item \all{Ich weiß nicht, \dots\ die Nachricht wahr ist.} \quad a) \all{wenn} \quad b) \all{ob} \quad c) \all{dass}
\item \all{Er fragt, \dots\ der Journalist kommt.} \quad a) \all{wann kommt} \quad b) \all{wann} \quad c) \all{ob wann}
\item \all{Wir warten \dots\ eine Antwort.} \quad a) \all{für} \quad b) \all{über} \quad c) \all{auf}
\item \all{Sie informiert sich \dots\ die Ereignisse in Kamerun.} \quad a) \all{über + Akkusativ} \quad b) \all{von + Dativ} \quad c) \all{auf + Akkusativ}
\item Le pluriel de \all{die Zeitung} est : \quad a) \all{die Zeitungen} \quad b) \all{die Zeitunger} \quad c) \all{die Zeitung}
\end{enumerate}
\end{exercice}

\begin{exercice}[Vrai ou faux ? Justifie ta réponse]
Dis si les affirmations suivantes sont vraies (V) ou fausses (F) et justifie ta réponse par une règle ou un exemple.
\begin{enumerate}
\item Dans une interrogative indirecte, le verbe conjugué se place à la fin de la subordonnée.
\item \all{ob} sert à introduire une question indirecte qui commence par un mot interrogatif (\all{wer}, \all{was}, \all{wo}...).
\item Le verbe \all{sich freuen} se construit toujours avec la préposition \all{auf}.
\item \all{die Quelle} signifie « le journal ».
\item \all{warten auf} est suivi de l'accusatif.
\end{enumerate}
\end{exercice}

\begin{exercice}[Vocabulaire : à chaque mot son article et son pluriel]
Complète le tableau suivant avec l'article, le pluriel et la traduction française de chaque mot.
\begin{center}
\begin{tabularx}{\linewidth}{|X|l|l|X|}
\hline
\rowcolor{popblueL} Mot allemand & Article & Pluriel & Traduction \\
\hline
Information & & & \\
\hline
Nachricht & & & \\
\hline
Quelle & & & \\
\hline
Journalist & & & \\
\hline
Zeitung & & & \\
\hline
Internet & & & \\
\hline
Bericht & & & \\
\hline
Sender & & & \\
\hline
Artikel & & & \\
\hline
Leser & & & \\
\hline
\end{tabularx}
\end{center}
\end{exercice}

\begin{exercice}[Conjugaison : complète les verbes]
Conjugue les verbes entre parenthèses au \all{Präsens}, avec la bonne préposition et le bon cas.
\begin{enumerate}
\item Ich (warten) \dots\ \dots\ den Bus nach Douala.
\item Wir (sich informieren) \dots\ \dots\ die Nachrichten des Tages.
\item (du / sich freuen) \dots\ \dots\ dein Geburtstagsgeschenk?
\item Der Journalist (fragen) \dots\ \dots\ der Wahrheit der Information.
\item Die Schüler (sprechen) \dots\ \dots\ die neue Zeitung der Schule.
\item Sie (sich interessieren) \dots\ \dots\ Politik und Sport.
\end{enumerate}
\end{exercice}

\courssub{Je m'entraîne}

\begin{exercice}[Thème : traduis en allemand]
Traduis les phrases suivantes en allemand correct (attention à la place du verbe et aux prépositions).
\begin{enumerate}
\item Le journaliste attend une confirmation.
\item Je ne sais pas si la nouvelle est vraie.
\item Nous nous informons sur Internet.
\item Elle demande où le journaliste habite.
\item Les élèves se réjouissent des vacances.
\end{enumerate}
\end{exercice}

\begin{exercice}[Transformation : de la question directe à l'interrogative indirecte]
Transforme chaque question directe en interrogative indirecte introduite par \all{Er fragt, \dots} ou \all{Ich weiß nicht, \dots}.
\begin{enumerate}
\item \all{Woher kommt die Nachricht?} $\rightarrow$ \all{Er fragt, \dots}
\item \all{Ist die Information richtig?} $\rightarrow$ \all{Ich weiß nicht, \dots}
\item \all{Wann erscheint die Zeitung?} $\rightarrow$ \all{Er fragt, \dots}
\item \all{Hat der Journalist die Quelle genannt?} $\rightarrow$ \all{Ich weiß nicht, \dots}
\item \all{Wer hat den Artikel geschrieben?} $\rightarrow$ \all{Er fragt, \dots}
\item \all{Warum liest sie keine Zeitung?} $\rightarrow$ \all{Ich weiß nicht, \dots}
\end{enumerate}
\end{exercice}

\begin{exercice}[Texte à trous : prépositions et cas]
Complète le texte suivant avec la préposition qui convient (\all{auf}, \all{über}, \all{an}, \all{für}, \all{von}) et mets l'article au bon cas.

\all{Jeden Morgen informiere ich mich \dots\ (die) Nachrichten. Ich warte \dots\ (die) Zeitung, die um sieben Uhr kommt. Mein Bruder interessiert sich \dots\ (der) Sport, aber ich denke oft \dots\ (meine) Zukunft. Meine Mutter spricht gern \dots\ (die) Politik, und mein Vater hält nichts \dots\ (die) sozialen Netzwerke.}
\end{exercice}

\begin{exercice}[Appariements : le verbe et sa préposition]
Relie chaque verbe de la colonne de gauche à la préposition et au cas qui lui correspondent dans la colonne de droite, puis construis une phrase complète avec deux des verbes.
\begin{center}
\begin{tabularx}{\linewidth}{|X|X|}
\hline
\rowcolor{popblueL} Verbe & Préposition + cas \\
\hline
\all{warten} & \all{über + Akkusativ} \\
\hline
\all{sich freuen} (avant l'événement) & \all{auf + Akkusativ} \\
\hline
\all{sich informieren} & \all{an + Dativ} \\
\hline
\all{denken} & \all{für + Akkusativ} \\
\hline
\all{sich interessieren} & \all{auf + Akkusativ} \\
\hline
\all{sich freuen} (après l'événement) & \all{über + Akkusativ} \\
\hline
\end{tabularx}
\end{center}
\end{exercice}

\courssub{Je me prépare au Bac}

\begin{exercice}[Compréhension écrite : texte et questions]
Lis attentivement le texte suivant, puis réponds aux questions en allemand, par des phrases complètes. Pour les questions 3, 4 et 5, justifie ta réponse par une citation exacte du texte.

\begin{textebox}[Soziale Netzwerke in Deutschland]
In Deutschland informieren sich viele junge Menschen nicht mehr über die klassische Zeitung, sondern über soziale Netzwerke. Sie warten nicht auf die Abendnachrichten im Fernsehen, denn das Handy bringt die Information sofort. Aber diese schnelle Information hat auch Gefahren: Nicht jede Nachricht ist wahr, und viele Nutzer fragen nicht, woher die Meldung kommt. Journalisten warnen vor falschen Nachrichten, den sogenannten „Fake News“. Sie erklären, dass man immer die Quelle prüfen muss. Eine Lehrerin aus München sagt: „Ich weiß oft nicht, ob meine Schüler eine Nachricht wirklich verstehen. Deshalb sprechen wir im Unterricht über die Medien.“ Die Schüler lernen, wie man eine seriöse Quelle erkennt und wie man sich über ein Thema richtig informiert. Viele junge Deutsche interessieren sich trotzdem für Politik: Sie lesen Artikel im Internet, diskutieren mit Freunden und denken über die Zukunft ihres Landes nach. Die Medienkompetenz, also der kluge Umgang mit Informationen, ist heute ein wichtiges Schulfach geworden.
\end{textebox}

\textbf{Questions :}
\begin{enumerate}
\item De quoi parle ce texte ? Résume le thème général en une ou deux phrases. (Question globale)
\item Quel est le problème principal que présente le texte ? (Question globale)
\item Comment les jeunes Allemands s'informent-ils aujourd'hui ? Justifie par une citation.
\item Que conseillent les journalistes ? Justifie par une citation.
\item Pourquoi la professeure de Munich parle-t-elle des médias en classe ? Justifie par une citation.
\end{enumerate}

\textbf{Questions de langue :}
\begin{enumerate}
\item Relève dans le texte deux interrogatives indirectes et souligne le verbe conjugué de chaque subordonnée.
\item Relève trois verbes à complément prépositionnel et indique pour chacun la préposition et le cas.
\end{enumerate}
\end{exercice}

\begin{exercice}[Thème et version]
\textbf{A. Thème.} Traduis les phrases suivantes en allemand.
\begin{enumerate}
\item Le journaliste attend une confirmation de sa source.
\item Je ne sais pas si la nouvelle est vraie.
\item Nous nous informons sur Internet et dans les journaux.
\item Elle demande quand l'article paraît.
\item Beaucoup de jeunes s'intéressent à la politique, mais ils ne pensent pas toujours aux dangers des réseaux sociaux.
\end{enumerate}

\textbf{B. Version.} Traduis le paragraphe suivant en français soigné.

\begin{textebox}[Zum Übersetzen]
Ein Schüler aus Yaoundé fragt seinen Lehrer, ob alle Nachrichten im Internet wahr sind. Der Lehrer antwortet, dass man immer auf die Quelle achten muss. Er weiß nicht, woher viele Meldungen kommen, deshalb informiert er sich über mehrere Zeitungen. Die Schüler warten auf seine Erklärung und freuen sich über den interessanten Unterricht.
\end{textebox}
\end{exercice}

\begin{exercice}[Expression écrite : sujet de type Baccalauréat]
\textbf{Sujet :} \all{Wie informierst du dich über die Welt?}

Rédige un texte de 8 à 10 lignes en allemand dans lequel tu expliques comment tu t'informes sur l'actualité (à Yaoundé, au Cameroun, dans le monde). Tu dois obligatoirement employer :
\begin{itemize}
\item au moins \textbf{trois verbes à complément prépositionnel} différents (par exemple : \all{sich informieren über}, \all{warten auf}, \all{sich interessieren für}, \all{denken an}, \all{sich freuen über}) ;
\item au moins \textbf{une interrogative indirecte} avec \all{ob} ou avec un mot interrogatif (\all{Ich weiß nicht, ob \dots} / \all{Ich frage mich, woher \dots}) ;
\item le vocabulaire du thème des médias (\all{die Nachricht}, \all{die Quelle}, \all{die Zeitung}, \all{das Internet}...).
\end{itemize}
Soigne la place du verbe, la majuscule des noms et l'accord des articles.
\end{exercice}

\newpage

\coursec{Corrigés des exercices}

\begin{corrige}[Exercice 1 — QCM : le bon choix]
\begin{enumerate}
\item \textbf{b)} \all{Ich weiß nicht, \textbf{ob} die Nachricht wahr ist.} — La question directe serait \all{Ist die Nachricht wahr?} (question sans mot interrogatif) : on introduit donc l'interrogative indirecte par \all{ob}. \all{dass} introduit un fait, pas une question ; \all{wenn} est temporel ou conditionnel.
\item \textbf{b)} \all{Er fragt, \textbf{wann} der Journalist kommt.} — Question avec mot interrogatif : on reprend le mot interrogatif (\all{wann}) et le verbe conjugué se place \textbf{à la fin} (\all{kommt}). La réponse a) est fausse car le verbe n'est pas en fin de phrase.
\item \textbf{c)} \all{Wir warten \textbf{auf} eine Antwort.} — \all{warten auf + Akkusativ} : préposition figée à apprendre avec le verbe.
\item \textbf{a)} \all{Sie informiert sich \textbf{über} die Ereignisse in Kamerun.} — \all{sich informieren über + Akkusativ} (= se renseigner sur).
\item \textbf{a)} \all{die Zeitung, die Zeitungen} — Les noms féminins en \all{-ung} forment leur pluriel en \all{-en}.
\end{enumerate}
\end{corrige}

\begin{corrige}[Exercice 2 — Vrai ou faux ? Justifie ta réponse]
\begin{enumerate}
\item \textbf{Vrai.} Dans toute subordonnée allemande (donc aussi dans l'interrogative indirecte), le verbe conjugué occupe la \textbf{dernière place} : \all{Ich weiß nicht, ob die Nachricht wahr \emph{ist}.}
\item \textbf{Faux.} C'est le contraire : \all{ob} introduit une question indirecte \textbf{sans} mot interrogatif (réponse oui/non : \all{Kommt er?} $\rightarrow$ \all{Er fragt, ob er kommt.}). Avec un mot interrogatif, on reprend ce mot : \all{Er fragt, \emph{wer} kommt.}
\item \textbf{Faux.} \all{sich freuen} a deux constructions : \all{sich freuen \textbf{auf} + Akkusativ} (avant l'événement, joie anticipée : \all{Ich freue mich auf die Ferien.}) et \all{sich freuen \textbf{über} + Akkusativ} (après l'événement : \all{Ich freue mich über das Geschenk.}).
\item \textbf{Faux.} \all{die Quelle} signifie « la source » (d'une information). « Le journal » se dit \all{die Zeitung}.
\item \textbf{Vrai.} \all{warten auf} est toujours suivi de l'\textbf{accusatif} : \all{Ich warte auf den Bus.} (\all{der Bus} $\rightarrow$ \all{den Bus}).
\end{enumerate}
\end{corrige}

\begin{corrige}[Exercice 3 — Vocabulaire : à chaque mot son article et son pluriel]
\begin{center}
\begin{tabularx}{\linewidth}{|X|l|l|X|}
\hline
\rowcolor{popblueL} Mot allemand & Article & Pluriel & Traduction \\
\hline
Information & die & die Informationen & l'information \\
\hline
Nachricht & die & die Nachrichten & la nouvelle, le message \\
\hline
Quelle & die & die Quellen & la source \\
\hline
Journalist & der & die Journalisten & le journaliste \\
\hline
Zeitung & die & die Zeitungen & le journal \\
\hline
Internet & das & die Internets (rare, souvent sans pluriel) & l'internet \\
\hline
Bericht & der & die Berichte & le reportage, le rapport \\
\hline
Sender & der & die Sender & la chaîne, l'émetteur \\
\hline
Artikel & der & die Artikel & l'article \\
\hline
Leser & der & die Leser & le lecteur \\
\hline
\end{tabularx}
\end{center}
\textbf{Remarques :} les noms en \all{-ung}, \all{-tion}, \all{-e} (féminins) prennent \all{-en} ou \all{-n} au pluriel ; \all{der Journalist} est un nom faible : \all{den/dem Journalisten} à l'accusatif et au datif ; les noms masculins en \all{-er} (\all{Sender}, \all{Leser}) sont invariables au pluriel.
\end{corrige}

\begin{corrige}[Exercice 4 — Conjugaison : complète les verbes]
\begin{enumerate}
\item \all{Ich \textbf{warte auf} den Bus nach Douala.} — \all{warten auf + Akkusativ} (\all{der Bus} $\rightarrow$ \all{den Bus}).
\item \all{Wir \textbf{informieren uns über} die Nachrichten des Tages.} — verbe réfléchi : \all{uns} ; \all{sich informieren über + Akkusativ}.
\item \all{\textbf{Freust du dich auf} dein Geburtstagsgeschenk?} — \all{sich freuen auf + Akkusativ} (joie anticipée) ; attention : \all{du freust} (et non *\all{freuest}).
\item \all{Der Journalist \textbf{fragt nach} der Wahrheit der Information.} — \all{fragen nach + Dativ} (= se renseigner sur) ; \all{die Wahrheit} $\rightarrow$ datif féminin \all{der Wahrheit}.
\item \all{Die Schüler \textbf{sprechen über} die neue Zeitung der Schule.} — \all{sprechen über + Akkusativ}.
\item \all{Sie \textbf{interessiert sich für} Politik und Sport.} — \all{sich interessieren für + Akkusativ}.
\end{enumerate}
\end{corrige}

\begin{corrige}[Exercice 5 — Thème : traduis en allemand]
\begin{enumerate}
\item \all{Der Journalist wartet auf eine Bestätigung.} — \all{warten auf + Akkusativ} ; « confirmation » = \all{die Bestätigung}.
\item \all{Ich weiß nicht, ob die Nachricht wahr ist.} — question indirecte sans mot interrogatif $\rightarrow$ \all{ob}, verbe \all{ist} à la fin.
\item \all{Wir informieren uns im Internet.} — ou \all{über das Internet} si l'on veut dire « au sujet d'Internet » ; « sur Internet » (lieu/moyen) = \all{im Internet} (\all{in dem}).
\item \all{Sie fragt, wo der Journalist wohnt.} — mot interrogatif repris, verbe conjugué \all{wohnt} en dernière position.
\item \all{Die Schüler freuen sich auf die Ferien.} — joie \emph{anticipée} (les vacances ne sont pas encore là) $\rightarrow$ \all{auf + Akkusativ}. Si les vacances sont déjà commencées : \all{freuen sich über die Ferien}.
\end{enumerate}
\end{corrige}

\begin{corrige}[Exercice 6 — Transformation : de la question directe à l'interrogative indirecte]
\begin{enumerate}
\item \all{Er fragt, \textbf{woher die Nachricht kommt}.} — le verbe \all{kommt} passe en fin de subordonnée.
\item \all{Ich weiß nicht, \textbf{ob die Information richtig ist}.} — question sans mot interrogatif $\rightarrow$ \all{ob}.
\item \all{Er fragt, \textbf{wann die Zeitung erscheint}.}
\item \all{Ich weiß nicht, \textbf{ob der Journalist die Quelle genannt hat}.} — au \all{Perfekt}, le verbe conjugué (\all{hat}) reste à la fin, le participe juste avant : \all{\dots genannt hat}.
\item \all{Er fragt, \textbf{wer den Artikel geschrieben hat}.} — même structure au \all{Perfekt}.
\item \all{Ich weiß nicht, \textbf{warum sie keine Zeitung liest}.}
\end{enumerate}
\textbf{Règle rappelée :} mot interrogatif ou \all{ob} + sujet + \dots + verbe conjugué \textbf{en dernière position}.
\end{corrige}

\begin{corrige}[Exercice 7 — Texte à trous : prépositions et cas]
\all{Jeden Morgen informiere ich mich über die Nachrichten. Ich warte auf die Zeitung, die um sieben Uhr kommt. Mein Bruder interessiert sich für den Sport, aber ich denke oft an meine Zukunft. Meine Mutter spricht gern über die Politik, und mein Vater hält nichts von den sozialen Netzwerken.}

\textbf{Justifications :}
\begin{itemize}
\item \all{sich informieren über + Akk.} : \all{die Nachrichten} (pluriel, accusatif = forme identique au nominatif).
\item \all{warten auf + Akk.} : \all{auf die Zeitung}.
\item \all{sich interessieren für + Akk.} : \all{der Sport} $\rightarrow$ \all{für den Sport}.
\item \all{denken an + Akk.} : \all{an meine Zukunft} (féminin accusatif).
\item \all{sprechen über + Akk.} : \all{über die Politik}.
\item \all{nichts halten von + Dat.} : \all{die sozialen Netzwerke} $\rightarrow$ datif pluriel \all{von den sozialen Netzwerken}.
\end{itemize}
\end{corrige}

\begin{corrige}[Exercice 8 — Appariements : le verbe et sa préposition]
\textbf{Appariements corrects :}
\begin{itemize}
\item \all{warten} $\rightarrow$ \all{auf + Akkusativ}
\item \all{sich freuen} (avant l'événement) $\rightarrow$ \all{auf + Akkusativ}
\item \all{sich informieren} $\rightarrow$ \all{über + Akkusativ}
\item \all{denken} $\rightarrow$ \all{an + Akkusativ}
\item \all{sich interessieren} $\rightarrow$ \all{für + Akkusativ}
\item \all{sich freuen} (après l'événement) $\rightarrow$ \all{über + Akkusativ}
\end{itemize}
\textbf{Phrases possibles (exemples) :}
\begin{enumerate}
\item \all{Ich warte auf die Zeitung und informiere mich über die Nachrichten.} « J'attends le journal et je me renseigne sur les nouvelles. »
\item \all{Mein Bruder interessiert sich für Sport und denkt oft an die nächsten Spiele.} « Mon frère s'intéresse au sport et pense souvent aux prochains matchs. »
\end{enumerate}
\end{corrige}

\begin{corrige}[Exercice 9 — Compréhension écrite : texte et questions]
\textbf{Barème indicatif (sur 20) :} questions de compréhension 10 pts (2 pts par question : 1 pt pour le sens, 1 pt pour la langue et la citation exacte) ; questions de langue 6 pts ; qualité générale de l'allemand 4 pts.

\textbf{Questions de compréhension — réponses modèles :}
\begin{enumerate}
\item \all{Der Text handelt von der Mediennutzung junger Menschen in Deutschland: Sie informieren sich über soziale Netzwerke und lernen in der Schule den richtigen Umgang mit Informationen.} « Le texte traite de l'usage des médias par les jeunes en Allemagne : ils s'informent via les réseaux sociaux et apprennent à l'école le bon usage de l'information. »
\item \all{Das Hauptproblem ist, dass nicht jede Nachricht im Internet wahr ist: Es gibt falsche Nachrichten, die sogenannten „Fake News“, und viele Nutzer prüfen die Quelle nicht.} « Le problème principal : toutes les nouvelles sur Internet ne sont pas vraies ; il existe des « fake news » et beaucoup d'utilisateurs ne vérifient pas la source. »
\item \all{Die jungen Deutschen informieren sich heute über soziale Netzwerke und das Handy.} Citation : \all{„In Deutschland informieren sich viele junge Menschen nicht mehr über die klassische Zeitung, sondern über soziale Netzwerke.“}
\item \all{Die Journalisten raten, immer die Quelle zu prüfen.} Citation : \all{„Sie erklären, dass man immer die Quelle prüfen muss.“}
\item \all{Die Lehrerin spricht über die Medien, weil ihre Schüler die Nachrichten oft nicht wirklich verstehen.} Citation : \all{„Ich weiß oft nicht, ob meine Schüler eine Nachricht wirklich verstehen. Deshalb sprechen wir im Unterricht über die Medien.“}
\end{enumerate}

\textbf{Questions de langue :}
\begin{enumerate}
\item Deux interrogatives indirectes (verbe conjugué souligné) :
\begin{itemize}
\item \all{\dots und viele Nutzer fragen nicht, \textbf{woher} die Meldung \emph{kommt}.} (mot interrogatif + verbe final)
\item \all{Ich weiß oft nicht, \textbf{ob} meine Schüler eine Nachricht wirklich \emph{verstehen}.} (\all{ob} + verbe final)
\end{itemize}
\item Trois verbes à complément prépositionnel :
\begin{itemize}
\item \all{sich informieren über + Akkusativ} : \all{informieren sich \dots über soziale Netzwerke}
\item \all{warten auf + Akkusativ} : \all{Sie warten nicht auf die Abendnachrichten}
\item \all{warnen vor + Dativ} : \all{Journalisten warnen vor falschen Nachrichten}
\end{itemize}
(On accepte aussi : \all{sprechen über + Akk.}, \all{sich interessieren für + Akk.}, \all{denken über \dots nach + Akk.})
\end{enumerate}

\textbf{Points de vigilance :} recopier les citations \emph{mot à mot} avec les guillemets allemands „\dots“ ; répondre par des phrases complètes avec le verbe en deuxième position ; ne pas confondre \all{über} (au sujet de) et \all{vor} (contre, devant).
\end{corrige}

\begin{corrige}[Exercice 10 — Thème et version]
\textbf{Barème indicatif (sur 20) :} thème 10 pts (2 pts par phrase : préposition et cas 1 pt, place du verbe et lexique 1 pt) ; version 10 pts (fidélité 6 pts, qualité du français 4 pts).

\textbf{A. Thème — traductions modèles :}
\begin{enumerate}
\item \all{Der Journalist wartet auf eine Bestätigung seiner Quelle.} — \all{warten auf + Akk.} ; génitif \all{seiner Quelle}.
\item \all{Ich weiß nicht, ob die Nachricht wahr ist.} — \all{ob} + verbe final.
\item \all{Wir informieren uns im Internet und in den Zeitungen.} — \all{in + Dativ} (lieu) : \all{im} = \all{in dem} ; \all{in den Zeitungen} (datif pluriel).
\item \all{Sie fragt, wann der Artikel erscheint.} — mot interrogatif \all{wann} repris, verbe \all{erscheint} à la fin.
\item \all{Viele junge Leute interessieren sich für Politik, aber sie denken nicht immer an die Gefahren der sozialen Netzwerke.} — \all{sich interessieren für + Akk.} ; \all{denken an + Akk.} : \all{an die Gefahren}.
\end{enumerate}

\textbf{B. Version — traduction modèle :}

« Un élève de Yaoundé demande à son professeur si toutes les nouvelles sur Internet sont vraies. Le professeur répond qu'il faut toujours faire attention à la source. Il ne sait pas d'où viennent beaucoup d'informations ; c'est pourquoi il se renseigne dans plusieurs journaux. Les élèves attendent son explication et se réjouissent de ce cours intéressant. »

\textbf{Points de vigilance :}
\begin{itemize}
\item \all{ob alle Nachrichten \dots wahr sind} : interrogative indirecte $\rightarrow$ « si » en français (et non « quand »).
\item \all{auf die Quelle achten} = « faire attention à la source » (faux ami : \all{achten} n'est pas « acheter »).
\item \all{sich freuen über + Akk.} = « se réjouir de » (événement présent ou passé).
\item \all{woher viele Meldungen kommen} = « d'où viennent beaucoup d'informations » : respecter l'ordre des mots français (sujet avant verbe).
\end{itemize}
\end{corrige}

\begin{corrige}[Exercice 11 — Expression écrite : sujet de type Baccalauréat]
\textbf{Barème indicatif (sur 20) :} respect du sujet et de la longueur (8 à 10 lignes) 4 pts ; emploi d'au moins trois verbes à complément prépositionnel correctement construits 6 pts ; au moins une interrogative indirecte correcte 3 pts ; vocabulaire des médias 3 pts ; correction grammaticale générale (place du verbe, majuscules, cas) 4 pts.

\textbf{Proposition de rédaction modèle :}

\all{Ich informiere mich jeden Tag über die Nachrichten in Kamerun und in der Welt. Am Morgen lese ich die Zeitung oder höre ich das Radio, und am Abend warte ich auf die Nachrichten im Fernsehen. Oft benutze ich auch mein Handy, denn das Internet bringt die Information sehr schnell. Aber ich weiß nicht immer, ob alle Meldungen im Internet wahr sind. Deshalb frage ich mich, woher eine Nachricht kommt, und ich achte auf die Quelle. Ich interessiere mich besonders für Sport und Politik, und ich denke oft an die Zukunft meines Landes. Wenn ein Artikel gut ist, freue ich mich über die gute Arbeit der Journalisten. Meine Lehrer sagen, dass man die Medien klug benutzen muss, und ich finde das richtig.}

(Traduction : « Je me renseigne chaque jour sur les nouvelles du Cameroun et du monde. Le matin, je lis le journal ou j'écoute la radio, et le soir j'attends les informations à la télévision. Souvent, j'utilise aussi mon portable, car Internet apporte l'information très vite. Mais je ne sais pas toujours si toutes les informations sur Internet sont vraies. C'est pourquoi je me demande d'où vient une nouvelle et je fais attention à la source. Je m'intéresse surtout au sport et à la politique, et je pense souvent à l'avenir de mon pays. Quand un article est bon, je me réjouis du bon travail des journalistes. Mes professeurs disent qu'il faut utiliser les médias intelligemment, et je trouve cela juste. »)

\textbf{Vérification des contraintes :}
\begin{itemize}
\item Verbes à complément prépositionnel : \all{sich informieren über}, \all{warten auf}, \all{achten auf}, \all{sich interessieren für}, \all{denken an}, \all{sich freuen über} (six, soit plus que le minimum exigé).
\item Interrogatives indirectes : \all{ich weiß nicht immer, ob alle Meldungen im Internet wahr sind} ; \all{ich frage mich, woher eine Nachricht kommt}.
\item Vocabulaire du thème : \all{die Nachrichten}, \all{die Zeitung}, \all{das Internet}, \all{die Quelle}, \all{der Artikel}, \all{die Journalisten}.
\end{itemize}

\textbf{Points de vigilance :}
\begin{itemize}
\item Ne jamais écrire *\all{ich informiere mich auf} : c'est \all{über + Akkusativ}.
\item Dans l'interrogative indirecte, verbe conjugué \textbf{à la fin} : *\all{ob sind alle Meldungen wahr} est une faute éliminatoire.
\item Majuscule à tous les noms : \all{das Handy}, \all{die Quelle}, \all{das Radio}.
\item Éviter la répétition de \all{und} : varier avec \all{denn}, \all{deshalb}, \all{aber}.
\end{itemize}
\end{corrige}

\newpage

\coursec{Fiche bilan}

\begin{popfigure}
\begin{tikzpicture}[mindmap, grow cyclic, every node/.style={concept}, concept color=popblue,
root concept/.append style={concept color=popdark, font=\bfseries\large, text=white},
level 1/.append style={level distance=3.6cm, sibling angle=51.4, font=\small},
level 2/.append style={level distance=2.2cm, font=\footnotesize}]
\node[root concept, align=center]{AL16\\La gestion\\de l'information}
  child[concept color=popblue!70]{ node[align=center]{Texte support\\Informationen im digitalen Zeitalter}
    child { node {global puis détaillé} }
    child { node {questions + résumé} }
  }
  child[concept color=poporange!80]{ node[align=center]{Lexique\\Medien und Information}
    child { node {die Quelle, -n} }
    child { node {die Nachricht, -en} }
  }
  child[concept color=popgreen!80]{ node[align=center]{Grammaire 1\\Interrogation indirecte}
    child { node {ob = si} }
    child { node {w-Fragen : verbe en fin} }
  }
  child[concept color=poppurple!80]{ node[align=center]{Grammaire 2\\Verbes à préposition}
    child { node {warten auf + Akk.} }
    child { node {sich informieren über + Akk.} }
  }
  child[concept color=poppink!80]{ node[align=center]{Conjugaison\\verbes + prépositions}
    child { node {denken an, glauben an} }
    child { node {sich freuen auf/über} }
  }
  child[concept color=popgold!90]{ node[align=center]{Méthode\\commentaire}
    child { node {relever -- expliquer -- juger} }
  }
  child[concept color=popgreen!60]{ node[align=center]{Production\\guidée}
    child { node {court paragraphe} }
    child { node {avis argumenté} }
  };
\end{tikzpicture}
\legende{Carte mentale de la leçon AL16 : la gestion de l'information}
\end{popfigure}

\begin{bilan}
\textbf{1. Lexique clé (à connaître avec article et pluriel)}
\begin{itemize}
\item \all{die Information, -en} : l'information ; \all{die Nachricht, -en} : la nouvelle, l'information (médias)
\item \all{die Quelle, -n} : la source ; \all{der Journalist, -en / die Journalistin, -nen} : le/la journaliste
\item \all{die Zeitung, -en} : le journal ; \all{das Internet, -s} : l'internet ; \all{die Medien} (pl.) : les médias
\item \all{die Nachrichten} (pl.) : les informations (télévisées) ; \all{die Fake News} (pl.) : les fausses informations
\item \all{veröffentlichen} : publier ; \all{verbreiten} : diffuser ; \all{überprüfen} : vérifier ; \all{berichten} : rapporter
\item \all{die Meinung, -en} : l'opinion ; \all{glaubwürdig} : crédible ; \all{zuverlässig} : fiable
\end{itemize}

\textbf{2. Grammaire à connaître par c\oe ur}
\begin{center}
\begin{tabularx}{\linewidth}{|X|X|}
\hline
\rowcolor{popblueL} \textbf{Règle} & \textbf{Exemple} \\
\hline
Interrogation indirecte avec \all{ob} (= si) : le verbe conjugué va en \textbf{fin de phrase}. & \all{Ich weiß nicht, ob die Quelle zuverlässig ist.} \\
\hline
Interrogation indirecte avec mot interrogatif (\all{wer, was, wann, wo, warum, wie}) : mot interrogatif + sujet + verbe en fin. & \all{Er fragt, woher die Information kommt.} \\
\hline
Verbe + préposition figée : la préposition impose son cas. & \all{warten auf + Akk., sich informieren über + Akk., denken an + Akk., glauben an + Akk.} \\
\hline
Renvoi à une chose : \all{worauf, worüber, daran, darüber} (jamais \all{auf es}). & \all{Worauf wartest du? -- Daran glaube ich nicht.} \\
\hline
\all{sich freuen auf + Akk.} (futur) / \all{sich freuen über + Akk.} (présent ou passé). & \all{Wir freuen uns auf die Sendung. / über die gute Nachricht.} \\
\hline
\end{tabularx}
\end{center}

\textbf{3. Méthodes express}
\begin{itemize}
\item \textbf{Compréhension} : lire le titre et le premier paragraphe (sens global), puis repérer les mots-clés du lexique des médias avant de répondre aux questions détaillées.
\item \textbf{Résumé} : une phrase d'introduction (thème + source), puis les idées principales avec ses propres mots, sans citer de phrases entières.
\item \textbf{Traduction} : repérer d'abord la préposition du verbe allemand, puis vérifier le cas exigé ; placer le verbe en fin dans toute subordonnée avec \all{ob} ou mot interrogatif.
\item \textbf{Production} : donner son avis avec \all{Meiner Meinung nach... / Ich bin der Ansicht, dass...} et justifier par un exemple concret (Cameroun, réseaux sociaux, presse).
\end{itemize}
\end{bilan}

\begin{aretenir}[Checklist avant le Bac]
\begin{itemize}
\item[$\square$] Je connais le lexique des médias avec article et pluriel : \all{die Information, die Nachricht, die Quelle, der Journalist, die Zeitung, das Internet}.
\item[$\square$] Je distingue \all{die Nachricht} (une nouvelle) et \all{die Nachrichten} (le journal télévisé).
\item[$\square$] Je place le verbe conjugué en fin de phrase dans une interrogative indirecte avec \all{ob}.
\item[$\square$] Je place le verbe en fin après un mot interrogatif : \all{..., warum er das sagt.}
\item[$\square$] Je sais que \all{ob} traduit « si » (condition indirecte), jamais « si » de condition (\all{wenn}).
\item[$\square$] Je connais par c\oe ur : \all{warten auf, sich informieren über, denken an, glauben an, sich freuen auf/über, berichten über}.
\item[$\square$] J'emploie \all{worauf / worüber} pour interroger sur une chose et \all{daran / darüber} pour y répondre.
\item[$\square$] Je distingue \all{sich freuen auf} (avant) et \all{sich freuen über} (après).
\item[$\square$] Je mets une majuscule à tous les noms allemands et les guillemets allemands „\ldots“ dans mes citations.
\item[$\square$] Je sais résumer un texte en trois ou quatre phrases avec mes propres mots.
\item[$\square$] Je sais exprimer mon opinion sur les médias et les fausses informations avec \all{Meiner Meinung nach} et un argument.
\end{itemize}
\end{aretenir}

\findecours

\end{document}
